\documentclass[11pt]{article}
\usepackage[T1]{fontenc}
\usepackage[utf8]{inputenc}
\usepackage{lmodern}
\usepackage[margin=1in]{geometry}
\usepackage{amsmath,amssymb,amsthm}
\usepackage{enumitem}
\usepackage{booktabs,longtable,array}
\usepackage[hidelinks]{hyperref}
\allowdisplaybreaks
\emergencystretch=3em
\tolerance=2000

\newtheorem{theorem}{Theorem}
\newtheorem*{thmA}{Theorem A}
\newtheorem*{thmB}{Theorem B}
\newtheorem{lemma}{Lemma}[section]
\newtheorem{proposition}[lemma]{Proposition}
\newtheorem{corollary}[lemma]{Corollary}
\theoremstyle{definition}
\newtheorem{definition}[lemma]{Definition}
\newtheorem{remark}[lemma]{Remark}
\newtheorem{procedure}[lemma]{Procedure}

\newcommand{\E}{\mathbb E}
\newcommand{\Z}{\mathbb Z}
\newcommand{\R}{\mathbb R}
\newcommand{\Tc}{\mathbb T}
\newcommand{\eps}{\varepsilon}
\newcommand{\bc}{\bar c}
\newcommand{\abs}[1]{\lvert #1\rvert}
\newcommand{\floor}[1]{\lfloor #1\rfloor}
\newcommand{\ceil}[1]{\lceil #1\rceil}
\newcommand{\poly}{\operatorname{poly}}
\newcommand{\Var}{\operatorname{Var}}
\newcommand{\av}[1]{\langle #1\rangle}
\newcommand{\cG}{\mathcal G}
\newcommand{\cP}{\mathcal P}
\newcommand{\cE}{\mathcal E}
\newcommand{\cA}{\mathcal A}
\newcommand{\cD}{\mathcal D}
\newcommand{\cC}{\mathcal C}
\newcommand{\cT}{\mathcal T}
\newcommand{\Ntr}{N_{\mathrm{tr}}}
\newcommand{\key}{\mathsf k}
\newcommand{\pat}{\operatorname{pat}}
\DeclareMathOperator{\KL}{KL}

\title{A self-contained proof of the exponent threshold $0.4875754036$ for worst-case Subset Sum\\[2pt]
\large Preprint --- not refereed}
\author{Jennifer Voitle\\ Independent researcher}
\date{November 2026}
\hypersetup{pdftitle={A self-contained proof of the exponent threshold 0.4875754036 for worst-case Subset Sum},
 pdfauthor={Jennifer Voitle}}

\begin{document}
\maketitle
\begin{abstract}
We give a complete written proof that Subset Sum on $n$ positive integers of polynomial bit length is decided by a
randomized word-RAM algorithm in time $O(2^{T'n})$ for every fixed $T'>T=0.4875754036$, in particular in time
$O(2^{0.4875755n})$. The threshold $T$ was reported earlier by the author, conditionally on the non-numerical
lemmas of two preprints; the contribution of this paper is the derivation that removes that condition.
The algorithm is the one introduced by OpenAI (isolation, a capped preliminary search, a
filtered modular checksum, subtraction of aliases) with the alias step of A.~Kapila (one large shared block, sparse
masks, alias phase sums) and a single asymmetric checksum filter. Those ideas are theirs. What is done here is
different in kind: every non-numerical lemma that the result uses is stated with its hypotheses and proved in this
text, so that the argument can be checked without accepting unproved statements from other preprints. The numerical
inequalities that the proof uses are listed explicitly and matched one by one to a ball-arithmetic certificate,
whose analytic error bounds are also proved here. The theorem has no mathematical hypothesis beyond two cited
classical theorems and those certified inequalities.
The manuscript also records, statement by statement, how each proof was obtained and which independent checks are
still open. It has been reviewed and audited, with the sources open, by a second
AI system; it has not been refereed.
\end{abstract}

\tableofcontents

\section{Main result, provenance and credits}\label{sec:result}

\paragraph{Related results.} From 1974 until October 2026 the best known exponent for worst-case Subset Sum was
$1/2$, the exponent of the meet-in-the-middle algorithm of Horowitz and Sahni~\cite{HS}; the first worst-case
improvement on that algorithm, by Chen, Jin, Randolph and Servedio~\cite{CJRS}, was by a polynomial factor. In
October 2026 OpenAI~\cite{OAI} and, independently and by a different method, Li and Ye~\cite{LY} gave algorithms
with exponent below $1/2$. Within a few weeks a sequence of manuscripts lowered the exponent of the OpenAI
construction to about $0.4876$. The table lists these results. An algorithm with exponent $c$ runs in time about
$2^{cn}$, so a smaller exponent is a stronger result.
\begin{center}
\small
\begin{tabular}{@{}>{\raggedright\arraybackslash}p{0.15\linewidth}>{\raggedright\arraybackslash}p{0.21\linewidth}%
>{\raggedright\arraybackslash}p{0.27\linewidth}>{\raggedright\arraybackslash}p{0.27\linewidth}@{}}
\toprule
Exponent & Work & Method & Status of the proof\\
\midrule
$1/2$ & Horowitz and Sahni~\cite{HS}, 1974 & meet in the middle & published\\
$1/2$, with time $2^{n/2}/\poly(n)$ & Chen, Jin, Randolph and Servedio~\cite{CJRS}, 2023 & meet in the middle,
improved by a polynomial factor & published\\
$1/2-1/450000$ $\approx0.4999978$ & Li and Ye~\cite{LY}, 7~October 2026 & representation technique; partial match
by depth-2 linear circuits & preprint\\
$0.49$ & OpenAI~\cite{OAI}, 4~October 2026 & isolation, preliminary search, filtered modular checksum, subtraction
of aliases & manuscript\\
$0.4877$ & Kapila~\cite{Base}, 14~October 2026 & OpenAI's construction with a new alias step and 43 filters & manuscript; uses
lemmas of~\cite{OAI}\\
$0.487676$ & Wharton~\cite{W}, 20~October 2026 & Kapila's construction with retuned parameters & manuscript; conditional
on~\cite{Base}\\
$0.4876723$ & Bills~\cite{B}, 21~October 2026 & Kapila's construction with the certificate boundary of one filter moved &
manuscript; conditional on~\cite{Base}\\
$0.4875754036$ & the author~\cite{V}, November 2026 & Kapila's construction with one asymmetric filter in place of 43 & preprint;
numerical premises certified; conditional on the non-numerical lemmas of~\cite{OAI} and~\cite{Base}\\
$0.4875754036$ & this paper & the algorithm and parameter point of~\cite{V} & every lemma proved in the text;
numerical premises certified; not refereed\\
\bottomrule
\end{tabular}
\end{center}
In the last two rows the number is a threshold: the theorem gives every exponent above it.

\paragraph{Why this paper was written.} The exponents below $0.49$ were obtained within a few weeks, and each rests
on the ones before it. The proof in~\cite{Base} uses lemmas of~\cite{OAI}. The results of~\cite{W} and~\cite{B}
retune the construction of~\cite{Base} and are conditional on its proof. The author's earlier preprint~\cite{V}
certifies the numerical conditions at a new parameter point with a single filter, and states its theorem under the
assumption that the non-numerical lemmas of~\cite{OAI} and~\cite{Base} hold. To the author's knowledge none of
these manuscripts has been refereed, and they differ in notation and conventions. A reader who wants to check the
exponent $0.4875754036$ therefore has to work through three manuscripts and to verify, wherever one relies on
another, that the hypotheses of the cited statement are met at the new parameters.

This paper removes that dependence. It does not lower the exponent and it introduces no new algorithmic idea: the
algorithm and the parameter point are those of~\cite{V}, with the small changes recorded in
Section~\ref{sec:sources}. What it adds is a single document in which the algorithm is specified in full, every
non-numerical lemma that the result uses is stated with its hypotheses and proved in one notation, and the numerical
premises are listed explicitly and certified. The threshold of~\cite{V} thereby becomes a theorem that assumes no
statement of another preprint (Theorem~\ref{thm:main}). To the author's knowledge it is the smallest exponent
claimed for worst-case Subset Sum at the date of this manuscript. A complete written proof is not the same as a
verified one: Sections~\ref{sec:status} and~\ref{sec:verification} state what has been checked, by whom and how, and what
has not.

\subsection{The theorem}
An instance of Subset Sum consists of $n\ge2$ positive integers $a_1,\dots,a_n$ and a nonnegative integer $t$; the
question is whether some $I\subseteq[n]=\{1,\dots,n\}$ has $a(I):=\sum_{i\in I}a_i=t$. Equal integers and the empty
set are allowed. Let $b\ge1$ be the largest bit length among $t,a_1,\dots,a_n$. The machine model is fixed in
Section~\ref{sec:model}. Throughout, $\log$ is the logarithm to base $2$ and
$h(x)=-x\log x-(1-x)\log(1-x)$ is the binary entropy.

\begin{theorem}[main result]\label{thm:main}
Let $T=0.4875754036$. There is a single randomized program for Subset Sum on the word RAM of
Section~\ref{sec:model} with the following properties.
\begin{enumerate}[label=\textup{(\alph*)}]
\item On every input with $n\ge2$ it halts on every execution, and it returns the correct answer with probability at
least $2/3$ over its own random words. This holds for \textsc{yes} instances and for \textsc{no} instances.
\item For every fixed positive integer $c$ and every fixed real $T'>T$ there are constants $C$ and $N$ such that
every input with $n\ge N$ and $b\le n^c$ is processed in at most $C\,2^{T'n}$ word operations, on every outcome of
the random choices. In particular the running time is $O(2^{0.4875755\,n})$. The program depends on neither $c$ nor
$T'$.
\end{enumerate}
\end{theorem}

\paragraph{Status of the proof.} A complete written proof of Theorem~\ref{thm:main} is given in this paper and its
numerical premises are certified. The written proof has had a cross-model review and audit, and it has not been
refereed. Section~\ref{sec:status} says what was checked, by which system and how.

The intermediate bound, valid on every execution on an input that passes the guard of Section~\ref{sec:params}, is
$K_0\,2^{(\tau_0+20\eps)n+2n^{0.7}}(n+b+2)^{K}$ with $\tau_0+20\eps=T$ exactly (Proposition~\ref{prop:time}).
The exponent is never rounded downward: $0.4875754$ would \emph{not} be a valid statement of the result. The
number $T$ is a threshold: the theorem gives every exponent $T'>T$, not $T$ itself.

\paragraph{Integer inputs of either sign.} The theorem is stated for positive $a_i$ and $t\ge0$. Arbitrary integers
reduce to this case by a standard preprocessing step, which is not used elsewhere in the paper. Let
$\mathcal N=\{i:a_i<0\}$. For $I\subseteq[n]$ one has $\sum_{i\in I\triangle\mathcal N}\abs{a_i}=a(I)-a(\mathcal N)$,
and $I\mapsto I\triangle\mathcal N$ is a bijection of the subsets of $[n]$; so the instance $(a,t)$ is equivalent
to the instance with weights $\abs{a_i}$ and target $t-a(\mathcal N)$, which is answered \textsc{no} at once if
that target is negative. Indices with $a_i=0$ are deleted (if fewer than two weights remain the instance is decided
by inspection). The bit length grows by at most $\ceil{\log(n+1)}$, so polynomial bit length is preserved.

\paragraph{Structure of the proof.} The argument has four parts, and each depends only on the ones before it.
\begin{enumerate}[label=\textup{(\roman*)}]
\item \emph{A dichotomy.} Isolation (Lemma~\ref{lem:iso}) replaces the instance by polynomially many targets $t_s$,
one of which, on a \textsc{yes} instance, has a unique solution with probability at least $11/12$. The calculus of
deficiencies (Lemmas~\ref{lem:def}--\ref{lem:fmono}) shows that for each fixed choice of tie weights either the
low-deficiency condition \textup{(LD)} holds, or the capped preliminary search returns a correct answer except with
probability at most $8/(\zeta^2n)+n2^{-n^3}$, for $n>1/\zeta$ (Proposition~\ref{prop:P}).
\item \emph{Theorem~A} (Section~\ref{sec:alias}). Under \textup{(LD)}, for a random prime $P$ the phase sums over all
\emph{aliases} of every target (subsets whose sum agrees with the target modulo $P$ only) are computed within the
work budget. The inputs are the averaged small-layer bound (Corollary~\ref{cor:D}), collisions modulo random primes
(Lemma~\ref{lem:primes}) and the trial count (Lemma~\ref{lem:A3}).
\item \emph{Theorem~B} (Sections~\ref{sec:filter}--\ref{sec:checksum}). For random phases, the phase sum over
\emph{all} subsets in a residue class modulo $P$ is computed within the budget. The inputs are the bounds for the
filter functions, the two certified moments \textup{(M1)}, \textup{(M2)} transferred to the finite circle
(Lemma~\ref{lem:B1}), and the tables and inverse-multiplicity sampling of Section~\ref{sec:checksum}.
\item \emph{Assembly} (Sections~\ref{sec:final}--\ref{sec:time}). The difference of the two sums is the phase sum
over exact solutions (Lemma~\ref{lem:cancel}), which is $0$ on a \textsc{no} instance and has modulus $1$ for an
isolated target. Caps on every loop turn the expected-work bounds into bounds valid on every execution.
\end{enumerate}
Section~\ref{sec:cert} certifies the numerical premises and Section~\ref{sec:limits} states the limits of the
analysis.

\subsection{What is proved here, what is certified, what is cited}\label{sec:status}
The proof of Theorem~\ref{thm:main} has three kinds of ingredients.
\begin{description}[leftmargin=1.5em]
\item[Written proofs.] Every lemma, proposition and theorem in Sections~\ref{sec:model}--\ref{sec:time} is proved
in the text. No step is of the form ``as shown in'' another preprint.
\item[Certified numerics.] The proof uses the scalar inequalities \textup{(E0)}, \textup{(C1)--(C8)},
\textup{(N1)--(N14)} and the
two moment inequalities \textup{(M1)}, \textup{(M2)} of Section~\ref{sec:params}, at the frozen parameter point
given there. They are verified in ball arithmetic by the certificate script supplied with this paper
(Section~\ref{sec:cert}). The independently written implementations of the moment computation are those reported
in~\cite{V}; they are historical checks and are not supplied here. The analytic bounds on which the moment certificate relies are proved in
Lemma~\ref{lem:strip}. A certificate is a proof of a numerical inequality; it is not a proof of the lemmas.
\item[Cited classical results.] Two published theorems are used and not re-proved: an explicit lower bound for the
number of primes in $(x,2x]$ (Rosser and Schoenfeld~\cite{RS}; used in Lemma~\ref{lem:primes}) and, inside the
certificate only, the trapezoid-rule error bound for analytic periodic functions (Trefethen and
Weideman~\cite{TW}). Elementary inequalities (Markov, Chebyshev, Cauchy--Schwarz, Jensen) are used freely; the
Azuma--Hoeffding inequality is proved in the form we need (Lemma~\ref{lem:azuma}). The proof of
Lemma~\ref{lem:strip} also uses textbook analysis without citation: the Fourier series of an entire periodic
function, Parseval's identity and Cauchy's theorem.
\end{description}
The open combinatorial statement called \textup{[S]} in~\cite{V} (``small slices are robust'') is \emph{not} used
anywhere in this paper; the conditional exponent $0.4873407528$ of~\cite{V} is not claimed here.

\paragraph{Logical status and verification status.} These are different things and we keep them apart.
\emph{Logically}, Theorem~\ref{thm:main} has no hypothesis: its proof rests on the three ingredients above and on
nothing else, and in particular it assumes no statement of~\cite{OAI} or~\cite{Base}. Whether the written proofs
are correct is a question of \emph{verification}, and a review is evidence about that question, not a premise of
the theorem. Section~\ref{sec:verification} gives the verification status of every statement group. In brief: all
proofs in this text were reconstructed by an AI system, Claude Fable~5.2 (Anthropic; below ``Claude''), working from
the sources~\cite{OAI,Base}, and then
rewritten; three groups (the small-layer lemmas, the one-trial lemma, the domination step) had earlier been
re-derived by two different AI systems with partial blindness. A source-aware reconstruction is not an independent
proof. Blind re-derivations from statement-only
packets were made for most statement groups (Section~\ref{sec:verification}). The manuscript was then reviewed by a different AI system, GPT-6
Astra (OpenAI, at its high reasoning setting; below ``Astra''), with the sources open: it traced the main proof chain, checked the analytic basis of the certificate
against the code and re-ran the certificate. It then re-derived the small-layer lemmas and the one-trial lemma
from their statements and audited the rest of the proof chain with the manuscript open
(Section~\ref{sec:verification}). None of this is a blind re-derivation or a formal verification, and no human referee
has read the manuscript. That is the basis of the status stated after Theorem~\ref{thm:main}.

\subsection{Credits}
The architecture (tie weights and isolation, the capped preliminary search, the deficiency dichotomy, the
random-phase modular checksum with block tables, patterns, inverse-multiplicity sampling, and alias subtraction) is
due to OpenAI~\cite{OAI}; it builds on the isolation lemma of Mulmuley, Vazirani and Vazirani~\cite{MVV}, on meet in
the middle~\cite{HS}, and on inverse-multiplicity sampling in the sense of Karp, Luby and Madras~\cite{KLM}. The alias step used here (one shared block with sparse masks, phase sums over alias
classes, the sampled trial count, right classes, and the averaged small-layer bound, Corollary~\ref{cor:D}) and the
treatment of the checksum error through nonnegative multipliers are due to A.~Kapila~\cite{Base}, whose manuscript
proves the exponent $0.4877$. The parameter point and the single asymmetric filter are from~\cite{V}. The present
text contributes the complete self-contained proof at that point, the explicit list of numerical premises, a small
number of changes relative to the sources and explicit constants (Section~\ref{sec:sources}), and the disclosure of what has and has not
been independently checked. No sentence of~\cite{OAI} or~\cite{Base} is reproduced; statements are re-expressed in
the notation of this paper, which differs from both.

\subsection{Parameters and numerical premises}\label{sec:params}
The following decimals are exact rational numbers (the file \texttt{candidate\_base.json} of~\cite{V}, supplied
again with this manuscript; its decimal strings are the authoritative values).
\begin{center}
\begin{tabular}{llp{0.56\linewidth}}
\toprule
Symbol & Value & Role\\
\midrule
$\tau_0$ & $0.4875752036$ & work exponent\\
$\eta$ & $0.5162721413$ & the large prime $P$ has $\floor{\eta n}+1$ bits\\
$\mu$ & $0.4752421539$ & the shared block $M$ has $m\approx\mu n$ elements\\
$a$ & $0.0082379228$ & left masks have $k\approx am$ elements\\
$\gamma$ & $0.0605824120$ & the small prime $P'$ has $\floor{\gamma m}+1$ bits\\
$l$ & $0.4836347962$ & the left part $L$ has $\floor{ln}$ elements\\
$r$ & $1-\mu-l=0.0411230499$ & the right part $R$ has about $rn$ elements\\
$\kappa$ & $1.03465831$ & power applied to the filter densities\\
$p_1,\ p_2$ & $8.45902160,\ 8.27015425$ & multipliers of the two moments ($p$ and $q_m$ in~\cite{V})\\
$\eps,\ \zeta$ & $10^{-8},\ 10^{-8}$ & pattern accuracy; slack in the preliminary threshold\\
$G_0$ & $10^{8}$ & guard coefficient\\
$\rho_f,\ \rho_c$ & $\eta-1.5\cdot10^{-6},\ \rho_f-10^{-6}$ & finite-circle and continuous moment thresholds\\
\bottomrule
\end{tabular}
\end{center}
Derived constants: $\Delta=\eta-\tau_0=0.0286969377$, $\Gamma=2\tau_0-\eta=0.4588782659$, $a_\Gamma=\Delta/\Gamma$,
$\delta=\tfrac12\bigl(\eta-\tfrac\mu2\log3\bigr)$, and $T=\tau_0+20\eps=0.4875754036$. The filter is the pair of real
trigonometric polynomials $r,s$ of degree $16$ (sixteen cosine and sixteen sine coefficients each, eight decimals)
listed in Appendix~\ref{app:filter}; Section~\ref{sec:filter} uses only the norms
$\|r\|_{1}\le3.64$, $\|s\|_1\le3.60$ of their coefficient vectors, and the moment premises below depend on the
coefficients themselves.

\paragraph{Sizes.} For each $n$,
\[
 m=2\floor{\mu n/2},\quad \sigma=m/2,\quad k=\floor{am},\quad v=m-2k,\quad \abs L=\floor{ln},\quad \abs R=n-m-\abs L .
\]

\paragraph{Guard.} An input passes the guard if $G_0\ceil{\log(b+1)}\le n$. Then $b+2\le 2^{n/G_0+1}$. For each fixed
$c$, every input with $b\le n^c$ passes the guard once $n\ge G_0\ceil{\log(n^c+1)}$.

\paragraph{Scalar premises.} The proof uses exactly the following strict inequalities.
\begin{align*}
&\textup{(C1)}\ \gamma<\tfrac12h(2a) &&\textup{(C2)}\ l+(h(a)-\gamma)\mu<\tau_0\\
&\textup{(C3)}\ r+(1-\gamma)\mu<\tau_0 &&\textup{(C4)}\ 1-\mu+\bigl(h(a)+h(\tfrac12-a)-\gamma\bigr)\mu-\eta<\tau_0\\
&\textup{(C5)}\ h(\tfrac12-a)\mu+\tfrac12-\tau_0<\tau_0 &&\textup{(C6)}\ 0<r,\quad r<\tau_0,\quad l<\tau_0\\
&\textup{(C7)}\ \mu(1-2a)<l &&\textup{(C8)}\ \tfrac\mu2\log3<\eta
\end{align*}
The \emph{slack} of an inequality is its right side minus its left side. Put
\[
 \eps_0=0.0000013156 ,
\]
an exact decimal. The proof uses \textup{(C1)--(C8)} in the following quantitative form.
\begin{itemize}[leftmargin=3.2em]
\item[\textup{(E0)}] Each of the inequalities \textup{(C1)--(C8)} has slack greater than $\eps_0$.
\end{itemize}
(Only \textup{(C2)--(C5)} and the last two parts of \textup{(C6)} are needed with the slack $\eps_0$; for the
others positivity suffices.) The number $\eps_0$ is the least of these slacks, that of \textup{(C4)}, rounded down
to five significant digits. It is fixed as a constant, rather than defined as the least slack, so that every
inequality containing $\eps_0$ is an inequality between explicitly given numbers. Further:
\begin{align*}
&\textup{(N1)}\ \zeta+1/G_0<\eps_0 &&\textup{(N2)}\ \mu<\tau_0-\eps_0\\
&\textup{(N3)}\ \tfrac12-\tau_0+\zeta<\delta &&\textup{(N4)}\ 4a<1\\
&\textup{(N5)}\ \tfrac12h(2a)+h(a)<1-2a &&\textup{(N6)}\ \mu(1-2a)<\tfrac12\\
&\textup{(N7)}\ h(a)<h(\tfrac12-a) &&\textup{(N8)}\ \tfrac12-\tau_0+\tfrac{\gamma\mu}2<\tau_0-\eps_0,\quad r+\tfrac12-\tau_0<\tau_0-\eps_0\\
&\textup{(N9)}\ \eta/2<\tau_0 &&\textup{(N10)}\ \Delta>0,\quad\Gamma>0\\
&\textup{(N11)}\ a_\Gamma\le\tfrac34 &&\textup{(N12)}\ p_1\ge1,\quad p_2\ge1\\
&\textup{(N13)}\ \bigl(18+6\max(p_1,p_2)\bigr)\eps<\eta-\rho_f &&\textup{(N14)}\ \rho_c<\rho_f
\end{align*}
\paragraph{Moment premises.} With the functions $u,\omega$ of Section~\ref{sec:filter} on the continuous circle and
$\Lambda(x)=2+2\cos x$,
\[
\textup{(M1)}\ \int_{\Tc}\Lambda(x)\,2^{p_1(\Delta-u(x))}\,dx<2^{\rho_c},\qquad
\textup{(M2)}\ \int_{\Tc}\Lambda(x)\,2^{p_2(a_\Gamma\omega(x)-u(x))}\,dx<2^{\rho_c}.
\]
Section~\ref{sec:cert} reports the certified slacks; the least are $8.1\cdot10^{-7}$ for \textup{(N13)} and
$1.3\cdot10^{-6}$ for \textup{(M1)} (as a logarithmic margin), \textup{(C4)} and \textup{(N1)}. (The margin of \textup{(E0)} is
$7\cdot10^{-11}$ by construction; it is the distance from $\eps_0$ to the slack of \textup{(C4)} and plays no role
in any threshold.) \textup{(N1)} and \textup{(N14)} are facts of exact decimal arithmetic:
$\zeta+1/G_0=2\cdot10^{-8}$.
Two consequences are used repeatedly. From \textup{(N3)} and $\tau_0<\tfrac12$: $\delta>\zeta=1/G_0$. From
\textup{(E0)} for \textup{(C6)}, and from \textup{(C7)}: $l<\tau_0-\eps_0$, $r<\tau_0-\eps_0$ and
$\mu(1-2a)<l<\tau_0-\eps_0$.

\paragraph{Consequences for the sizes.} We collect the facts about the integers $m,\sigma,k,v,\abs L,\abs R$ that
are used later; ``large $n$'' means $n\ge n_0$ for a threshold depending only on the constants above.
\begin{enumerate}[label=\textup{(S\arabic*)}]
\item $4k\le m$ (by \textup{(N4)}), hence $2k\le\sigma$; and $2\sigma=v+2k$.
\item $v<\abs L$ and $v\le\floor{n/2}$ for large $n$. Indeed $v=m-2k<(1-2a)m+2\le\mu(1-2a)n+2$ (using $k>am-1$
and $m\le\mu n$), while $\abs L\ge ln-1$ and
$\floor{n/2}\ge n/2-1$; use \textup{(C7)} and \textup{(N6)}.
\item $\binom\sigma k\binom vk\le2^{v}$ for all $n$ with $m\ge1$. Indeed $k/\sigma\le2a\le\tfrac12$ gives
$\binom\sigma k\le2^{(m/2)h(2a)}$, and $\binom vk\le\binom mk\le2^{mh(a)}$ because $v\le m$ and $k/m\le a$; and $m\bigl(\tfrac12h(2a)+h(a)\bigr)<m(1-2a)\le v$
by \textup{(N5)}.
\item $\binom\sigma k\ge2^{\gamma m}$ for large $n$. Indeed $\binom\sigma k\ge2^{\sigma h(k/\sigma)}/(\sigma+1)$,
$k/\sigma>2a-2/m$, and $h$ is concave with $h'(x)=\log\frac{1-x}{x}\le\log(1/a)$ for $x\ge a$; so
$\sigma h(k/\sigma)\ge\tfrac m2h(2a)-\log(1/a)$, and \textup{(C1)} gives the claim once
$m\bigl(\tfrac12h(2a)-\gamma\bigr)\ge\log(1/a)+\log(\sigma+1)$.
\item $\binom mk\le2^{h(a)m}$ and $\binom m{\sigma-k}\le2^{h(1/2-a)m+1}$; and $h(a)<h(\tfrac12-a)$ by
\textup{(N7)}. For the second bound, $(\sigma-k)/m\le\min\bigl(\tfrac12,\tfrac12-a+\tfrac1m\bigr)$, $h$ is increasing on
$[0,\tfrac12]$ with $h'\le1$ on $[\tfrac13,\tfrac12]$, and $\tfrac12-a\ge\tfrac13$ because $a\le\tfrac16$.
\item $\abs R\le rn+3$, $\abs L\le ln$, $\mu n-2<m\le\mu n$, and $\abs R\ge1$ for large $n$ by \textup{(C6)}.
\end{enumerate}
The entropy bounds $\binom Nj\le2^{Nh(j/N)}$ and $\binom Nj\ge2^{Nh(j/N)}/(N+1)$ used here are standard: the first
follows from $1=\sum_i\binom Ni x^i(1-x)^{N-i}\ge\binom Njx^j(1-x)^{N-j}$ at $x=j/N$, and the second because at
$x=j/N$ the $j$-th term is the largest of the $N+1$ terms.

\subsection{Differences from the sources, and explicit constants}\label{sec:sources}
Relative to the sources, this text makes the following changes, none of which alters a parameter or the exponent.
(i) The prime-collision bound is stated with an explicit constant and without the factor $L$
(Lemma~\ref{lem:primes}). (ii) The one-trial lemma is stated for the \emph{ideal} procedure, without cap and with
exact samplers; the capped procedure is compared with it afterwards (Section~\ref{sec:thmA}). (iii) The bad-set
argument needs $\delta>1/G_0$ in addition to \textup{(N3)}; both are recorded. (iv) With one filter the domination
step is written with the exact rounding costs $3p_1\eps$ and $p_2(4a_\Gamma+3)\eps$; the uniform bound
$6\max(p_1,p_2)\eps$ used by the certificate needs $a_\Gamma\le\tfrac34$, which is \textup{(N11)}. (v) All uniform
bounds for the filter functions are derived for densities with sine terms. (vi) The cutoff $n_*$ beyond which the
asymptotic argument applies is astronomically large; Section~\ref{sec:limits} says how large and why.

\section{Computational model, isolation and deficiency}\label{sec:model}

\subsection{The machine}
We use the word RAM of~\cite{OAI}. The word length is
\begin{equation}\label{eq:word}
 w=\bigl\lceil 4\bigl(n+b+\log(n+2)\bigr)\bigr\rceil .
\end{equation}
Each input integer occupies one word of read-only memory. Unit-cost operations are: reading and writing a word,
indirect addressing, comparison, addition, subtraction, multiplication, integer quotient and remainder by a nonzero
divisor, bitwise Boolean operations and shifts, each with a constant number of $w$-bit operands; and the generation
of one fresh uniformly random $w$-bit word. Arithmetic on longer integers is carried out by multiword routines and
every word operation is charged. Randomness is one-way: a random word can be used again only if it has been stored.
The algorithm is one finite program for all $n$ and $b$; nothing depending on $n$ is supplied to it.

\emph{Conventions.} A ``step on integers of $B$ bits'' costs $\poly(B)$ word operations. All integers handled below
have $O(b+n^2)$ bits, so every step costs $\poly(n+b)$ word operations; we say so where it matters.
The integers $\ceil{\tau_0n}$, $\floor{\eta n}$ and the like are computed exactly from the decimal parameters, and
the integers $\ceil{\sqrt n}$, $\ceil{2n^{0.6}}$ and $\ceil{n^{0.7}}$ used below are computed exactly in integer
arithmetic, as the least integers $e$ with $e^2\ge n$, $e^5\ge32n^3$ and $e^{10}\ge n^7$ respectively; this costs
$\poly(\log n)$ word operations.

\subsection{Tie weights and isolation}
Let $N_0$ be the power of two with $12n\le N_0<24n$. Draw $u_1,\dots,u_n$ independently and uniformly from
$\{1,\dots,N_0\}$ (each is one plus the value of $\log N_0$ fair bits). Put
\begin{equation}\label{eq:transform}
 C=1+nN_0,\qquad c_i=Ca_i+u_i,\qquad t_s=Ct+s\quad(0\le s\le nN_0).
\end{equation}
There are $nN_0+1\le24n^2+1$ targets $t_s$. Since $a_i,t<2^b$ and $C\le24n^2+1$, every $c_i$, every $t_s$ and every
signed sum $\sum_i\lambda_ic_i$ with $\lambda\in\{-1,0,1\}^n$ has absolute value below
\begin{equation}\label{eq:bits}
 2^{\,b+3\log n+6}.
\end{equation}

\begin{lemma}[isolation]\label{lem:iso}
\textup{(a)} For every value of the tie weights, some $I\subseteq[n]$ has $a(I)=t$ if and only if some
$s\in\{0,\dots,nN_0\}$ and some $I\subseteq[n]$ have $c(I)=t_s$; more precisely, $c(I)=t_s$ holds if and only if
$a(I)=t$ and $u(I)=s$.
\textup{(b)} If the instance has a solution, then with probability at least $11/12$ over the tie weights there is an
$s$ for which exactly one $I\subseteq[n]$ has $c(I)=t_s$.
\end{lemma}
\begin{proof}
(a) $c(I)-t_s=C\,(a(I)-t)+(u(I)-s)$. Both $u(I)$ and $s$ lie in $[0,nN_0]$, so $\abs{u(I)-s}\le nN_0<C$. Hence
$c(I)=t_s$ forces $a(I)=t$ and then $u(I)=s$; the converse is immediate, and every solution $I$ of the original
instance satisfies $c(I)=t_{u(I)}$ with $0\le u(I)\le nN_0$.

(b) Let $\mathcal F=\{I:a(I)=t\}\ne\varnothing$ and let $u_{\min}=\min_{I\in\mathcal F}u(I)$. Fix an index $i$ and the
values $u_j$, $j\ne i$. The least value of $u(I)$ over the members of $\mathcal F$ that contain $i$ is $\alpha_i+u_i$,
where $\alpha_i$ does not depend on $u_i$ (it is $+\infty$ if there is no such member); the least value over the
members that do not contain $i$ is a number $\beta_i$ not depending on $u_i$. Call $i$ \emph{ambiguous} if both are
finite and $\alpha_i+u_i=\beta_i$. For fixed $u_j$ ($j\ne i$) this happens for at most one value of $u_i$, so
$\Pr[i\text{ ambiguous}]\le1/N_0$. If two distinct members $I\ne I'$ of $\mathcal F$ both attain $u_{\min}$, take
$i\in I\triangle I'$; one of them contains $i$, the other does not, and both attain the overall minimum, so
$\alpha_i+u_i=u_{\min}=\beta_i$ and $i$ is ambiguous. By the union bound the minimum is attained by a unique set
$I^\ast$ with probability at least $1-n/N_0\ge11/12$. In that case let $s=u_{\min}$. By (a), the sets $I$ with
$c(I)=t_s$ are the members of $\mathcal F$ with $u(I)=u_{\min}$, that is, only $I^\ast$.
\end{proof}

\subsection{Deficiency}
Let $w$ be integer weights on a finite index set; signs and repetitions are allowed. For a finite set $V$ of
indices,
\begin{equation}\label{eq:def}
 \Sigma(V)=\{w(I):I\subseteq V\},\qquad d(V)=\abs V-\log\abs{\Sigma(V)} .
\end{equation}
So $\abs{\Sigma(V)}=2^{\abs V-d(V)}$: the deficiency $d(V)$ measures how far the subset sums of $V$ are from being
distinct.

\begin{lemma}[deficiency]\label{lem:def}
For arbitrary integer weights:
\begin{enumerate}[label=\textup{(\alph*)}]
\item $d(\varnothing)=0$ and $d(V)\ge0$.
\item If $i\notin V$ then $0\le d(V\cup\{i\})-d(V)\le1$. Consequently $d(U)\le d(V)$ and $\Sigma(U)\subseteq\Sigma(V)$
whenever $U\subseteq V$.
\item If $U\cap V=\varnothing$ then $d(U\cup V)\ge d(U)+d(V)$.
\item If $V'=(V\setminus\{i\})\cup\{j\}$ with $i\in V$, $j\notin V$, then $\abs{d(V)-d(V')}\le1$.
\item \textup{(signs)} Let $F$ be a set of indices and let $w^F$ be obtained from $w$ by negating $w_i$ for
$i\in F$. For $Z,F\subseteq V$,
\begin{equation}\label{eq:sign}
 w^F(Z)=w(Z\triangle F)-w(F).
\end{equation}
Hence, for any $F$, the set of subset sums of $V$ under $w^F$ is the translate $\Sigma(V)-w(F\cap V)$, and $d(V)$
is the same for $w$ and for $w^F$.
\end{enumerate}
\end{lemma}
\begin{proof}
(a) $\Sigma(\varnothing)=\{0\}$, and $1\le\abs{\Sigma(V)}\le2^{\abs V}$. (b) $\Sigma(V\cup\{i\})=\Sigma(V)\cup(w_i+\Sigma(V))$
has between $\abs{\Sigma(V)}$ and $2\abs{\Sigma(V)}$ elements, so $\log\abs\Sigma$ grows by an amount in $[0,1]$
while $\abs V$ grows by $1$. Monotonicity follows by adding the elements of $V\setminus U$ one at a time.
(c) $\Sigma(U\cup V)=\Sigma(U)+\Sigma(V)$ has at most $\abs{\Sigma(U)}\abs{\Sigma(V)}$ elements. (d) With
$U=V\setminus\{i\}$, both $d(V)$ and $d(V')$ lie in $[d(U),d(U)+1]$ by (b). (e) For $Z,F\subseteq V$,
$w(Z\triangle F)-w(F)=w(Z\setminus F)+w(F\setminus Z)-w(F)=w(Z\setminus F)-w(Z\cap F)=w^F(Z)$. For general $F$
replace $F$ by $F\cap V$, which does not change $w^F$ on $V$. As $Z\mapsto Z\triangle(F\cap V)$ is a bijection of the
subsets of $V$, the sums $w^F(Z)$, $Z\subseteq V$, are exactly the numbers $x-w(F\cap V)$, $x\in\Sigma(V)$.
\end{proof}

\subsection{Functions of a uniformly random subset}
\begin{lemma}[Azuma--Hoeffding]\label{lem:azuma}
Let $q\ge1$ and let $M_0,M_1,\dots,M_q$ be a martingale with $\abs{M_r-M_{r-1}}\le1$ for all $r$. Then for every
$\lambda>0$,
$\Pr[M_q-M_0\ge\lambda]\le e^{-\lambda^2/(2q)}$ and $\Pr[\abs{M_q-M_0}\ge\lambda]\le2e^{-\lambda^2/(2q)}$. Moreover
$\E(M_q-M_0)^2=\sum_{r}\E(M_r-M_{r-1})^2\le q$; in particular, if $M_0$ is deterministic, then $\Var M_q\le q$.
\end{lemma}
\begin{proof}
For $\abs y\le1$ and real $\theta$, convexity gives $e^{\theta y}\le\frac{1+y}2e^{\theta}+\frac{1-y}2e^{-\theta}$. If
$Y$ has conditional mean zero and $\abs Y\le1$ then $\E[e^{\theta Y}\mid\text{past}]\le\cosh\theta\le e^{\theta^2/2}$.
Applying this to the increments successively gives $\E e^{\theta(M_q-M_0)}\le e^{q\theta^2/2}$. By Markov's inequality
$\Pr[M_q-M_0\ge\lambda]\le e^{q\theta^2/2-\theta\lambda}$; take $\theta=\lambda/q$. The two-sided bound follows by
applying this to $-M$. Martingale increments are orthogonal in $L^2$, so
$\E(M_q-M_0)^2=\sum_r\E(M_r-M_{r-1})^2$. If $M_0$ is deterministic then $M_0=\E M_q$ and the left side is
$\Var M_q$. (Without that hypothesis the variance statement is false: a constant martingale $M_r=Z$ has zero
increments and variance $\Var Z$.)
\end{proof}

\begin{lemma}[exchange-Lipschitz functions]\label{lem:subsets}
Let $\Omega$ be a finite set, $0\le q\le\abs\Omega$, and let $g$ be a real function on the $q$-element subsets of
$\Omega$ with $\abs{g(V)-g(V')}\le1$ whenever $V'$ arises from $V$ by exchanging one element. Let $V_q$ be a uniformly
random $q$-subset of $\Omega$. Then $\Var g(V_q)\le q$, and for $q\ge1$
\[
 \Pr\bigl[g(V_q)-\E g(V_q)\ge\lambda\bigr]\le e^{-\lambda^2/(2q)},\qquad
 \Pr\bigl[\abs{g(V_q)-\E g(V_q)}\ge\lambda\bigr]\le2e^{-\lambda^2/(2q)}\quad(\lambda>0).
\]
For $q=0$ the set $V_0=\varnothing$ is deterministic and both probabilities are $0$.
\end{lemma}
\begin{proof}
For $q=0$ there is nothing to prove. Let $q\ge1$ and let
$X_1,\dots,X_q$ be a uniformly random ordered sample from $\Omega$ without replacement; then
$\{X_1,\dots,X_q\}$ is a uniform $q$-subset. Put $M_r=\E[g(\{X_1,\dots,X_q\})\mid X_1,\dots,X_r]$. This is a martingale
with the deterministic initial value $M_0=\E g(V_q)$ and with $M_q=g(V_q)$. Fix a prefix $(x_1,\dots,x_{r-1})$ and two admissible values $i\ne j$ of $X_r$.
Given $X_r=i$, the tail $(X_{r+1},\dots,X_q)$ is uniform over the ordered $(q-r)$-tuples of distinct elements of
$\Omega\setminus\{x_1,\dots,x_{r-1},i\}$. Applying the transposition $(i\,j)$ to each entry maps this law to the law
of the tail given $X_r=j$. Under this coupling the two final sets coincide if the tail contains $j$ (then both equal
the prefix together with $i$, $j$ and the rest of the tail), and otherwise they differ by exchanging $i$ for $j$.
Hence the two conditional expectations differ by at most $1$. So, given the prefix, the possible values of $M_r$
lie in an interval of length at most $1$ that contains their mean $M_{r-1}$; thus $\abs{M_r-M_{r-1}}\le1$. Apply
Lemma~\ref{lem:azuma}.
\end{proof}

\begin{lemma}[deficiency of random subsets]\label{lem:fmono}
Fix integer weights on a finite set $\Omega$ and let $V_q$ denote a uniform $q$-subset of $\Omega$.
\begin{enumerate}[label=\textup{(\alph*)}]
\item $\Var d(V_q)\le q$ and, for $q\ge1$ and $\lambda>0$,
$\Pr[\abs{d(V_q)-\E d(V_q)}\ge\lambda]\le2e^{-\lambda^2/(2q)}$.
\item $q\mapsto\E d(V_q)$ is nondecreasing.
\item A uniform $q'$-subset of a uniform $q$-subset of $\Omega$ ($q'\le q$) is a uniform $q'$-subset of $\Omega$. If
$M$ is a uniform $m$-subset of $\Omega$ and $L$ is a uniform $\ell$-subset of $\Omega\setminus M$, then $L$ is a
uniform $\ell$-subset of $\Omega$.
\end{enumerate}
\end{lemma}
\begin{proof}
(a) By Lemma~\ref{lem:def}(d), $d$ is exchange-Lipschitz; apply Lemma~\ref{lem:subsets}. (b) Let $\pi$ be a uniform
random ordering of $\Omega$. Its first $q$ entries form a uniform $q$-subset contained in the uniform
$(q+1)$-subset formed by its first $q+1$ entries; apply Lemma~\ref{lem:def}(b) and take expectations. (c) In both
cases the law of the resulting set is invariant under every permutation of $\Omega$, and the only
permutation-invariant law on the subsets of a given size is the uniform one.
\end{proof}

From now on, when the tie weights are fixed, $d$ denotes the deficiency for the weights $c_1,\dots,c_n$ and
\begin{equation}\label{eq:f}
 f(q)=\E\,d(V_q)\qquad(V_q\text{ a uniform }q\text{-subset of }[n]).
\end{equation}
\section{Random primes, bounded samplers and finite-precision sums}\label{sec:prims}

\subsection{Collisions modulo a random prime}
\begin{lemma}[random primes]\label{lem:primes}
Let $L\ge5$ be an integer and let $p$ be uniform among the primes of $[2^L,2^{L+1})$.
\begin{enumerate}[label=\textup{(\alph*)}]
\item The interval contains more than $\tfrac35\,2^L/(L\ln2)$ primes; so a uniform integer of the interval is prime
with probability greater than $c_\ast/L$, where $c_\ast=3/(5\ln2)>0.86$.
\item If $D$ is a nonzero integer with $\abs D<2^B$, $B>0$, then $\Pr[p\mid D]\le\dfrac{B}{c_\ast2^{L}}<1.16\,B\,2^{-L}$.
\end{enumerate}
\end{lemma}
\begin{proof}
(a) Rosser and Schoenfeld~\cite[Corollary~3]{RS}: $\pi(2x)-\pi(x)>3x/(5\ln x)$ for $x\ge20.5$. Take $x=2^L\ge32$.
Neither $2^L$ nor $2^{L+1}$ is prime, so the primes of $(2^L,2^{L+1}]$ are those of $[2^L,2^{L+1})$. (b) Let
$p_1<\dots<p_j$ be the distinct primes of the interval that divide $D$. Their product divides $D$, so
$2^{Lj}\le\abs D<2^B$ and $j<B/L$. Divide by the number of primes from (a).
\end{proof}
The sources state the bound as $K\,B\,L\,2^{-L}$ with unspecified $K$; the proof above shows that the factor $L$ is
not needed, because only prime divisors of size at least $2^L$ are counted.

\subsection{Bounded samplers}\label{sec:samplers}
The analysis uses exactly uniform draws; the machine has random words. We use the following bounded procedures, all
of which either return a sample with exactly the intended law or report failure.

\begin{lemma}[samplers]\label{lem:samplers}
Let $n\ge2$.
\begin{enumerate}[label=\textup{(\alph*)}]
\item \emph{Uniform integer.} For $1\le N\le2^{n^2}$ there is a procedure that returns an element of
$\{0,\dots,N-1\}$ or fails; it fails with probability at most $2^{-n^3}$, and conditional on not failing its output
is uniform. It uses $\poly(n)$ word operations.
\item \emph{Uniform subset, Bernoulli.} A uniform $j$-subset of a set of size $N\le n$ and a Bernoulli variable of
rational parameter $P_1/P_2$ ($P_2\le 2^{n^2}$) are obtained from (a), with failure probability at most $n2^{-n^3}$
and $2^{-n^3}$ respectively, and exact conditional laws.
\item \emph{Uniform prime.} For $5\le L\le n$ there is a procedure that returns a prime of $[2^L,2^{L+1})$ or fails; it
fails with probability at most $e^{-c_\ast n^2}+n^32^{-n^3}$; conditional on not failing its output is uniform among
the primes of the interval; and it uses at most $n^3\cdot2^{(L+1)/2}$ trial divisions on $(L+1)$-bit integers, plus
$\poly(n)$ word operations.
\end{enumerate}
\end{lemma}
\begin{proof}
(a) Let $\ell=\ceil{\log N}$. Draw $\ell$ fresh fair bits (low bits of random words), read them as an integer $x$,
accept if $x<N$, otherwise repeat, at most $n^3$ times; for $N=1$ return $0$. Each attempt accepts with probability
$N/2^\ell>\tfrac12$ and, given acceptance, $x$ is uniform on $\{0,\dots,N-1\}$; the attempts are independent. So the
failure probability is below $2^{-n^3}$ and the conditional law is uniform. (b) Choose the elements one after
another, each uniform among the remaining ones, by (a): the probability of not failing at each stage does not depend
on the element chosen, so conditional on no failure the set is uniform. For the Bernoulli variable test a uniform
element of $\{0,\dots,P_2-1\}$ against $P_1$. (c) Repeat at most $n^3$ times: draw a uniform integer of
$[2^L,2^{L+1})$ by (a) and test it for primality by dividing by $2,3,\dots,\floor{\sqrt{\cdot}}$ (the integer square
root is found by binary search); stop at the first prime. Conditional on the draws not failing, each candidate is
uniform on the interval and independent of the earlier ones, so the first prime found is uniform among the primes of
the interval, and by Lemma~\ref{lem:primes}(a) all $n^3$ candidates are composite with probability at most
$(1-c_\ast/L)^{n^3}\le e^{-c_\ast n^2}$ because $L\le n$.
\end{proof}

\paragraph{Coupling with ideal draws.} Call a run of an algorithm \emph{ideal} if every call of
Lemma~\ref{lem:samplers} is replaced by an exact draw from the intended law. A real run and the ideal run can be
defined on one probability space so that they coincide unless some sampler call of the real run fails. If the real
run makes at most $2^{Kn}$ calls of types (a),(b) and at most $2^{Kn}$ calls of type (c) (the calls may be chosen
adaptively; the failure bound of each call holds conditionally on the past), the two runs differ with probability at
most
\begin{equation}\label{eq:couple}
 2^{Kn}\bigl(n2^{-n^3}+e^{-c_\ast n^2}+n^32^{-n^3}\bigr)=o(1),
\end{equation}
a quantity depending only on $n$ and $K$.

\subsection{Finite-precision complex sums}\label{sec:precision}
Write $e_P(x)=\exp(2\pi\mathrm i\,x/P)$. The algorithm never needs exact complex numbers.

\begin{lemma}[phases]\label{lem:phase}
There is a deterministic rational-arithmetic procedure which, given integers $P\ge2$, $x$ and $Q\ge1$, returns a
complex number with dyadic real and imaginary parts within $2^{-Q}$ of $e_P(x)$, in time polynomial in $Q$ and in the
bit lengths of $P$ and $x$.
\end{lemma}
\begin{proof}
Reduce $x$ modulo $P$. Compute a rational approximation of $\pi$ from
$\pi/4=\arctan\tfrac12+\arctan\tfrac13$ and the alternating series $\arctan y=\sum_{j\ge0}(-1)^jy^{2j+1}/(2j+1)$,
whose remainder after $J$ terms is below $y^{2J+1}$; hence the angle $\theta=2\pi x/P\in[0,2\pi)$ is obtained to
within $2^{-Q-4}$ with $O(Q)$ terms. Evaluate $\cos$ and $\sin$ by their Taylor series at the approximate angle
$\tilde\theta\in[0,7]$: the remainder after the term of degree $J$ is at most $7^{J+1}/(J+1)!$, which is below
$2^{-Q-4}$ once $J\ge 40+Q$: then $(J+1)!\ge((J+1)/e)^{J+1}$ and $7e/(J+1)<\tfrac12$. Since $\cos$ and $\sin$ are $1$-Lipschitz, the angle error contributes
at most $2^{-Q-4}$; a final rounding to $Q+3$ fractional bits contributes at most $2^{-Q-3}$. All intermediate
rationals have numerators and denominators of $\poly(Q+\log P)$ bits.
\end{proof}

\begin{lemma}[products and sums]\label{lem:sums}
Let $n\ge8$ and $\delta=2^{-2n^2}$.
\begin{enumerate}[label=\textup{(\alph*)}]
\item If $z_0,\dots,z_j$ ($j\le n$) are complex numbers with $\abs{z_i}\le2$ and $\abs{\tilde z_i-z_i}\le\delta$, then
$\bigl\lvert\prod_i\tilde z_i-\prod_iz_i\bigr\rvert\le(n+1)\,\delta\,(2+\delta)^{n}\le2^{-2n^2+n+\log(n+1)+1}$.
\item Let $X=\sum_{i=1}^{N}\omega_i\zeta_i$ with $N\le2^{Kn}$, exact rational weights $\abs{\omega_i}\le2^{Kn}$, and
complex $\zeta_i$ each known to within $2^{-2n^2+2n}$. If each product $\omega_i\tilde\zeta_i$ is rounded to $n^2$
fractional bits before the additions, the computed sum $\tilde X$ satisfies
$\abs{\tilde X-X}\le2^{Kn}\bigl(2^{Kn}2^{-2n^2+2n}+2^{-n^2+1}\bigr)\le2^{-n^2+(2K+3)n}$.
\end{enumerate}
In particular the error in \textup{(b)} is at most $1/16$ as soon as $n\ge2K+4$.
\end{lemma}
\begin{proof}
(a) Telescoping,
$\prod\tilde z_i-\prod z_i=\sum_{i}\bigl(\prod_{i'<i}\tilde z_{i'}\bigr)(\tilde z_i-z_i)\bigl(\prod_{i'>i}z_{i'}\bigr)$,
and each of the $j+1\le n+1$ terms is at most $\delta(2+\delta)^{n}$. (b) Each term has error at most
$2^{Kn}2^{-2n^2+2n}$ before rounding, and rounding the real and imaginary parts to $n^2$ fractional bits adds at
most $2^{-n^2+1}$. Summing over $N\le2^{Kn}$ terms gives at most $2^{-2n^2+(2K+2)n}+2^{-n^2+Kn+1}$, and each of
these two numbers is at most $2^{-n^2+(2K+3)n-1}$. The last sentence is
arithmetic: $-n^2+(2K+3)n\le-4$ when $n\ge2K+4$.
\end{proof}
Dyadic numbers with $n^2$ fractional bits and absolute value at most $2^{2Kn}$ are added exactly with $n^2+O(n)$
bits. Every comparison that the algorithm makes on approximate quantities is made on these stored dyadic values or
on the fixed rational approximations of Section~\ref{sec:filter}, so repeated evaluations give identical results.
\section{The preliminary search}\label{sec:prelim}
Fix the tie weights. The following procedure is meet in the middle on \emph{distinct} sums, with a cap.

\begin{procedure}[capped preliminary search]\label{proc:prelim}
\begin{enumerate}[label=\arabic*.]
\item Draw a uniform $\floor{n/2}$-subset $V$ of $[n]$ by Lemma~\ref{lem:samplers}(b). If the sampler fails, return
\textsc{continue}.
\item Build the sorted list of the distinct subset sums of $V$: start from the list $(0)$ and, for each $i\in V$ in
turn, merge the current sorted list with its translate by $c_i$, deleting duplicates. If a merge is about to produce
more than $B_n=2^{\ceil{\tau_0n}}$ distinct values, stop and return \textsc{continue}. Do the same for
$[n]\setminus V$.
\item For each $s=0,\dots,nN_0$, scan the first list upward and the second downward to decide whether some pair of
entries sums to $t_s$. Return \textsc{yes} if this happens for some $s$, and \textsc{no} otherwise.
\end{enumerate}
\end{procedure}

\begin{proposition}[preliminary search]\label{prop:P}
For every value of the tie weights:
\begin{enumerate}[label=\textup{(\alph*)}]
\item if Procedure~\ref{proc:prelim} returns \textsc{yes} or \textsc{no}, that is the correct answer for the
instance $(a,t)$;
\item for $n\ge8$, on every execution it performs at most $50\,n^3\,2^{\tau_0n}$ list operations (comparisons, additions and
moves of integers of at most $b+3\log n+7$ bits), plus $\poly(n)$ operations for the sampler; hence at most
$2^{\tau_0n}\poly(n+b+2)$ word operations;
\item if
\begin{equation}\label{eq:notLD}
 \E\,d(V)>\bigl(\tfrac12-\tau_0+\zeta\bigr)n\qquad(V\text{ a uniform }\floor{n/2}\text{-subset of }[n])
\end{equation}
and $n>1/\zeta$, then $\Pr[\textsc{continue}]\le\dfrac{8}{\zeta^2n}+n2^{-n^3}$.
\end{enumerate}
\end{proposition}
\begin{proof}
(a) If the procedure reaches Step~3, the two lists are $\Sigma(V)$ and $\Sigma([n]\setminus V)$ in full, because the
intermediate lists are the sets $\Sigma(V')$ for prefixes $V'\subseteq V$. A pair of entries with sum $t_s$ exists if
and only if some $I\subseteq[n]$ has $c(I)=t_s$. By Lemma~\ref{lem:iso}(a) such $s$ and $I$ exist if and only if the
original instance has a solution.

(b) There are $n$ merges; a merge of a list of at most $B_n$ entries with its translate takes at most $2B_n+2$
comparisons and moves before it completes or is interrupted. Each of the at most $24n^2+1$ scans takes at most
$2B_n$ steps. With $B_n\le2\cdot2^{\tau_0n}$ the total is at most $(2n+48n^2+2)(2B_n+2)\le50n^3\,2^{\tau_0n}$ for
$n\ge8$. Entries are subset sums of the $c_i$, bounded by~\eqref{eq:bits}, and their translates.

(c) Let $U$ be either $V$ or $[n]\setminus V$; it is a uniform subset of $[n]$ of size $\floor{n/2}$ or
$\ceil{n/2}$. By Lemma~\ref{lem:fmono}(b) and~\eqref{eq:notLD}, $\E d(U)\ge f(\floor{n/2})>(\tfrac12-\tau_0+\zeta)n$,
and by Lemma~\ref{lem:fmono}(a), $\Var d(U)\le n$. Chebyshev's inequality gives
\[
 \Pr\bigl[d(U)<(\tfrac12-\tau_0+\tfrac\zeta2)n\bigr]\le\frac{n}{(\zeta n/2)^2}=\frac4{\zeta^2n}.
\]
If this event occurs for neither side then, for each side,
$\log\abs{\Sigma(U)}=\abs U-d(U)\le\tfrac n2+\tfrac12-(\tfrac12-\tau_0+\tfrac\zeta2)n=\tau_0n+\tfrac12-\tfrac{\zeta n}2<\tau_0n$
when $n>1/\zeta$. By Lemma~\ref{lem:def}(b) no intermediate list is longer than the final one, so no merge is
interrupted and the procedure returns \textsc{yes} or \textsc{no}. Add the sampler's failure probability.
\end{proof}

\begin{definition}[the low-deficiency condition]\label{def:LD}
For fixed tie weights, \textup{(LD)} is the statement
\begin{equation}\label{eq:LD}
 f(\floor{n/2})=\E\,d(V)\le\bigl(\tfrac12-\tau_0+\zeta\bigr)n .\tag{LD}
\end{equation}
\end{definition}
So, for each fixed choice of tie weights, either \textup{(LD)} holds or the preliminary search returns a correct
answer except with probability at most $8/(\zeta^2n)+n2^{-n^3}$. The algorithm never evaluates $f$, and the outcome
\textsc{continue} is \emph{not} evidence that \textup{(LD)} holds; \textup{(LD)} is used only inside proofs.
With $\zeta=10^{-8}$ the bound in (c) is below $1/100$ only for $n\ge8\cdot10^{18}$; this is one of the reasons why
the cutoff $n_\ast$ of the final algorithm is so large (Section~\ref{sec:limits}).
\section{Small-layer estimates}\label{sec:layers}
In this section $w$ denotes arbitrary integer weights on a finite index set (equal weights and both signs allowed).
For a finite set $V$ of indices and $j\ge0$ put
\[
 \Sigma_j(V)=\{w(X):X\subseteq V,\ \abs X=j\},\qquad a_j(V)=\abs{\Sigma_j(V)}\Big/\binom{\abs V}{j}\quad(0\le j\le\abs V).
\]
Thus $a_j(V)\in(0,1]$ is the fraction of distinct values among the $j$-element subset sums, and
$\Sigma(V)=\bigcup_j\Sigma_j(V)$.

\begin{lemma}[covering]\label{lem:cover}
Let $0\le k\le t\le\abs V$ and let $W$ be a uniform $(t-k)$-subset of $V$. Then $a_t(V)\le\E_W\,a_k(V\setminus W)$.
\end{lemma}
\begin{proof}
For each $x\in\Sigma_t(V)$ choose one $t$-set $X_x\subseteq V$ with $w(X_x)=x$. Let $\Pi$ be the set of pairs
$(x,W)$ with $x\in\Sigma_t(V)$, $W\subseteq X_x$, $\abs W=t-k$; so $\abs\Pi=\binom tk\abs{\Sigma_t(V)}$. The map
$(x,W)\mapsto(W,\,x-w(W))$ is injective on $\Pi$, since $x$ is recovered as $(x-w(W))+w(W)$. Its second component
is $w(X_x\setminus W)$, the sum of a $k$-subset of $V\setminus W$. Hence
$\binom tk\abs{\Sigma_t(V)}\le\sum_{W}\abs{\Sigma_k(V\setminus W)}$, the sum over all $(t-k)$-subsets $W$ of $V$. With
$N=\abs V$, divide both sides by $\binom Nt\binom tk=\binom N{t-k}\binom{N-t+k}{k}$ (both equal
$N!/((N-t)!\,k!\,(t-k)!)$).
\end{proof}

\begin{lemma}[sampling]\label{lem:sample}
Let $k\le s\le\abs V$ and let $S$ be a uniform $s$-subset of $V$. Then $\E_S\,a_k(S)\ge a_k(V)$.
\end{lemma}
\begin{proof}
For each $x\in\Sigma_k(V)$ choose a $k$-set $X_x\subseteq V$ with $w(X_x)=x$. Then $x\in\Sigma_k(S)$ whenever
$X_x\subseteq S$, an event of probability $\binom{N-k}{s-k}/\binom Ns=\binom sk/\binom Nk$ ($N=\abs V$). Summing over
$x$, $\E\abs{\Sigma_k(S)}\ge\abs{\Sigma_k(V)}\binom sk/\binom Nk$; divide by $\binom sk$.
\end{proof}

\begin{lemma}[a good layer]\label{lem:good}
For every $V$ there is $t$ with $2t\le\abs V$ and $\abs{\Sigma(V)}\le(\abs V+1)\abs{\Sigma_t(V)}$; in particular
$a_t(V)\ge2^{-d(V)}/(\abs V+1)$.
\end{lemma}
\begin{proof}
The $\abs V+1$ layers cover $\Sigma(V)$, so some $j$ has $\abs{\Sigma_j(V)}\ge\abs{\Sigma(V)}/(\abs V+1)$.
Complementation $X\mapsto V\setminus X$ gives $\Sigma_{\abs V-j}(V)=w(V)-\Sigma_j(V)$, so $t=\min(j,\abs V-j)$ has the
same property and $2t\le\abs V$. Finally $\binom{\abs V}{t}\le2^{\abs V}$ and $\abs{\Sigma(V)}=2^{\abs V-d(V)}$.
\end{proof}

\begin{corollary}[averaged small layer; Corollary~D of~\cite{Base}]\label{cor:D}
Let $M$ be a finite index set with integer weights and let $k,\sigma,v$ be nonnegative integers with
\[
 2k\le\sigma,\qquad 2\sigma\le v+2k,\qquad v\le\abs M,\qquad \binom\sigma k\binom vk\le2^{v}.
\]
Let $S$ be a uniform $\sigma$-subset and $V$ a uniform $v$-subset of $M$. Then
\[
 \E_S\,a_k(S)\ \ge\ \frac{\E_V\,2^{-d(V)}}{v+1}.
\]
\end{corollary}
\begin{proof}
From $2k\le\sigma\le v/2+k$ we get $k\le v/2$ and $k\le\sigma\le v$. Fix a $v$-set $V\subseteq M$ and let $S_V$ be a
uniform $\sigma$-subset of $V$. We show
\begin{equation}\label{eq:Dfixed}
 \frac{2^{-d(V)}}{v+1}\le\E\,a_k(S_V).
\end{equation}
Take $t$ from Lemma~\ref{lem:good} for $V$; so $2t\le v$ and $a_t(V)\ge2^{-d(V)}/(v+1)$.

\emph{Case $t\ge k$.} By Lemma~\ref{lem:cover}, $a_t(V)\le\E_W\,a_k(V\setminus W)$ with $W$ a uniform
$(t-k)$-subset of $V$. Here $\abs{V\setminus W}=v-t+k\ge v/2+k\ge\sigma\ge k$, so Lemma~\ref{lem:sample} applies to
$V\setminus W$ with $s=\sigma$: $a_k(V\setminus W)\le\E\,a_k(S')$ for a uniform $\sigma$-subset $S'$ of
$V\setminus W$. The law of $S'$, with $W$ uniform, is invariant under the permutations of $V$, so $S'$ is a uniform
$\sigma$-subset of $V$. Hence $2^{-d(V)}/(v+1)\le a_t(V)\le\E\,a_k(S_V)$.

\emph{Case $t<k$.} Then $t<k\le v/2$, so $\abs{\Sigma_t(V)}\le\binom vt\le\binom vk$. By Lemma~\ref{lem:good},
$2^{v-d(V)}=\abs{\Sigma(V)}\le(v+1)\binom vk$, and with the last hypothesis
$2^{-d(V)}/(v+1)\le\binom vk2^{-v}\le1/\binom\sigma k$. Every $\sigma$-set $S$ has $a_k(S)\ge1/\binom\sigma k$ because
$\Sigma_k(S)\ne\varnothing$ ($k\le\sigma$). So~\eqref{eq:Dfixed} holds.

Now average~\eqref{eq:Dfixed} over a uniform $v$-subset $V$ of $M$: by Lemma~\ref{lem:fmono}(c), $S_V$ is then a
uniform $\sigma$-subset of $M$.
\end{proof}

\begin{remark}\label{rem:signs}
Negating some weights leaves $d$ unchanged (Lemma~\ref{lem:def}(e)) but can change $a_k$. In Section~\ref{sec:alias}
the corollary is applied to one signing of the block, fixed before the random set $S$ is drawn. The hypothesis
$\binom\sigma k\binom vk\le2^v$ is used only in the case $t<k$ and cannot be dropped from this proof. The
corollary controls layers of size $k$ inside a random half-size set through the deficiency of sets of size $v$; it
does not assert that small layers of the whole block are robust (the open statement \textup{[S]} of~\cite{V}), and
nothing of that kind is used in this paper.
\end{remark}
\section{Alias phase sums and Theorem A}\label{sec:alias}
Fix the tie weights, hence $c_1,\dots,c_n$ and the targets $t_s$. Let $P$ be a prime of
$[2^{\floor{\eta n}},2^{\floor{\eta n}+1})$ and put
\[
 \bc_i=c_i\bmod P\in\{0,\dots,P-1\},\qquad \bc(I)=\sum_{i\in I}\bc_i\ \text{(an integer, not reduced)}.
\]
Phases $\varphi\in(\Z/P\Z)^n$ are given, and $\chi(I)=e_P\bigl(\sum_{i\in I}\varphi_i\bigr)$. The \emph{aliases} of
$t_s$ are
\[
 \cA_s(P)=\{I\subseteq[n]:\ c(I)\equiv t_s\!\!\pmod P,\ c(I)\ne t_s\}.
\]
If $I\in\cA_s(P)$ then $\bc(I)\equiv t_s\pmod P$ and $0\le\bc(I)\le n(P-1)$, so there is a unique \emph{carry}
$j\in\{0,\dots,n-1\}$ with
\begin{equation}\label{eq:carry}
 \bc(I)=T_{s,j}:=(t_s\bmod P)+jP .
\end{equation}

\subsection{Statement}
\begin{thmA}[alias phase sums]
There are a randomized procedure and an integer $n_A$ with the following properties for all $n\ge n_A$. The input
is the transformed weights and targets~\eqref{eq:transform} of an instance that passes the guard, a prime
$P\in[2^{\floor{\eta n}},2^{\floor{\eta n}+1})$ and phases $\varphi\in(\Z/P\Z)^n$. The output is \textsc{fail} or
complex numbers $\hat A_s$, $0\le s\le nN_0$.
\begin{enumerate}[label=\textup{(\arabic*)}]
\item Assume \textup{(LD)}. Let $P$ be uniform among the primes of its interval and independent of the procedure's
random words. There is an event $\mathcal H$, determined by $P$ and the procedure's random words (for the given
weights and targets) and not involving $\varphi$, with $\Pr[\mathcal H]\ge1-\epsilon_A(n)$, where
$\epsilon_A(n)\to0$ depends only on $n$, such that on $\mathcal H$ the following holds \emph{simultaneously for
every} $\varphi\in(\Z/P\Z)^n$: the procedure run with the phases $\varphi$ does not fail and
\[
 \Bigl\lvert\hat A_s-\sum_{I\in\cA_s(P)}\chi(I)\Bigr\rvert\le\frac1{16}\qquad\text{for every }s .
\]
In particular the conclusion holds with probability at least $1-\epsilon_A(n)$ for phases chosen in any way
whatever; they may depend on $P$. (The event that the output is accurate does depend on $\varphi$ in general, for
instance through cancellation among aliases that were missed; $\mathcal H$ is a sufficient event that does not.)
\item On every execution the procedure uses at most $2^{\tau_0n+2n^{0.7}}\poly(n+b+2)$ word operations.
\end{enumerate}
\end{thmA}

\subsection{Signed weights, keys and classes}
The partition $[n]=M\cup L\cup R$ is chosen below. For $F\subseteq M$ and $i\in M$ put
$c^F_i=(-1)^{[i\in F]}c_i$ and $\bc^F_i=(-1)^{[i\in F]}\bc_i$. By~\eqref{eq:sign}, for all $Z,F\subseteq M$,
\begin{equation}\label{eq:signcF}
 c^F(Z)=c(Z\triangle F)-c(F),\qquad \bc^F(Z)=\bc(Z\triangle F)-\bc(F).
\end{equation}
The \emph{key} of a set $A$ is the integer pair $\key(A)=(c(A),\bc(A))$. For a key $\xi=(x,y)$ write
$c(\xi)=x$, $\bc(\xi)=y$; note $y\equiv x\pmod P$. A \emph{key of $L$} is a key $\key(A)$ with $A\subseteq L$; likewise
for $R$. Define
\[
 \Phi_L(\xi)=\sum_{A\subseteq L:\ \key(A)=\xi}\chi(A),\qquad \Phi_R(\xi)=\sum_{B\subseteq R:\ \key(B)=\xi}\chi(B).
\]
For keys $\xi_L$ of $L$, $\xi_R$ of $R$ and $Z\subseteq M$, the \emph{class} named $(\xi_L,\xi_R,Z)$ is
\[
 \cC(\xi_L,\xi_R,Z)=\{A\cup Z\cup B:\ A\subseteq L,\ B\subseteq R,\ \key(A)=\xi_L,\ \key(B)=\xi_R\}.
\]
Each $I\subseteq[n]$ lies in exactly one class, namely that of $(\key(I\cap L),\key(I\cap R),I\cap M)$. All members
of a class have the same $c(I)$ and the same $\bc(I)$; so a class consists entirely of aliases of $t_s$ (an
\emph{alias class} of $t_s$) or contains none. Since $\chi$ is multiplicative over disjoint unions,
\begin{equation}\label{eq:classsum}
 \sum_{I\in\cC(\xi_L,\xi_R,Z)}\chi(I)=\Phi_L(\xi_L)\,\Phi_R(\xi_R)\,\chi(Z).
\end{equation}

\subsection{The procedure}\label{sec:procA}
We first describe the \emph{ideal} procedure: exact uniform draws, exact complex arithmetic, no cap.
Let $e_n$ be the integer $\ceil{2n^{0.6}}$.

\begin{procedure}[ideal alias procedure]\label{proc:A}
\begin{enumerate}[label=\arabic*.]
\item \emph{Partition.} Draw a uniform $m$-subset $M$ of $[n]$, then a uniform $\abs L$-subset $L$ of $[n]\setminus M$;
put $R=[n]\setminus(M\cup L)$.
\item \emph{Value maps.} Enumerate all $A\subseteq L$ with $\key(A)$ and $\chi(A)$; sort by key; store each distinct
key $\xi$ once together with $\Phi_L(\xi)$. Do the same for $R$.
\item \emph{Trial count.} Draw $n^2$ independent uniform $v$-subsets $V_1,\dots,V_{n^2}$ of $M$ and compute each
$\abs{\Sigma(V_i)}$ exactly, by the merging of Procedure~\ref{proc:prelim} without a cap. Put
\[
 \hat S=\frac1{n^2}\sum_{i=1}^{n^2}2^{-d(V_i)},\qquad \hat F=\log\frac1{\hat S}+e_n,\qquad
 \Ntr=n^7\,2^{e_n}\Bigl\lceil\frac1{\hat S}\Bigr\rceil .
\]
\item \emph{Trials.} For every target index $s$ and every carry $j\in\{0,\dots,n-1\}$ run $\Ntr$ independent trials.
A trial for $(s,j)$:
\begin{enumerate}[label=(\alph*)]
\item draw a uniform $F\subseteq M$; put $T^F=T_{s,j}-\bc(F)$ and $t^F=t_s-c(F)$;
\item draw a uniform prime $P'$ of $[2^{\floor{\gamma m}},2^{\floor{\gamma m}+1})$ and a uniform $y\in\Z/P'\Z$;
\item \emph{left records}: all pairs $(\xi_L,X)$ with $\xi_L$ a key of $L$, $X\subseteq M$, $\abs X=k$ and
$\bc(\xi_L)+\bc^F(X)\equiv y\pmod{P'}$. (Sort the keys of $L$ by $\bc(\xi_L)\bmod P'$; for each $X$ find the keys of
the required residue by binary search.)
\item \emph{right classes}: list all $Y\subseteq M$ with $\abs Y=\sigma-k$ and group them by
$\lambda(Y)=(\bc^F(Y),c^F(Y))$; a right class is a value $\lambda=(\bc_\lambda,c_\lambda)$ that occurs, with the list
of its members $Y$;
\item \emph{right class records}: all pairs $(\xi_R,\lambda)$ with $\xi_R$ a key of $R$, $\lambda$ a right class and
$\bc(\xi_R)+\bc_\lambda\equiv T^F-y\pmod{P'}$ (listed as in (c));
\item \emph{matching}: group the right class records by the integer $\bc(\xi_R)+\bc_\lambda$ and sort each group by
the integer $c(\xi_R)+c_\lambda$. For each left record $(\xi_L,X)$ find the group with
$\bc(\xi_R)+\bc_\lambda=T^F-\bc(\xi_L)-\bc^F(X)$; in it find, by two binary searches, the interval of records with
$c(\xi_R)+c_\lambda=t^F-c(\xi_L)-c^F(X)$; skip that interval and \emph{visit} every other record of the group;
\item \emph{expansion}: for each visited $(\xi_R,\lambda)$ and each member $Y$ of $\lambda$, test $X\cap Y=\varnothing$;
if the test passes, put $Z=(X\cup Y)\triangle F$ and insert $(\xi_L,\xi_R,Z)$ into the dictionary $\cD_s$ unless it
is already there.
\end{enumerate}
\item \emph{Output.} $A^\circ_s=\sum_{(\xi_L,\xi_R,Z)\in\cD_s}\Phi_L(\xi_L)\Phi_R(\xi_R)\chi(Z)$ for each $s$.
\end{enumerate}
\end{procedure}
Since $1/\hat S\ge1$ we have $1/\hat S\le\ceil{1/\hat S}\le2/\hat S$, so
\begin{equation}\label{eq:Ntr}
 n^7\,2^{\hat F}\le\Ntr\le2n^7\,2^{\hat F}.
\end{equation}
An \emph{elementary step} is one of: a subset enumerated, a merge operation, a mask listed, a key or record
compared or moved during sorting or binary search, a record formed, a record visited, a member expanded, a
dictionary operation, a term added, a trial division, a sampler call. Each elementary step costs $\poly(n+b)$ word
operations: keys have at most $b+3\log n+7$ and $\eta n+\log n+2$ bits, masks have $m$ bits, dictionaries are
balanced search trees with at most $2^n$ entries, and the complex numbers are the dyadic approximations of
Section~\ref{sec:precision}. A sampler call of type (a) or (b) of Lemma~\ref{lem:samplers} is one elementary step;
a call of type (c) for a prime of $[2^L,2^{L+1})$ consists of at most $n^3$ candidates, each of which is one
sampler call of type (a) followed by at most $2^{(L+1)/2}$ trial divisions.

\paragraph{Charging convention for the ideal procedure.} The ideal procedure has no samplers, so the number of
elementary steps of its exact draws has to be fixed by convention. We \emph{charge} every exact draw of the ideal
procedure the largest number of elementary steps that the corresponding bounded sampler can perform on any
execution: one step for a uniform integer, a uniform subset or a Bernoulli variable, and
\begin{equation}\label{eq:charge}
 n^3\bigl(1+2^{(L+1)/2}\bigr)\ \text{steps for a uniform prime of }[2^L,2^{L+1}).
\end{equation}
All other operations of the ideal procedure are counted as they are performed. The resulting number is the
\emph{charged step count} of the ideal run; ``the number of elementary steps of the ideal procedure'' always means
this number. The convention is chosen so that the following holds on every outcome. Couple a run of the capped
procedure (Procedure~\ref{proc:Acap} below) with the ideal run as in Section~\ref{sec:samplers}. As long as no sampler of the capped run has failed and its cap has not
intervened, the capped run has made the same draws and performed the same operations as the ideal run, and each of
its sampler calls has cost at most the charge of the corresponding ideal draw. Hence, at every moment of a capped
run in which no sampler has failed,
\begin{equation}\label{eq:dominate}
 \#\{\text{steps performed so far by the capped run}\}\ \le\ \#\{\text{charged steps of the ideal run}\}.
\end{equation}
In particular, if no sampler fails and the charged step count of the ideal run is at most the cap, the capped run
never reaches its cap and builds the same dictionaries as the ideal run.

\begin{procedure}[capped alias procedure]\label{proc:Acap}
Run Procedure~\ref{proc:A} with the bounded samplers of Lemma~\ref{lem:samplers} and with the arithmetic of
Lemmas~\ref{lem:phase}--\ref{lem:sums} (each phase to within $2^{-2n^2}$, each summed term rounded to $n^2$
fractional bits), and count elementary steps. Return \textsc{fail} as soon as the count would exceed
$2^{\ceil{\tau_0n}+\ceil{n^{0.7}}}$, or if a sampler fails. Otherwise return the computed values $\hat A_s$.
\end{procedure}
Part~(2) of Theorem~A is immediate: at most $2^{\tau_0n+n^{0.7}+2}$ elementary steps, each of
$\poly(n+b)$ word operations. The rest of the section proves part~(1).

\subsection{Soundness and counts}
\begin{lemma}[soundness]\label{lem:A1}
In any run (ideal or capped), every triple inserted into $\cD_s$ names an alias class of $t_s$. Consequently, if at
the end $\cD_s$ contains the names of all alias classes of $t_s$, then $A^\circ_s=\sum_{I\in\cA_s(P)}\chi(I)$.
\end{lemma}
\begin{proof}
An insertion arises from a left record $(\xi_L,X)$, a visited record $(\xi_R,\lambda)$ and a member $Y$ of $\lambda$
with $X\cap Y=\varnothing$; and $Z=(X\cup Y)\triangle F$. Let $A\subseteq L$, $B\subseteq R$ with $\key(A)=\xi_L$,
$\key(B)=\xi_R$, and $I=A\cup Z\cup B$. By~\eqref{eq:signcF} and disjointness,
$\bc^F(X)+\bc^F(Y)=\bc^F(X\cup Y)=\bc(Z)-\bc(F)$, and likewise for $c$. The record $(\xi_R,\lambda)$ lies in the group
of $(\xi_L,X)$, so $\bc(A)+\bc(B)+\bc(Z)-\bc(F)=T^F=T_{s,j}-\bc(F)$; hence $\bc(I)=T_{s,j}$ and
$c(I)\equiv\bc(I)\equiv t_s\pmod P$. The record is outside the skipped interval, so
$c(A)+c(B)+c(Z)-c(F)\ne t^F=t_s-c(F)$, that is $c(I)\ne t_s$. So $I$ is an alias of $t_s$, and then so is every
member of its class. For the second statement: distinct names are distinct, hence disjoint, classes; the alias
classes of $t_s$ partition $\cA_s(P)$; apply~\eqref{eq:classsum}.
\end{proof}

\begin{lemma}[counts]\label{lem:A2}
\textup{(a)} $L$ has at most $(n+1)\abs{\Sigma(L)}=(n+1)2^{\abs L-d(L)}$ keys, and each value of $c(\xi)$ occurs in
at most $n+1$ keys of $L$; the same holds for $R$. \textup{(b)} In every trial there are at most
$(2m+1)\,2^{m-d(M)}$ right classes.
\end{lemma}
\begin{proof}
(a) A key $(x,y)$ of $L$ has $x\in\Sigma(L)$, $y\equiv x\pmod P$ and $0\le y\le\abs L(P-1)$, which leaves at most
$\abs L+1\le n+1$ values of $y$ for each $x$. (b) By~\eqref{eq:signcF}, $c^F(Y)=c(Y\triangle F)-c(F)$ takes at most
$\abs{\Sigma(M)}=2^{m-d(M)}$ values; for each of them $\bc^F(Y)\equiv c^F(Y)\pmod P$ and $\abs{\bc^F(Y)}<mP$, which
leaves at most $2m+1$ values.
\end{proof}

\subsection{The trial count}
For the block $M$ let $V$ denote a uniform $v$-subset of $M$ and put
\[
 g_M=-\log\E_V\,2^{-d(V)},\qquad e_M=\E_V\,d(V).
\]
By Jensen's inequality $g_M\le e_M$, and by Lemma~\ref{lem:def}(b) $e_M\le d(M)$.

\begin{lemma}[trial count]\label{lem:A3}
Let $n$ be large enough for \textup{(S2)}. With probability at least $1-\epsilon_3(n)$ over $M$, $L$ and Step~3,
where $\epsilon_3(n)=2^{-2n^{0.6}}+(2n^2+3)e^{-n^{0.2}/2}$, all of the following hold.
\begin{enumerate}[label=\textup{(\roman*)}]
\item $\hat F\ge g_M$.
\item $\hat F\le e_M+3n^{0.6}+1\le d(M)+3n^{0.6}+1$.
\item $\hat F\le d(L)+5n^{0.6}+1$.
\item If \textup{(LD)} holds: $g_M\le e_M\le(\tfrac12-\tau_0+\zeta)n+n^{0.6}$ and
$\hat F\le(\tfrac12-\tau_0+\zeta)n+4n^{0.6}+1$.
\end{enumerate}
Call this event $\cE_{\mathrm{tr}}$. It is determined by $M$, $L$ and the draws of Step~3; in particular it is
independent of $P$, of the phases and of the trials.
\end{lemma}
\begin{proof}
(i) $\hat S$ is the mean of $n^2$ independent copies of $2^{-d(V)}$, so $\E[\hat S\mid M]=2^{-g_M}$. By Markov's
inequality, $\hat S<2^{2n^{0.6}}2^{-g_M}$ except with probability $2^{-2n^{0.6}}$; then
$\hat F=\log(1/\hat S)+e_n>g_M-2n^{0.6}+e_n\ge g_M$.

(ii) Fix $M$ and apply Lemma~\ref{lem:fmono}(a) with $\Omega=M$, $q=v\le n$, $\lambda=n^{0.6}$: each $d(V_i)$ is at
most $e_M+n^{0.6}$ except with probability $2e^{-n^{1.2}/(2v)}\le2e^{-n^{0.2}/2}$. If this holds for all $n^2$
samples then $\hat S\ge2^{-e_M-n^{0.6}}$ and $\hat F\le e_M+n^{0.6}+e_n\le e_M+3n^{0.6}+1$.

Next, $e_M\le f(v)+n^{0.6}$ except with probability $e^{-n^{0.2}/2}$. Indeed $e_M$ is a function of the uniform
$m$-subset $M$ of $[n]$, and it is exchange-Lipschitz: if $M'=(M\setminus\{i\})\cup\{i'\}$, couple a uniform
$v$-subset $V$ of $M$ with $V'=V$ if $i\notin V$ and $V'=(V\setminus\{i\})\cup\{i'\}$ otherwise; then $V'$ is a
uniform $v$-subset of $M'$ and $\abs{d(V)-d(V')}\le1$ by Lemma~\ref{lem:def}(d), so $\abs{e_M-e_{M'}}\le1$. By
Lemma~\ref{lem:fmono}(c), $\E_M\,e_M=f(v)$. Apply the one-sided bound of Lemma~\ref{lem:azuma} through the
martingale of Lemma~\ref{lem:subsets} with $q=m\le n$.

(iii) By Lemma~\ref{lem:fmono}(c), $L$ is a uniform $\abs L$-subset of $[n]$, so $d(L)\ge f(\abs L)-n^{0.6}$ except
with probability $2e^{-n^{0.2}/2}$. By \textup{(S2)}, $v<\abs L$, and $f$ is nondecreasing; hence
$\hat F\le e_M+3n^{0.6}+1\le f(v)+4n^{0.6}+1\le f(\abs L)+4n^{0.6}+1\le d(L)+5n^{0.6}+1$.

(iv) By \textup{(S2)}, $v\le\floor{n/2}$, so $f(v)\le f(\floor{n/2})\le(\tfrac12-\tau_0+\zeta)n$ under \textup{(LD)}.
Combine with the bounds above. Adding the exceptional probabilities gives $\epsilon_3(n)$.
\end{proof}

\subsection{Reduction modulo $P$ on half-blocks}
Call a $\sigma$-subset $S\subseteq M$ \emph{bad for $P$} if there is $\lambda\in\{-1,0,1\}^S$ with
$\sum_{i\in S}\lambda_ic_i\ne0$ and $P\mid\sum_{i\in S}\lambda_ic_i$; otherwise $S$ is \emph{good}. Recall
$\delta=\tfrac12(\eta-\tfrac\mu2\log3)>0$ by \textup{(C8)}.

\begin{lemma}[bad half-blocks]\label{lem:A4}
Let $M$ be fixed and let $P$ be uniform among the primes of $[2^{\floor{\eta n}},2^{\floor{\eta n}+1})$, with
$\floor{\eta n}\ge5$. With probability at least $1-\epsilon_4(n)$, where
$\epsilon_4(n)=2.4\,(b+3\log n+6)\,2^{-\delta n}$, at most a $2^{-\delta n}$ fraction of the $\sigma$-subsets of $M$ are
bad for $P$. Call this event $\cE_P$. Under the guard, $\epsilon_4(n)\le 2^{(1/G_0-\delta)n+O(\log n)}\to0$, because
$\delta>1/G_0$.

If $S$ is good, $\epsilon_i\in\{\pm1\}$ for $i\in S$ and $X,X'\subseteq S$, then
\begin{equation}\label{eq:keep}
 \sum_{i\in X}\epsilon_ic_i\ne\sum_{i\in X'}\epsilon_ic_i\ \Longrightarrow\ \sum_{i\in X}\epsilon_i\bc_i\ne\sum_{i\in X'}\epsilon_i\bc_i .
\end{equation}
\end{lemma}
\begin{proof}
Fix $S$. There are fewer than $3^\sigma$ vectors $\lambda$, and each nonzero $\sum_i\lambda_ic_i$ has absolute value
below $2^{B}$ with $B=b+3\log n+6$ by~\eqref{eq:bits}. By Lemma~\ref{lem:primes}(b) with $L=\floor{\eta n}$, $P$
divides a given one with probability at most $1.16\,B\,2^{-\floor{\eta n}}\le2.4\,B\,2^{-\eta n}$. So
$\Pr[S\text{ bad}]\le2.4\,B\,3^{\sigma}2^{-\eta n}\le2.4\,B\,2^{-2\delta n}$, using $\sigma\le\mu n/2$. The expected
fraction of bad sets has the same bound, and Markov's inequality gives the first claim. For the guard: $b\le2^{n/G_0}$.

For~\eqref{eq:keep}: the difference of the two left sides is $\sum_{i\in S}\lambda_ic_i$ with
$\lambda_i=\epsilon_i([i\in X]-[i\in X'])\in\{-1,0,1\}$; it is nonzero, so $P$ does not divide it. The difference of
the two right sides is congruent to it modulo $P$, hence nonzero.
\end{proof}

\subsection{Coverage}
\begin{lemma}[one trial]\label{lem:A5}
There is $n_5$ such that the following holds for $n\ge n_5$. Assume \textup{(LD)}. Fix $P$, $M$, $L$ and the
outcome of Step~3 such that $\cE_{\mathrm{tr}}$ and $\cE_P$ hold. Fix $s$, an alias $I\in\cA_s(P)$ with carry $j$, and
one trial of the ideal procedure for $(s,j)$. Over the trial's own draws $(F,P',y)$, the trial inserts the name of
the class of $I$ into $\cD_s$ (or finds it already there) with probability at least $2^{-g_M}n^{-6}$.
\end{lemma}
\begin{proof}
Put $A=I\cap L$, $B=I\cap R$ and $S=(I\cap M)\triangle F$.

\emph{Step 1 (a half of the block).} $F$ is uniform among the subsets of $M$, hence so is $S$. So
$\Pr[\abs S=\sigma]=\binom m\sigma2^{-m}\ge1/(m+1)$, the middle binomial coefficient being the largest of $m+1$
coefficients that sum to $2^m$. Conditional on $\abs S=\sigma$, $S$ is a uniform $\sigma$-subset of $M$, and
$F=S\triangle(I\cap M)$.

\emph{Step 2 (the signs on $S$ do not depend on $S$).} Define weights on $M$ by $c^I_i=c_i$ for $i\in I\cap M$ and
$c^I_i=-c_i$ for $i\in M\setminus I$; define $\bc^I$ likewise. If $i\in S$, then either $i\in I$ and $i\notin F$, or
$i\notin I$ and $i\in F$; in both cases $c^F_i=c^I_i$ and $\bc^F_i=\bc^I_i$. The weights $c^I$ depend on $I$ only.

\emph{Step 3 (many distinct layer values).} Apply Corollary~\ref{cor:D} to the weights $c^I$ on $M$ with the
integers $k,\sigma,v$ of Section~\ref{sec:params}; its hypotheses are \textup{(S1)} and \textup{(S3)} and
$v\le m$. By Lemma~\ref{lem:def}(e) the deficiency of every $V\subseteq M$ is the same for $c^I$ and for $c$. Hence,
writing $a_k(S;c^I)$ for the layer density under $c^I$,
\[
 \E_S\,a_k(S;c^I)\ \ge\ \frac{\E_V2^{-d(V)}}{v+1}=\frac{2^{-g_M}}{v+1}.
\]
Let $W_S=\{\bc(A)+\bc^F(X):X\subseteq S,\ \abs X=k\}$, a set of integers, and $\beta_S=\abs{W_S}$. By Step~2,
$W_S=\{\bc(A)+\bc^I(X)\}$ depends on $S$ (and on $I$, $P$, $L$) but not on $P'$ or $y$. If $S$ is good,
\eqref{eq:keep} with $\epsilon_i=c^I_i/c_i$ shows that $k$-subsets of $S$ with different $c^I$-sums have different
$\bc^I$-sums, so $\beta_S\ge\abs{\Sigma_k(S;c^I)}=a_k(S;c^I)\binom\sigma k$. On $\cE_P$ a uniform $\sigma$-subset is
bad with probability at most $2^{-\delta n}$, and $a_k\le1$, so
\[
 \E_S\Bigl[\mathbf 1_{S\text{ good}}\,\frac{\beta_S}{\binom\sigma k}\Bigr]\ \ge\ \frac{2^{-g_M}}{v+1}-2^{-\delta n}
 \ \ge\ \frac{2^{-g_M}}{2(v+1)} .
\]
The last step needs $2^{-\delta n}\le2^{-g_M}/(2(v+1))$. By Lemma~\ref{lem:A3}(iv),
$g_M\le(\tfrac12-\tau_0+\zeta)n+n^{0.6}$, so it holds as soon as
$\bigl(\delta-(\tfrac12-\tau_0+\zeta)\bigr)n\ge n^{0.6}+\log(2n+2)$, which is true for large $n$ by \textup{(N3)}.

\emph{Step 4 (the small prime keeps many residues).} Fix a good $S$. The elements of $W_S$ lie in $(-nP,nP)$, so
the difference of two of them is a nonzero integer of absolute value below $2nP<2^{\eta n+\log n+2}\le2^{n}$ (large
$n$). By Lemma~\ref{lem:primes}(b) with $L=\floor{\gamma m}\ge5$, for distinct $x,x'\in W_S$,
$\Pr_{P'}[P'\mid x-x']\le1.16\,n\,2^{-\floor{\gamma m}}\le3n\,2^{-\gamma m}$. Let $z_\rho$ be the number of elements of
$W_S$ in the residue class $\rho$ modulo $P'$. Counting ordered pairs,
\[
 \E_{P'}\sum_\rho z_\rho^2\ \le\ \beta_S+3n\,\beta_S^2\,2^{-\gamma m}\ \le\ 4n\max\bigl(\beta_S,\ \beta_S^2\,2^{-\gamma m}\bigr).
\]
By Markov's inequality, with probability at least $\tfrac12$ over $P'$ the sum $\sum_\rho z_\rho^2$ is at most twice
this bound. On that event the Cauchy--Schwarz inequality
$\beta_S^2=(\sum_\rho z_\rho)^2\le\#\{\rho:z_\rho>0\}\cdot\sum_\rho z_\rho^2$ shows that at least
$\min(\beta_S,2^{\gamma m})/(8n)$ residues are occupied. Since $P'<2^{\gamma m+1}$, a uniform $y\in\Z/P'\Z$ is occupied
with probability at least
\[
 \frac{\min(\beta_S2^{-\gamma m},1)}{16n}\ \ge\ \frac{\beta_S}{16n\binom\sigma k},
\]
where the last step uses both $\beta_S\le\binom\sigma k$ and $2^{\gamma m}\le\binom\sigma k$ \textup{(S4)}.

\emph{Step 5 (an occupied residue is recovered).} Suppose $\abs S=\sigma$ and $y\equiv\bc(A)+\bc^F(X)\pmod{P'}$ for
some $X\subseteq S$ with $\abs X=k$. Then $(\key(A),X)$ is a left record. Put $Y=S\setminus X$, so $\abs Y=\sigma-k$
and $X\cap Y=\varnothing$. Since $S\triangle F=I\cap M$, \eqref{eq:signcF} gives
\[
 \bc(A)+\bc^F(X)+\bc(B)+\bc^F(Y)=\bc(A)+\bc(B)+\bc(I\cap M)-\bc(F)=T_{s,j}-\bc(F)=T^F .
\]
Hence $\bc(B)+\bc^F(Y)\equiv T^F-y\pmod{P'}$, so $(\key(B),\lambda(Y))$ is a right class record, and it lies in the
group of the left record $(\key(A),X)$. The same computation with $c$ gives $c(A)+c^F(X)+c(B)+c^F(Y)=c(I)-c(F)$,
which differs from $t^F=t_s-c(F)$ because $I$ is an alias. So the record is outside the skipped interval and is
visited; the expansion reaches its member $Y$, the test $X\cap Y=\varnothing$ passes, and the triple
$(\key(A),\key(B),(X\cup Y)\triangle F)=(\key(A),\key(B),I\cap M)$, the name of the class of $I$, is inserted unless
already present.

\emph{Step 6 (product).} $F$ is independent of the pair $(P',y)$, and given $P'$ the residue $y$ is uniform on
$\Z/P'\Z$; Step~4 used exactly this, first over $P'$ and then over $y$. Steps 1, 3, 4 and 5 give a success probability of at
least
\[
 \frac1{m+1}\ \E_S\Bigl[\mathbf 1_{S\text{ good}}\cdot\frac12\cdot\frac{\beta_S}{16n\binom\sigma k}\Bigr]
 \ \ge\ \frac{2^{-g_M}}{64\,n\,(m+1)(v+1)}\ \ge\ 2^{-g_M}\,n^{-6}
\]
for $n\ge8$. The thresholds on $n$ (Steps~3 and~4, \textup{(S1)--(S4)}) depend only on the fixed constants.
\end{proof}

\begin{corollary}[coverage]\label{cor:A6}
Assume \textup{(LD)}, and fix $P$, $M$, $L$ and the outcome of Step~3 such that $\cE_{\mathrm{tr}}$ and $\cE_P$ hold,
$n\ge n_5$. Then, except with probability at most $(24n^2+1)(2/e)^n$ over the trials, at the end of Step~4 of the
ideal procedure every alias class of every target $t_s$ is named in $\cD_s$.
\end{corollary}
\begin{proof}
Fix $s$ and an alias $I$ of $t_s$ with carry $j$. The $\Ntr$ trials for $(s,j)$ are independent, and
$\Ntr\ge n^72^{\hat F}\ge n^72^{g_M}$ by~\eqref{eq:Ntr} and Lemma~\ref{lem:A3}(i). By Lemma~\ref{lem:A5}, none of them
produces the class of $I$ with probability at most $(1-2^{-g_M}n^{-6})^{\Ntr}\le e^{-n}$. There are at most
$(nN_0+1)2^n$ pairs $(s,I)$.
\end{proof}

\subsection{Expected work}\label{sec:work}
This is the step where an accounting error is most likely, so each category of elementary step is treated
separately and the random choices averaged over are named each time.

\begin{lemma}[expected work]\label{lem:A7}
Assume \textup{(LD)}. Fix $M$, $L$ and the outcome of Step~3 such that $\cE_{\mathrm{tr}}$ holds. Let $P$ be uniform
among the primes of its interval (independent of $M$, $L$, Step~3), let the phases be arbitrary, and let $n$ be
large. Then the expected number of elementary steps of the ideal procedure (its charged step count, in the sense
of Section~\ref{sec:procA}), over $P$ and over all draws of the trials, is at most
\[
 (b+2)\,2^{(\tau_0-\eps_0+\zeta)n+c_7n^{0.6}}
\]
for an absolute constant $c_7$.
\end{lemma}
\begin{proof}
The number of trials is $N_{\mathrm{all}}=(nN_0+1)\,n\,\Ntr\le50\,n^{10}\,2^{\hat F}$ by~\eqref{eq:Ntr}; it is
determined by Step~3. Below, $\pi_n$ denotes a factor $2^{O(n^{0.6})}$ with an absolute implied constant; it absorbs
all polynomial factors in $n$ and $N_{\mathrm{all}}2^{-\hat F}$. We use $m\le\mu n$, $m>\mu n-2$,
$\abs L\le ln$, $\abs R\le rn+3$ \textup{(S6)}, and the four bounds on $\hat F$ of Lemma~\ref{lem:A3}. Each of
\textup{(C2)--(C5)} is used in the form given by \textup{(E0)}: its left side is less than $\tau_0-\eps_0$.

\emph{1.\ Value maps and Step~3 (no randomness needed).} Enumerating and sorting the subsets of $L$ and of $R$
takes $O(n(2^{\abs L}+2^{\abs R}))$ steps; Step~3 takes at most $n^2\cdot v\cdot(2^{v+1}+2)$ merge operations. The
exponents $l$, $r$ and $\mu(1-2a)$ are below $\tau_0-\eps_0$ (Section~\ref{sec:params}).

\emph{2.\ Masks (deterministic per trial).} A trial lists $\binom mk+\binom m{\sigma-k}\le2^{h(1/2-a)m+2}$ masks
\textup{(S5)}, computes $\bc^F$ and $c^F$ of each, and sorts the right masks into classes. Over all trials this is
$\pi_n\,2^{\hat F+h(1/2-a)\mu n}\le\pi_n\,2^{(h(1/2-a)\mu+\frac12-\tau_0+\zeta)n}$ by Lemma~\ref{lem:A3}(iv), and
$h(\tfrac12-a)\mu+\tfrac12-\tau_0\le\tau_0-\eps_0$ by \textup{(C5)}.

\emph{3.\ Small primes (deterministic charge per trial).} The exact draw of $P'$ is charged~\eqref{eq:charge}
with $L=\floor{\gamma m}$: $n^3$ candidates, each with at most $2^{(\gamma m+1)/2}$ trial divisions. This is a
convention about the ideal procedure, not a property of an exact draw; by~\eqref{eq:dominate} it dominates what the
bounded sampler of the capped run actually does. Over all trials: $\pi_n2^{(\frac12-\tau_0+\zeta+\gamma\mu/2)n}$,
and the exponent is
below $\tau_0-\eps_0+\zeta$ by \textup{(N8)}. The other draws of a trial (the set $F$, the residue $y$) are charged
one step each, and the $n^3$ candidate draws for $P'$ are included in~\eqref{eq:charge}; the draws of Steps~1
and~3 number $n^2+2$ in all.

\emph{4.\ Sorting keys by residue (deterministic per trial).} A trial sorts at most
$(n+1)(2^{\abs L-d(L)}+2^{\abs R})$ keys (Lemma~\ref{lem:A2}(a)). With Lemma~\ref{lem:A3}(iii) for the first term and
Lemma~\ref{lem:A3}(iv) for the second, the total is at most
$\pi_n\bigl(2^{ln}+2^{(r+\frac12-\tau_0+\zeta)n}\bigr)$; use $l<\tau_0-\eps_0$ and \textup{(N8)}.

\emph{5.\ Left records (average over $y$).} Fix everything in a trial except $y$. A pair $(\xi_L,X)$ is a left
record for exactly one residue $y$, hence with probability $1/P'\le2^{1-\gamma m}$. So the expected number of left
records in a trial is at most $(n+1)2^{\abs L-d(L)}\binom mk2^{1-\gamma m}$, and over all trials at most
$\pi_n\,2^{\hat F-d(L)}\,2^{\abs L+(h(a)-\gamma)m}\le\pi_n\,2^{(l+(h(a)-\gamma)\mu)n}$, by Lemma~\ref{lem:A3}(iii),
\textup{(S5)} and $\mu n-2<m\le\mu n$ (which gives $(h(a)-\gamma)m\le(h(a)-\gamma)\mu n+2$ whatever the sign of
$h(a)-\gamma$). Use \textup{(C2)}.
The binary searches of Step~(c) cost $O(n)$ steps per mask and per record and are covered by items 2 and 5.

\emph{6.\ Right class records (average over $y$).} By Lemma~\ref{lem:A2}(b) a trial has at most
$(2m+1)2^{m-d(M)}$ right classes; handling the classes (one binary search each) costs
$\pi_n2^{\hat F+m-d(M)}\le\pi_n2^{\mu n}$ over all trials by Lemma~\ref{lem:A3}(ii), and $\mu<\tau_0-\eps_0$ by
\textup{(N2)}. A pair $(\xi_R,\lambda)$ is a right class record for exactly one $y$; so the expected number of right
class records in a trial is at most $(n+1)2^{\abs R}(2m+1)2^{m-d(M)}2^{1-\gamma m}$, and over all trials at most
$\pi_n\,2^{\hat F-d(M)}\,2^{\abs R+(1-\gamma)m}\le\pi_n\,2^{(r+(1-\gamma)\mu)n}$. Use \textup{(C3)}.

\emph{7.\ Visits and expansions (average over $P$ and $y$).} Call a tuple $(\xi_L,X,\xi_R,Y)$ \emph{visited} in a
trial if $(\xi_L,X)$ is a left record, $(\xi_R,\lambda(Y))$ is a right class record visited from it, and $Y$ is a
member of $\lambda(Y)$. Every expansion step handles one visited tuple, and every visited record has at least one
member, so the visits and expansions of a trial number at most twice its visited tuples.

Fix a trial's $s$, $j$, $F$ and $P'$. For values $x_L\in\Sigma(L)$, $x_R\in\Sigma(R)$ and masks $X$, $Y$ of sizes
$k$ and $\sigma-k$ put
\[
 D=D(x_L,X,x_R,Y)=x_L+c^F(X)+x_R+c^F(Y)-t^F .
\]
$D$ is an integer that depends on neither $P$ nor $y$ (the sets $\Sigma(L)$, $\Sigma(R)$ and the integers $c_i$,
$t_s$ do not involve $P$), and $\abs D<2^{B}$ with $B=b+3\log n+9$ by~\eqref{eq:bits}. Suppose a tuple with
$c(\xi_L)=x_L$ and $c(\xi_R)=x_R$ is visited. Then: (i) $D\ne0$, because the skipped interval consists exactly of
the records with $D=0$; (ii) $P\mid D$, because the group condition says
$\bc(\xi_L)+\bc^F(X)+\bc(\xi_R)+\bc^F(Y)=T^F$, while $\bc(\xi)\equiv c(\xi)$, $\bc^F\equiv c^F$ and
$T^F\equiv t^F\pmod P$; (iii) $y\equiv\bc(\xi_L)+\bc^F(X)\pmod{P'}$, because $(\xi_L,X)$ is a left record. The set
of keys depends on $P$, but for every $P$ at most $(n+1)^2$ pairs $(\xi_L,\xi_R)$ have $c(\xi_L)=x_L$ and
$c(\xi_R)=x_R$ (Lemma~\ref{lem:A2}(a)), and for each such pair condition (iii) holds for one residue $y$. As $y$ is
uniform and independent of $P$, the expected number of visited tuples with given $(x_L,X,x_R,Y)$ is at most
\[
 \mathbf 1_{D\ne0}\cdot\Pr_P[P\mid D]\cdot(n+1)^2\cdot\frac1{P'}
 \ \le\ 1.16\,B\,2^{-\floor{\eta n}}\,(n+1)^2\,2^{1-\gamma m},
\]
by Lemma~\ref{lem:primes}(b). Summing over $(x_L,X,x_R,Y)$, the expected number of visited tuples in one trial is
at most
\[
 5B(n+1)^2\;2^{\abs L-d(L)}\,2^{\abs R}\binom mk\binom m{\sigma-k}\,2^{-\eta n-\gamma m}.
\]
Multiply by $N_{\mathrm{all}}\le50n^{10}2^{\hat F}$ and use Lemma~\ref{lem:A3}(iii), $\abs L+\abs R=n-m$ and
\textup{(S5)}: the total is at most
\[
 B\,\pi_n\,2^{\,n-m+(h(a)+h(1/2-a)-\gamma)m-\eta n}\ \le\ B\,\pi_n\,2^{(1-\mu+(h(a)+h(1/2-a)-\gamma)\mu-\eta)n},
\]
where the last step writes the exponent as $(1-\eta)n+\vartheta m$ with
$\vartheta=h(a)+h(\tfrac12-a)-\gamma-1\in[-2,2]$ and uses $\mu n-2<m\le\mu n$, so that $\vartheta m\le\vartheta\mu n+4$
whatever the sign of $\vartheta$. By \textup{(C4)} with \textup{(E0)} the exponent is at most $(\tau_0-\eps_0)n$.
Finally $B\le(b+2)(3\log n+9)$.

\emph{8.\ Matching, dictionaries, output.} The matching step performs $O(n)$ elementary steps per left record, per
right class record and per visited tuple; each expansion causes at most one dictionary operation of $O(n)$ steps;
the output sums one term per dictionary entry, and there are at most as many entries as insertions.

Every item is at most $(b+2)\,\pi_n\,2^{(\tau_0-\eps_0+\zeta)n}$.
\end{proof}

\begin{remark}\label{rem:A7}
Item~7 averages over $P$ under its unconditional law; it does not condition on $\cE_P$, and it does not use
$\cE_P$. Only coverage uses $\cE_P$. Items 2--4 are worst-case bounds given $\cE_{\mathrm{tr}}$ and
\textup{(LD)}; items 5--6 average only over $y$. The hypothesis \textup{(LD)} enters through
Lemma~\ref{lem:A3}(iv) in items 2, 3 and 4; without it the trial count $\Ntr$ is not controlled. The lemma bounds an
expectation; it is turned into a statement about the capped procedure only through Markov's inequality in the
proof below, never into a worst-case time bound.
\end{remark}

\subsection{Proof of Theorem~A}\label{sec:thmA}
Part~(2) was shown after Procedure~\ref{proc:Acap}. For part~(1) assume \textup{(LD)} and $n\ge n_A$, where $n_A$
exceeds $n_5$ and the thresholds of \textup{(S1)--(S6)}. Couple the capped procedure with the ideal procedure
(Section~\ref{sec:samplers}). By the charging convention and~\eqref{eq:dominate}, they build the same dictionaries
unless a sampler of the capped run fails or the charged step count of the ideal procedure exceeds
$2^{\ceil{\tau_0n}+\ceil{n^{0.7}}}$. Consider the following events.
\begin{itemize}
\item $\cE_{\mathrm{tr}}$ fails: probability at most $\epsilon_3(n)$ (Lemma~\ref{lem:A3}).
\item $\cE_P$ fails: probability at most $\epsilon_4(n)$ (Lemma~\ref{lem:A4}; $M$ is independent of $P$).
\item $\cE_{\mathrm{tr}}$ and $\cE_P$ hold but some alias class is missing from the ideal dictionaries:
probability at most $(24n^2+1)(2/e)^n$ (Corollary~\ref{cor:A6}).
\item $\cE_{\mathrm{tr}}$ holds and the ideal procedure exceeds the step cap: by Lemma~\ref{lem:A7} and Markov's
inequality, conditionally on $(M,L,\text{Step 3})\in\cE_{\mathrm{tr}}$ this has probability at most
\[
 \frac{(b+2)\,2^{(\tau_0-\eps_0+\zeta)n+c_7n^{0.6}}}{2^{\tau_0n+n^{0.7}}}\ \le\ 2^{(1/G_0+\zeta-\eps_0)n+c_7n^{0.6}-n^{0.7}+1},
\]
using the guard. By \textup{(N1)} the coefficient of $n$ is negative, so this tends to $0$.
\item A sampler of the capped run fails: the capped run makes at most $2^{n}$ sampler calls, so by~\eqref{eq:couple}
this has probability $o(1)$.
\end{itemize}
All five bounds depend only on $n$ (the guard bounds $b$); let $\epsilon_A(n)$ be their sum and let $\mathcal H$
be the complement of the union of the five events.

\emph{$\mathcal H$ does not involve the phases.} The control flow of both procedures never reads $\varphi$: the
phases enter only the stored values $\Phi_L(\xi)$, $\Phi_R(\xi)$, the factors $\chi(Z)$ and the output sums, while
the partition, the trial count, the records, the visits, the dictionaries, every sampler call and the step count
are functions of the weights, the targets, $P$ and the random words. On the event that no sampler fails, the
coupled ideal draws are the outputs of the samplers, so the ideal run is itself a function of $P$ and the random
words. Hence $\mathcal H$ is determined by $P$ and the random words, and $\Pr[\mathcal H]\ge1-\epsilon_A(n)$.

\emph{On $\mathcal H$ the output is accurate for every $\varphi$.} On $\mathcal H$ the capped procedure does not
fail, and its dictionary
$\cD_s$ contains every alias class of $t_s$ and nothing else (Lemma~\ref{lem:A1}); so the exact value of its output
sum is $\sum_{I\in\cA_s(P)}\chi(I)$, whatever the phases are. It remains to bound the arithmetic error, uniformly
in $\varphi$.

\emph{Arithmetic error.} Each $\chi(A)$ is evaluated to within $2^{-2n^2}$ (Lemma~\ref{lem:phase}) and rounded to
$n^2$ fractional bits, so each computed $\tilde\Phi_L(\xi)$ and $\tilde\Phi_R(\xi)$ is within
$\epsilon'=2^{n}(2^{-2n^2}+2^{-n^2+1})\le2^{-n^2+n+2}$ of the true value, and $\abs{\Phi_L},\abs{\Phi_R}\le2^n$. A
product of three factors of modulus at most $2^n+1$, each known to within $\epsilon'$, is known to within
$3\epsilon'(2^n+1)^2$. Each dictionary has at most $2^n$ entries (they name disjoint nonempty classes), and each term
is rounded to $n^2$ fractional bits. The total error is at most
$2^n\bigl(3\epsilon'(2^n+1)^2+2^{-n^2+1}\bigr)\le2^{-n^2+4n+7}\le\tfrac1{16}$ for $n\ge8$. \qed
\section{Filter construction and finite-circle convergence}\label{sec:filter}
Everything in this section is fixed independently of the input.

\subsection{Definitions on a circle}
A \emph{circle} is either $\Tc=\R/2\pi\Z$ with normalized Lebesgue measure, or, for an integer $P\ge2$, the finite
circle $\Z/P\Z$ with normalized counting measure, the residue $x$ being identified with the angle $2\pi x/P$. On a
circle, $\av H$ denotes the mean of $H$ (an integral with the factor $1/2\pi$, or $P^{-1}\sum_x$).

Fix real trigonometric polynomials with cosine and sine terms,
\[
 r(x)=\sum_{j=1}^{J}\bigl(r_j\cos jx+r'_j\sin jx\bigr),\qquad s(x)=\sum_{j=1}^{J}\bigl(s_j\cos jx+s'_j\sin jx\bigr),
\]
with rational coefficients, and a rational $\kappa>0$. Write
\[
 R_r=\sum_j(\abs{r_j}+\abs{r'_j}),\quad R^1_r=\sum_jj(\abs{r_j}+\abs{r'_j}),\quad R=R_r+R_s,\quad R^1=R^1_r+R^1_s,\quad \kappa^+=\max(1,\kappa),
\]
so that $\abs r\le R_r$, $\abs{r'}\le R^1_r$, and likewise for $s$. On either circle put
\begin{equation}\label{eq:fg}
\begin{gathered}
 f_\ast(y)=\frac{2^{\kappa r(y)}}{\av{2^{\kappa r}}},\qquad g_\ast(y)=\frac{2^{\kappa s(y)}}{\av{2^{\kappa s}}},\\
 k(x,w)=2^{\,r(x+w)+s(w)},\qquad q(x)=\av{k(x,\cdot)},\qquad K_x(w)=\frac{k(x,w)}{q(x)},
\end{gathered}
\end{equation}
and
\begin{equation}\label{eq:DAE}
\begin{gathered}
 D(x)=\av{K_x\log K_x},\qquad A(x)=\av{K_x(\cdot)\log f_\ast(x+\cdot)},\qquad E(x)=\av{K_x\log g_\ast},\\
 u=A+E-D,\qquad \omega=2D-A-E .
\end{gathered}
\end{equation}
These are the functions of~\cite[Section~3]{OAI}: with $f=2^r/\av{2^r}$ and $g=2^s/\av{2^s}$ one has
$K_x(w)=f(x+w)g(w)/\av{f(x+\cdot)g}$ and $f_\ast=f^\kappa/\av{f^\kappa}$, the normalizers of $f$ and $g$ cancelling.
The function called $v$ in~\cite{OAI} and $w$ in~\cite{Base} is our $\omega$. Note the convention: the first density
is evaluated at $x+w$ and the second at $w$. No symmetry of $r$ or $s$ is assumed. On a finite circle every mean,
including those in the normalizers, is the finite mean; we write $D_P,A_P,\dots$ when the circle must be specified.

\begin{lemma}[entropy order and bounds]\label{lem:order}
On either circle, for every $x$: $K_x$, $f_\ast(x+\cdot)$ and $g_\ast$ are strictly positive with mean $1$, and
\begin{equation}\label{eq:order}
 D(x)\ge0,\qquad D(x)-A(x)\ge0,\qquad D(x)-E(x)\ge0 .
\end{equation}
Hence $\omega\ge0$ and $u+\omega=D\ge0$. Moreover, with $C_0=2\kappa^+R$,
\begin{equation}\label{eq:C0}
\begin{gathered}
 \abs{\log K_x(w)},\ \abs{\log f_\ast},\ \abs{\log g_\ast},\ \abs D,\ \abs A,\ \abs E\ \le\ C_0,\\
 \abs u\le U_\infty:=2(1+\kappa)R,\qquad 0\le\omega\le W_\infty:=(4+2\kappa)R .
\end{gathered}
\end{equation}
For the filter of Section~\ref{sec:params}, $R\le7.24$ and $C_0\le15$.
\end{lemma}
\begin{proof}
Positivity and the normalizations are clear from~\eqref{eq:fg}; for $f_\ast(x+\cdot)$ use shift invariance of the
mean, which holds on both circles. For positive functions $\pi_1,\pi_2$ with mean $1$, concavity of the logarithm
gives $\av{\pi_1\ln(\pi_2/\pi_1)}\le\ln\av{\pi_2}=0$ (Jensen's inequality for the probability measure with density
$\pi_1$). With $\pi_1=K_x$ and $\pi_2=1$, $f_\ast(x+\cdot)$, $g_\ast$ in turn this gives the three inequalities
of~\eqref{eq:order}. For the bounds: $2^{-R}\le k\le2^{R}$, so $2^{-R}\le q\le2^R$ and $\abs{\log K_x}\le2R$;
$2^{-\kappa R_r}\le2^{\kappa r},\av{2^{\kappa r}}\le2^{\kappa R_r}$ gives $\abs{\log f_\ast}\le2\kappa R_r$, and likewise
$\abs{\log g_\ast}\le2\kappa R_s$. $D$, $A$, $E$ are averages of these logarithms against the probability density
$K_x$. Then $\abs u\le\abs A+\abs E+\abs D\le2\kappa R+2R$ and $\omega\le2\abs D+\abs A+\abs E$.
\end{proof}

It is convenient to write the functions through three means. Put
$Y_r(x)=\av{k(x,\cdot)\,r(x+\cdot)}$, $Y_s(x)=\av{k(x,\cdot)\,s}$, $N_r=\av{2^{\kappa r}}$, $N_s=\av{2^{\kappa s}}$. Then
\begin{equation}\label{eq:means}
 D=\frac{Y_r+Y_s}{q}-\log q,\qquad A=\frac{\kappa Y_r}{q}-\log N_r,\qquad E=\frac{\kappa Y_s}{q}-\log N_s ,
\end{equation}
and therefore
\begin{equation}\label{eq:uomega}
 u=(\kappa-1)\frac{Y_r+Y_s}{q}+\log q-\log(N_rN_s),\qquad
 \omega=(2-\kappa)\frac{Y_r+Y_s}{q}-2\log q+\log(N_rN_s).
\end{equation}
These are the formulas that the certificate evaluates (Section~\ref{sec:cert}).

\subsection{Convergence of the finite circle}
\begin{lemma}[Riemann sums]\label{lem:riemann}
Let $H$ be a continuously differentiable $2\pi$-periodic function and $N\ge1$. Then
$\bigl\lvert N^{-1}\sum_{i=0}^{N-1}H(2\pi i/N)-\frac1{2\pi}\int_0^{2\pi}H\bigr\rvert\le\pi\,\|H'\|_\infty/N$.
\end{lemma}
\begin{proof}
On $[\theta_i,\theta_{i+1}]$ with $\theta_i=2\pi i/N$, $\abs{H(\theta)-H(\theta_i)}\le\|H'\|_\infty(\theta-\theta_i)$, so
$\bigl\lvert\frac1{2\pi}\int_{\theta_i}^{\theta_{i+1}}(H-H(\theta_i))\bigr\rvert\le\frac{\|H'\|_\infty}{2\pi}\cdot\frac{(2\pi/N)^2}{2}$.
Sum over $i$.
\end{proof}

\begin{lemma}[finite-circle functions]\label{lem:conv}
Let $C_1=64\,(\kappa^+)^2\,2^{2\kappa^+R}\,R^1\,(1+R)^2$. For every integer $N\ge1$ and every real $x$, let
$D^{(N)}(x),A^{(N)}(x),E^{(N)}(x)$ be defined by~\eqref{eq:means} with every mean over $w$ (and the means $N_r,N_s$)
replaced by the $N$-point average over $w\in\{2\pi i/N\}$. Then for $\Xi\in\{D,A,E,u,\omega\}$,
\begin{equation}\label{eq:conv}
 \bigl\lvert\Xi^{(N)}(x)-\Xi(x)\bigr\rvert\le\frac{C_1}{N}\quad\text{for all real }x,\qquad
 \abs{\log f^{(N)}_\ast-\log f_\ast},\ \abs{\log g^{(N)}_\ast-\log g_\ast}\le\frac{C_1}{N},
\end{equation}
where the functions without superscript are those of the continuous circle. In particular, for $P\ge2$ and
$x\in\Z/P\Z$, $\abs{\Xi_P(x)-\Xi(2\pi x/P)}\le C_1/P$. Moreover $u$ and $\omega$ are continuously differentiable on
$\Tc$ with $\abs{u'},\abs{\omega'}\le C_1$.
\end{lemma}
\begin{proof}
For fixed $x$ the $w$-derivatives of $k(x,w)$, $k(x,w)r(x+w)$ and $k(x,w)s(w)$ are bounded in absolute value by
$2^RR^1(1+R)$ (use $\partial_wk=\ln2\,(r'(x+w)+s'(w))\,k$ and $\ln2<1$). By Lemma~\ref{lem:riemann}, each of
$q^{(N)},Y^{(N)}_r,Y^{(N)}_s$ differs from $q,Y_r,Y_s$ by at most $e_1:=\pi2^RR^1(1+R)/N$. Both $q$ and $q^{(N)}$ lie
in $[2^{-R},2^R]$ and $\abs{Y_r/q}\le R_r$, $\abs{Y_s/q}\le R_s$ (weighted averages). Hence
\[
 \Bigl\lvert\frac{Y^{(N)}}{q^{(N)}}-\frac Yq\Bigr\rvert\le\frac{\abs{Y^{(N)}-Y}}{q^{(N)}}+\frac{\abs Y}{q}\cdot\frac{\abs{q-q^{(N)}}}{q^{(N)}}
 \le2^R(1+R)\,e_1
\]
for $Y=Y_r$ or $Y_s$, and $\abs{\log q^{(N)}-\log q}\le\abs{q^{(N)}-q}/(\ln2\cdot2^{-R})\le2^{R+1}e_1$. Similarly
$\abs{(2^{\kappa r})'}\le\kappa R^1_r2^{\kappa R_r}$ gives $\abs{N^{(N)}_r-N_r}\le\pi\kappa R^1_r2^{\kappa R_r}/N$ and
$\abs{\log N^{(N)}_r-\log N_r}\le5\kappa R^12^{2\kappa R}/N$; the same for $N_s$. Inserting these in~\eqref{eq:means}
and~\eqref{eq:uomega}: the error of $D$ is at most $2\cdot2^R(1+R)e_1+2^{R+1}e_1\le4\pi2^{2R}R^1(1+R)^2/N$, the errors
of $A$ and $E$ are at most $9\kappa^+2^{2\kappa^+R}R^1(1+R)^2/N$, and the errors of $u$ and $\omega$ are at most the
sum of two errors of $D$ and one each of $A$ and $E$, which is below $C_1/N$. The statement for $\log f_\ast$
follows from that for $\log N_r$. For $x$ on the grid of $\Z/P\Z$ and $N=P$, $\Xi^{(P)}(2\pi x/P)=\Xi_P(x)$.

For the derivatives in $x$ on $\Tc$: $\abs{\partial_xk}\le R^1_r2^R$, so $\abs{q'}\le R^12^R$ and
$\abs{Y_r'},\abs{Y_s'}\le2^RR^1(1+R)$; then $\abs{(Y/q)'}\le2\cdot2^{2R}R^1(1+R)$ and $\abs{(\log q)'}\le2\cdot2^{2R}R^1$,
and~\eqref{eq:uomega} gives $\abs{u'},\abs{\omega'}\le C_1$.
\end{proof}

\begin{lemma}[finite-circle moments]\label{lem:B1}
Let $\Lambda(x)=2+2\cos x$ (on $\Z/P\Z$: $\Lambda(x)=2+2\cos(2\pi x/P)$). Assume \textup{(M1)}, \textup{(M2)},
\textup{(N11)}, \textup{(N14)}. There is an absolute integer $P_0$ such that for every integer $P\ge P_0$,
\begin{equation}\label{eq:finmom}
 \frac1P\sum_{x\in\Z/P\Z}\Lambda(x)\,2^{p_1(\Delta-u_P(x))}<2^{\rho_f},\qquad
 \frac1P\sum_{x\in\Z/P\Z}\Lambda(x)\,2^{p_2(a_\Gamma\omega_P(x)-u_P(x))}<2^{\rho_f}.
\end{equation}
One may take
$P_0=\bigl\lceil 20\,(1+p_1+2p_2)\,C_1\,2^{\beta_0}/(2^{\rho_f}-2^{\rho_c})\bigr\rceil$ with
$\beta_0=\max\bigl(p_1(\Delta+U_\infty),\,p_2(a_\Gamma W_\infty+U_\infty)\bigr)$.
\end{lemma}
\begin{proof}
Consider the first moment; write $\Psi(x)=\Lambda(x)2^{p_1(\Delta-u(x))}$ on $\Tc$ and $\Psi_P$ for its finite-circle
analogue. For real $\alpha,\beta\le\beta_0$, $\abs{2^\alpha-2^\beta}\le2^{\beta_0}\abs{\alpha-\beta}$. By
Lemma~\ref{lem:conv} and $0\le\Lambda\le4$, $\abs{\Psi_P(x)-\Psi(2\pi x/P)}\le4\cdot2^{\beta_0}p_1C_1/P$ at every grid
point. Also $\abs{\Psi'}\le2\cdot2^{\beta_0}+4\cdot2^{\beta_0}p_1C_1$, so by Lemma~\ref{lem:riemann} the $P$-point
average of $\Psi$ differs from $\int_\Tc\Psi$ by at most $\pi2^{\beta_0}(2+4p_1C_1)/P$. The two errors add up to
$2^{\beta_0}\bigl(2\pi+(4+4\pi)p_1C_1\bigr)/P$, and $2\pi+(4+4\pi)p_1C_1\le20(1+p_1)C_1$ because $C_1\ge1$ (for our
filter $C_1>3\cdot10^9$, Section~\ref{sec:cert}). Together,
\[
 \Bigl\lvert\frac1P\sum_x\Psi_P(x)-\int_\Tc\Psi\Bigr\rvert\le\frac{20(1+p_1)C_1\,2^{\beta_0}}{P}\le2^{\rho_f}-2^{\rho_c}
\]
for $P\ge P_0$, and $\int_\Tc\Psi<2^{\rho_c}$ by \textup{(M1)}. The second moment is the same, with exponent
$p_2(a_\Gamma\omega-u)$, which is Lipschitz in $(u,\omega)$ with constant $p_2(1+a_\Gamma)\le2p_2$ by
\textup{(N11)}; so the same bound holds with $2p_2$ in place of $p_1$, and both bounds are at most
$20(1+p_1+2p_2)C_1\,2^{\beta_0}/P$.
\end{proof}
Strict positivity of the densities and smoothness enter only through the bounds $2^{-R}\le k\le2^R$ and through
$R^1$; sine terms change neither. $P_0$ is a (large) absolute constant: with our filter $\beta_0<270$ and $C_1<2^{32}$, so
$P\ge P_0$ holds once $\eta n\ge340$.

\subsection{Effective evaluation}
\begin{lemma}[effective evaluation]\label{lem:eval}
For every rational $\varrho>0$ there are an integer $P_\varrho$ and a deterministic procedure which, given an integer
$P\ge P_\varrho$ and a residue $x\in\Z/P\Z$, returns rational numbers within $\varrho$ of $D_P(x)$, $A_P(x)$,
$E_P(x)$, $\log f_{\ast,P}(x)$ and $\log g_{\ast,P}(x)$, in time polynomial in $\log P$, without enumerating
$\Z/P\Z$. Called twice with the same arguments it returns the same values.
\end{lemma}
\begin{proof}
Let $N=\ceil{8C_1/\varrho}$ and $P_\varrho=N$. By Lemma~\ref{lem:conv}, for $P\ge P_\varrho$ each of the five finite-circle
values at $x$ is within $\varrho/8$ of the continuous value at $\theta=2\pi x/P$, which is within $\varrho/8$ of the
$N$-point value $\Xi^{(N)}(\theta)$. The latter is a fixed finite expression: by~\eqref{eq:means} it is built from
the $2N$ numbers $r(\theta+2\pi i/N)$, $s(2\pi i/N)$ by finitely many additions, multiplications, one division by
$q^{(N)}\in[2^{-R},2^R]$, the function $z\mapsto2^z$ on $\abs z\le\kappa^+R$ and the function $\log$ on
$[2^{-\kappa^+R},2^{\kappa^+R}]$. On these compact domains all the operations are Lipschitz with constants depending
only on $(r,s,\kappa)$, so there is a fixed number of bits $Q_\varrho$ such that evaluating every elementary function
to within $2^{-Q_\varrho}$ gives the result to within $\varrho/4$. The elementary functions are evaluated in rational
arithmetic as in Lemma~\ref{lem:phase}: the angle from $x/P$ and a rational approximation of $\pi$; $\cos$ and
$\sin$ of arguments in $[0,7J+7]$ by Taylor series; $2^z=\exp(z\ln2)$ by the exponential series; and
$\ln y=2\sum_{i\ge0}\frac1{2i+1}\bigl(\frac{y-1}{y+1}\bigr)^{2i+1}$, which converges geometrically, uniformly for $y$
in a compact subinterval of $(0,\infty)$ (it also gives $\ln2$). The numbers of terms and the working precision are
constants depending on $\varrho$; the only dependence on $P$ is the division $x/P$ to a fixed number of bits. The
total error is below $\varrho/8+\varrho/8+\varrho/4<\varrho$. The procedure is deterministic.
\end{proof}
Since the approximations need not preserve the order relations~\eqref{eq:order} exactly, Section~\ref{sec:checksum}
rounds them in a fixed direction before use.
\section{The checksum estimator and Theorem B}\label{sec:checksum}
In this section $P$ is an integer with $\abs{\log P-\eta n}\le2$ (no primality is needed), $\bc_1,\dots,\bc_n$ are
residues modulo $P$, and for a residue $\rho$
\begin{equation}\label{eq:V}
 V_\rho=\sum_{I\subseteq[n]:\ \bc(I)\equiv\rho\ (\mathrm{mod}\ P)}\chi(I),\qquad \chi(I)=e_P\Bigl(\sum_{i\in I}\varphi_i\Bigr).
\end{equation}
All filter functions are the finite-circle functions on $\Z/P\Z$ of Section~\ref{sec:filter}; we drop the subscript
$P$. We write $\tau$ for $\tau_0$, and assume $n$ so large that $P\ge\max(P_0,P_{\eps/8})$
(Lemmas~\ref{lem:B1},~\ref{lem:eval}).

\begin{thmB}[modular checksum with one filter]
Fix a constant $d>0$. There are a randomized procedure and an integer $n_B$ with the following properties for
$n\ge n_B$. Its inputs are $P$, the residues $\bc_i$, phases $\varphi_1,\dots,\varphi_n\in\Z/P\Z$, and a list $\cT$ of at
most $n^d$ target residues. It returns \textsc{fail} or numbers $\hat V_\rho$ ($\rho\in\cT$) with dyadic real and
imaginary parts of $n^2+O(n)$ bits.
\begin{enumerate}[label=\textup{(\arabic*)}]
\item Let $P$, the residues $\bc_i$ and the list $\cT$ be fixed, that is, chosen without reference to the phases. If
the phases are independent and uniform on $\Z/P\Z$ and independent of the procedure's own draws, then with
probability at least $1-\epsilon_B(n)$, over the phases and those draws, the procedure does not fail and
$\abs{\hat V_\rho-V_\rho}\le\tfrac18$ for every $\rho\in\cT$. Here $\epsilon_B(n)\to0$ depends only on $n$ and $d$; the
bound is uniform in $P$, in the residues and in $\cT$.
\item On every execution the procedure uses at most $2^{(\tau+20\eps)n+c_B\sqrt n}\,\poly(n+\abs\cT)$ word operations,
$c_B$ an absolute constant.
\end{enumerate}
\end{thmB}

\subsection{Fourier form}
\begin{lemma}\label{lem:fourier}
For every residue $\rho$,
\begin{equation}\label{eq:fourier}
 V_\rho=\frac1P\sum_{k=0}^{P-1}J_\rho(k),\qquad J_\rho(k)=e_P(-k\rho)\prod_{i=1}^n\bigl(1+e_P(x_i(k))\bigr),\qquad
 x_i(k)=\varphi_i+k\bc_i\bmod P,
\end{equation}
and $\abs{J_\rho(k)}^2=\prod_{i=1}^n\Lambda(x_i(k))$. If the $\varphi_i$ are independent and uniform on $\Z/P\Z$ then, for
each fixed $k$, the residues $x_1(k),\dots,x_n(k)$ are independent and uniform, whatever the $\bc_i$ are.
\end{lemma}
\begin{proof}
Expanding the product, $J_\rho(k)=\sum_{I}\chi(I)\,e_P\bigl(k(\bc(I)-\rho)\bigr)$, and
$P^{-1}\sum_{k<P}e_P(kz)$ is $1$ if $P\mid z$ and $0$ otherwise. Further $\abs{1+e^{\mathrm i\theta}}^2=2+2\cos\theta$.
For fixed $k$ the map $\varphi\mapsto(\varphi_i+k\bc_i)_i$ is a translation of $(\Z/P\Z)^n$, which preserves the uniform
law.
\end{proof}
For two different frequencies $k\ne k'$ the vectors $x(k)$ and $x(k')$ are \emph{not} independent
($x_i(k)-x_i(k')=(k-k')\bc_i$ is deterministic). The argument below uses only the law of $x(k)$ for one $k$ at a
time, combined by linearity of expectation.

\subsection{Patterns, digits and tables}\label{sec:tables}
\paragraph{Blocks.} Let $N_b=\ceil{\sqrt n}$ and cut $[n]$ into consecutive blocks $\mathsf B_1,\dots,\mathsf B_{N_b}$
of sizes $n_\beta\in\{\floor{n/N_b},\ceil{n/N_b}\}$; so $\sqrt n/2\le n_\beta\le\sqrt n+1$ for $n\ge16$ and
$\sum_\beta n_\beta=n$.

\paragraph{Fixed approximations.} By Lemma~\ref{lem:eval} fix deterministic rational approximations of $D(x)$,
$A(x)$, $E(x)$ within $\eps/8$, and approximations $\tilde\ell_f(y)$, $\tilde\ell_g(y)$ of $\log f_\ast(y)$,
$\log g_\ast(y)$ within $\eps$. The same approximations are used at every occurrence.

\paragraph{Patterns.} Let $\mathbb G$ be the set of integer multiples of $\eps/4$ in $[-C_0-\eps,C_0+\eps]$; it has at
most $8C_0/\eps+9$ elements. For a frequency $k<P$ and a block $\beta$ let $\bar D_\beta,\bar A_\beta,\bar E_\beta$ be the
averages of $D,A,E$ over the coordinates $x_i(k)$, $i\in\mathsf B_\beta$, and let
$\tilde D_\beta,\tilde A_\beta,\tilde E_\beta$ be the averages of the fixed approximations (so each is within $\eps/8$ of
the true average). Define $D_\beta$ as the least element of $\mathbb G$ that is at least $\tilde D_\beta+\eps/8$, and
$A_\beta$, $E_\beta$ as the greatest elements of $\mathbb G$ that are at most $\tilde A_\beta-\eps/8$,
$\tilde E_\beta-\eps/8$. The \emph{pattern} $\pat(k)$ of $k$ is the tuple $(D_\beta,A_\beta,E_\beta)_{\beta\le N_b}$. It is a
deterministic function of $k$, the phases and the residues, computable in $\poly(n)$ time.

\begin{lemma}[pattern rule]\label{lem:pattern}
For every $k$ and $\beta$: $0\le D_\beta-\bar D_\beta<\eps/2$, $0\le\bar A_\beta-A_\beta<\eps/2$,
$0\le\bar E_\beta-E_\beta<\eps/2$; hence
\begin{equation}\label{eq:patorder}
 D_\beta\ge0,\qquad A_\beta\le D_\beta,\qquad E_\beta\le D_\beta .
\end{equation}
\end{lemma}
\begin{proof}
$\tilde D_\beta+\eps/8\ge\bar D_\beta$, and rounding up to the grid adds less than $\eps/4$, while
$\tilde D_\beta+\eps/8\le\bar D_\beta+\eps/4$; so $\bar D_\beta\le D_\beta<\bar D_\beta+\eps/2$. The same for $A,E$ with
directions reversed. All values stay in $[-C_0-\eps,C_0+\eps]$ by~\eqref{eq:C0}. By~\eqref{eq:order},
$\bar D_\beta\ge0$, $\bar A_\beta\le\bar D_\beta$ and $\bar E_\beta\le\bar D_\beta$, which gives~\eqref{eq:patorder}.
\end{proof}
Let $\cP$ be the set of \emph{all} tuples $(D_\beta,A_\beta,E_\beta)_\beta\in\mathbb G^{3N_b}$
satisfying~\eqref{eq:patorder}, whether or not some frequency attains them. Then
\begin{equation}\label{eq:Qn}
 Q_n:=\abs\cP\le(8C_0/\eps+9)^{3N_b}\le2^{c_{\cP}\sqrt n},\qquad c_{\cP}=6\log(8C_0/\eps+9).
\end{equation}
The set $\cP$ does not depend on the phases; only the map $k\mapsto\pat(k)$ does. For $\pi\in\cP$ put
\[
 (D_\pi,A_\pi,E_\pi)=\frac1n\sum_\beta n_\beta\,(D_\beta,A_\beta,E_\beta),\qquad U_\pi=A_\pi+E_\pi-D_\pi,\qquad
 W_\pi=2D_\pi-A_\pi-E_\pi\ \ (\ge0),
\]
\begin{equation}\label{eq:alpha}
 \alpha_\pi=\begin{cases}1,&W_\pi\le\Gamma,\\ \Gamma/W_\pi,&W_\pi>\Gamma,\end{cases}\qquad
 H_\pi=\tau-\alpha_\pi(D_\pi-A_\pi).
\end{equation}
All these are rationals with numerators and denominators of $O(\log n)$ bits.

\begin{lemma}[digit balance]\label{lem:balance}
$0<\alpha_\pi\le1$, $\alpha_\pi W_\pi\le\Gamma$, $\Delta\le H_\pi\le\tau$, and
$\eta-H_\pi+\alpha_\pi(D_\pi-E_\pi)=\Delta+\alpha_\pi W_\pi\le\tau$.
\end{lemma}
\begin{proof}
$D_\pi-A_\pi\ge0$ and $D_\pi-E_\pi\ge0$ by~\eqref{eq:patorder}, and their sum is $W_\pi$. So $H_\pi\le\tau$ and
$H_\pi\ge\tau-\alpha_\pi W_\pi\ge\tau-\Gamma=\Delta$. The identity is $\eta-\tau=\Delta$ and the last inequality is
$\Delta+\Gamma=\tau$.
\end{proof}

\paragraph{Digits.} Let $B_\pi=2^{\ceil{(\eta-H_\pi)n}}$. High digits are the integers $0\le h<\ceil{P/B_\pi}$, low
digits the integers $0\le\ell<B_\pi$; the pair $(h,\ell)$ \emph{represents} $k=hB_\pi+\ell$, and each $k<P$ is
represented by exactly one pair. The digit vectors are
\[
 y_i(h)=\varphi_i+hB_\pi\bc_i,\qquad z_i(\ell)=-\ell\,\bc_i\pmod P,\qquad\text{so}\quad y_i(h)-z_i(\ell)=x_i(k).
\]
Since $\abs{\log P-\eta n}\le2$, there are at most $2^{H_\pi n+3}\le2^{\tau n+3}$ high digits, at most
$2^{(\eta-H_\pi)n+1}\le2^{\tau n+1}$ low digits, and at most $B_\pi\ceil{P/B_\pi}<P+B_\pi\le2^{\eta n+2}+2^{\tau n+1}\le2^{\eta n+3}$ pairs (some of which represent
$k\ge P$).

\paragraph{Tables.} For each $\pi\in\cP$ and each block $\beta$ draw
\begin{equation}\label{eq:Rb}
 R_\beta=2^{\ceil{(\alpha_\pi D_\beta+4\eps)n_\beta}}
\end{equation}
independent \emph{block tables}. A block table assigns to each coordinate $i\in\mathsf B_\beta$, independently, the
label $\ast$ with probability $1-\alpha_\pi$ and otherwise a uniform label $\theta_i\in\Z/P\Z$. All block tables (for
all patterns and blocks) are independent of each other and of the phases, and they are stored.

\paragraph{Tests.} At a coordinate labelled $\ast$ both scores are $0$. At a coordinate labelled $\theta$, the exact
high score of the digit $h$ is $\log f_\ast(y_i(h)-\theta)$ and the exact low score of $\ell$ is
$\log g_\ast(z_i(\ell)-\theta)$; the approximate scores are $\tilde\ell_f(y_i(h)-\theta)$ and
$\tilde\ell_g(z_i(\ell)-\theta)$. A high digit \emph{passes} a block table if the average of its approximate scores
over the block is at least $\alpha_\pi A_\beta-3\eps$; a low digit passes if the average is at least
$\alpha_\pi E_\beta-3\eps$. A \emph{whole table} is a tuple of one block-table index per block; a digit passes it if it
passes each of its block tables. An \emph{entry} is a digit together with a whole table that it passes. A
\emph{match} is a high entry and a low entry with the same whole table. $L_\pi$ is the number of matches, and
$\mu_\pi(k)$ the number of matches whose digit pair represents $k$. These definitions do not require the represented
frequency to have pattern $\pi$.

\subsection{Coverage and sizes}
\begin{lemma}[coverage and sizes]\label{lem:cov}
Let $n$ be so large that $n_\beta\ge6C_0^2/\eps^2$ and $n_\beta\ge(\log n+1)/(3\eps)$ for all $\beta$.
\begin{enumerate}[label=\textup{(\roman*)}]
\item For arbitrary fixed phases and residues, the probability over the tables that some $k<P$ has
$\mu_{\pat(k)}(k)=0$ is at most $2^{\eta n+2}(\sqrt n+1)\exp(-2^{3\eps n_{\min}-1})$, where
$n_{\min}=\min_\beta n_\beta$.
\item For each $\pi\in\cP$: conditionally on arbitrary phases, the expected numbers of high and of low entries are
at most $2^{(\tau+8\eps)n+N_b+3}$; and, when the phases are independent and uniform,
\begin{equation}\label{eq:EL}
 \E\,L_\pi\le2^{(\eta-\alpha_\pi U_\pi+12\eps)n+N_b+3},
\end{equation}
the expectation being over the phases and the tables.
\item With independent uniform phases, with probability at least $1-\epsilon_{\cG}(n)$, where
$\epsilon_{\cG}(n)=2^{\eta n+2}(\sqrt n+1)\exp(-2^{3\eps n_{\min}-1})+3Q_n2^{-6\eps n+N_b+3}$, the following event $\cG$
holds: every $k<P$ has $\mu_{\pat(k)}(k)\ge1$, and for every $\pi\in\cP$ both entry counts are at most
$2^{(\tau+14\eps)n}$ and
\begin{equation}\label{eq:Lgood}
 L_\pi\le2^{(\eta-\alpha_\pi U_\pi+18\eps)n}.
\end{equation}
\end{enumerate}
\end{lemma}
\begin{proof}
(i) Fix the phases, a frequency $k<P$ with digits $(h,\ell)$, its pattern $\pi=\pat(k)$ with $\alpha=\alpha_\pi$, and a
block $\beta$. Write $x_i=x_i(k)$, $y_i=y_i(h)$, $z_i=z_i(\ell)$. Let $\Pr$ be the law of one block table and let
$\widetilde\Pr$ be the law under which the labels are still independent, $\ast$ still has probability $1-\alpha$,
but a non-$\ast$ label $\theta_i$ is such that $w_i=z_i-\theta_i$ has density $K_{x_i}$ with respect to the uniform
law on $\Z/P\Z$. Since $K>0$ the two laws are mutually absolutely continuous, with
$\log\frac{d\widetilde\Pr}{d\Pr}=\sum_{i\in\mathsf B_\beta}\Xi^{\mathrm{lr}}_i$, where
$\Xi^{\mathrm{lr}}_i=\log K_{x_i}(w_i)$ at non-$\ast$ coordinates and $0$ at $\ast$. At a non-$\ast$ coordinate the
exact high score is $\Xi^{\mathrm{hi}}_i=\log f_\ast(y_i-\theta_i)=\log f_\ast(x_i+w_i)$ and the exact low score is
$\Xi^{\mathrm{lo}}_i=\log g_\ast(w_i)$. By~\eqref{eq:DAE}, under $\widetilde\Pr$ the three variables have means
$\alpha A(x_i)$, $\alpha E(x_i)$, $\alpha D(x_i)$; they are bounded by $C_0$ in absolute value~\eqref{eq:C0} and
independent across $i$. By Chebyshev's inequality and the union bound, with $\widetilde\Pr$-probability at least
$1-3C_0^2/(n_\beta\eps^2)\ge\tfrac12$ the three block averages are within $\eps$ of their means
$\alpha\bar A_\beta,\alpha\bar E_\beta,\alpha\bar D_\beta$. Call this event $\mathcal Q$. On $\mathcal Q$, using
Lemma~\ref{lem:pattern} and $\alpha\le1$: the approximate high average is at least
$\alpha\bar A_\beta-\eps-\eps\ge\alpha A_\beta-3\eps$, the approximate low average is at least $\alpha E_\beta-3\eps$, so
the block table accepts both digits; and $\log\frac{d\widetilde\Pr}{d\Pr}\le(\alpha\bar D_\beta+\eps)n_\beta\le(\alpha D_\beta+\eps)n_\beta$.
Therefore
\[
 \Pr[\text{the block table accepts both digits}]\ \ge\ \Pr[\mathcal Q]=\widetilde\E\Bigl[\mathbf 1_{\mathcal Q}\frac{d\Pr}{d\widetilde\Pr}\Bigr]
 \ \ge\ \tfrac12\,2^{-(\alpha D_\beta+\eps)n_\beta}.
\]
The $R_\beta\ge2^{(\alpha D_\beta+4\eps)n_\beta}$ block tables are independent, so none accepts both digits with
probability at most $\exp(-\tfrac12R_\beta2^{-(\alpha D_\beta+\eps)n_\beta})\le\exp(-2^{3\eps n_\beta-1})$. If every block has
an accepting table, their tuple is a whole table passed by both digits, that is, a match of pattern $\pi$
representing $k$. A union bound over the $N_b\le\sqrt n+1$ blocks and the $P\le2^{\eta n+2}$ frequencies gives (i).

(ii) By~\eqref{eq:Rb} the number of whole tables is at most
$\prod_\beta R_\beta\le2^{(\alpha D_\pi+4\eps)n+N_b}$. Fix a high digit $h$ and a whole-table index tuple, and let
$\mathsf S$ be the sum over all $n$ coordinates of the exact high scores. The labels are independent and, by the
normalization of $f_\ast$ on $\Z/P\Z$,
$\E_{\text{tables}}2^{\mathsf S}=\prod_i\bigl((1-\alpha)+\alpha\av{f_\ast}\bigr)=1$ for every value of the phases. If
$h$ passes the whole table, each block average of approximate scores is at least $\alpha A_\beta-3\eps$, so each block
sum of exact scores is at least $(\alpha A_\beta-4\eps)n_\beta$ and $\mathsf S\ge(\alpha A_\pi-4\eps)n$. By Markov's
inequality this has probability at most $2^{-(\alpha A_\pi-4\eps)n}$. Multiplying by the numbers of high digits and of
whole tables (linearity of expectation; no independence between whole tables is used),
\[
 \E[\#\text{high entries}\mid\varphi]\le2^{H_\pi n+3}\cdot2^{(\alpha D_\pi+4\eps)n+N_b}\cdot2^{-(\alpha A_\pi-4\eps)n}=2^{(\tau+8\eps)n+N_b+3}.
\]
For low digits the same argument with $g_\ast$ gives
$2^{(\eta-H_\pi)n+1}2^{(\alpha D_\pi+4\eps)n+N_b}2^{-(\alpha E_\pi-4\eps)n}\le2^{(\tau+8\eps)n+N_b+1}$ by
Lemma~\ref{lem:balance}.

For matches fix a digit pair $(h,\ell)$ (representing some $k$, possibly $k\ge P$) and a whole-table index tuple, and
let $\mathsf S_h$, $\mathsf S_\ell$ be the two sums of exact scores. Now take the expectation over the tables
\emph{and} over independent uniform phases. The pairs $(\varphi_i,\text{label}_i)$ are independent across $i$. At a
non-$\ast$ coordinate, conditionally on $\theta_i$ the residue $y_i(h)-\theta_i=\varphi_i+hB_\pi\bc_i-\theta_i$ is
uniform, so the factor $f_\ast(y_i-\theta_i)$ averages to $1$; then $g_\ast(z_i(\ell)-\theta_i)$ averages to $1$ over
the uniform $\theta_i$. Hence
\begin{equation}\label{eq:joint}
 \E_{\varphi,\text{tables}}\,2^{\mathsf S_h+\mathsf S_\ell}=1 .
\end{equation}
A match requires $\mathsf S_h+\mathsf S_\ell\ge(\alpha(A_\pi+E_\pi)-8\eps)n$, which by Markov has probability at
most $2^{-(\alpha(A_\pi+E_\pi)-8\eps)n}$. Multiply by the at most $2^{\eta n+3}$ digit pairs and the number of whole
tables to get~\eqref{eq:EL}. It is essential that \emph{all} digit pairs are counted, whatever frequency they
represent and whatever its pattern, and that the phases have their original product law: \eqref{eq:joint} is false
conditionally on the phases.

(iii) Each of the three bounds defining $\cG$ exceeds the corresponding expectation in (ii) by a factor
$2^{6\eps n-N_b-3}$; apply Markov's inequality and a union bound over $\pi\in\cP$, and add (i).
\end{proof}

\subsection{Lists, multiplicities and sampling}\label{sec:lists}
For one pattern $\pi$ the lists are built block by block: for each digit, test its vector against the $R_\beta$
tables of each block and record the set of passing indices; if some set is empty the digit has no entry; otherwise
the entries of the digit are the elements of the Cartesian product of these sets, emitted one at a time by a nested
counter. By~\eqref{eq:C0} and~\eqref{eq:Rb},
\begin{equation}\label{eq:sumR}
 \sum_\beta R_\beta\le N_b\,2^{(C_0+6\eps)(\sqrt n+1)+1}\le2^{c_R\sqrt n}\qquad(c_R=2C_0+2),
\end{equation}
so the work per digit, apart from emitted entries, is $2^{c_R\sqrt n}\poly(n)$. Sort the high and low entry lists by
whole-table index tuple. If the group of the tuple $\nu$ has $a_\nu$ high and $b_\nu$ low entries, then
$L_\pi=\sum_\nu a_\nu b_\nu$; prefix sums of the products $a_\nu b_\nu$ allow a uniformly random match to be drawn
without enumerating the matches (draw a uniform integer below $L_\pi$, locate its group by binary search, and split
the offset by quotient and remainder), and the same structure enumerates the matches sequentially.

\begin{lemma}[multiplicity]\label{lem:mult}
For $k=hB_\pi+\ell$,
\[
 \mu_\pi(k)=\prod_{\beta}\#\{\text{block tables of block }\beta\text{ that both }h\text{ and }\ell\text{ pass}\}.
\]
It is computed by retesting the stored block tables in $2^{c_R\sqrt n}\poly(n)$ time. Every match that is sampled or
enumerated represents a frequency with $\mu_\pi\ge1$.
\end{lemma}
\begin{proof}
The digits of $k$ are unique, and a match representing $k$ is exactly a whole table passed by both digits, that is,
a choice in each block of a block table passed by both. The tests are deterministic functions of the stored labels
and of the fixed approximations, so the recomputation reproduces the list construction exactly.
\end{proof}

\paragraph{The estimator.} For each $\pi\in\cP$ and each $\rho\in\cT$ define the exact pattern contribution
$v_{\pi,\rho}=P^{-1}\sum_{k<P,\ \pat(k)=\pi}J_\rho(k)$; by~\eqref{eq:fourier}, $V_\rho=\sum_{\pi\in\cP}v_{\pi,\rho}$.
\begin{itemize}
\item \emph{Direct patterns} ($\alpha_\pi U_\pi\ge\Delta$): enumerate all matches; a match representing $k$
contributes $J_\rho(k)/(P\mu_\pi(k))$ if $k<P$ and $\pat(k)=\pi$, and $0$ otherwise.
\item \emph{Sampled patterns} ($\alpha_\pi U_\pi<\Delta$): if $L_\pi=0$ the estimate is $0$. Otherwise draw
$N=2^{\ceil{\tau n}}$ independent uniform matches and average the values
\begin{equation}\label{eq:X}
 X=\frac{L_\pi}{P\,\mu_\pi(k)}\,J_\rho(k)\ \text{ if }k<P\text{ and }\pat(k)=\pi,\qquad X=0\ \text{ otherwise},
\end{equation}
where $k$ is the frequency represented by the sampled match. (One set of $N$ matches serves all $\rho\in\cT$.)
\end{itemize}
The \emph{ideal estimate} $V^\circ_\rho$ is the sum over $\pi\in\cP$ of these pattern estimates, in exact arithmetic
and with exact uniform draws.

\begin{lemma}[conditional analysis]\label{lem:cond}
Fix phases and tables in $\cG$.
\begin{enumerate}[label=\textup{(\alph*)}]
\item A direct pattern's estimate equals $v_{\pi,\rho}$, and it enumerates $L_\pi\le2^{(\tau+18\eps)n}$ matches.
\item A sampled pattern's estimate has conditional mean $v_{\pi,\rho}$, and its conditional mean squared error is at
most
\begin{equation}\label{eq:mse}
 \frac1N\cdot\frac{L_\pi}{P^2}\sum_{k<P:\ \pat(k)=\pi}\abs{J_\rho(k)}^2 .
\end{equation}
If $L_\pi=0$ the estimate $0$ equals $v_{\pi,\rho}$.
\end{enumerate}
\end{lemma}
\begin{proof}
(a) On $\cG$ every $k<P$ with $\pat(k)=\pi$ is represented by $\mu_\pi(k)\ge1$ matches, whose contributions add up to
$J_\rho(k)/P$; all other matches contribute $0$. By~\eqref{eq:Lgood},
$L_\pi\le2^{(\eta-\Delta+18\eps)n}=2^{(\tau+18\eps)n}$. (b) A uniform match represents a given $k$ with probability
$\mu_\pi(k)/L_\pi$. So $\E X=\sum_{k<P,\,\pat(k)=\pi}\frac{\mu_\pi(k)}{L_\pi}\cdot\frac{L_\pi}{P\mu_\pi(k)}J_\rho(k)=v_{\pi,\rho}$
(on $\cG$ all these $k$ have $\mu_\pi(k)\ge1$), and
\[
 \E\abs X^2=\sum_{k<P,\,\pat(k)=\pi}\frac{\mu_\pi(k)}{L_\pi}\cdot\frac{L_\pi^2}{P^2\mu_\pi(k)^2}\abs{J_\rho(k)}^2
 \le\frac{L_\pi}{P^2}\sum_{k<P,\,\pat(k)=\pi}\abs{J_\rho(k)}^2 .
\]
The mean of $N$ independent copies has mean squared error $\Var(X)/N\le\E\abs X^2/N$ (independent centred complex
variables have vanishing cross terms). If $L_\pi=0$ then, on $\cG$, no frequency has pattern $\pi$.
\end{proof}

\subsection{Domination and the error of a sampled pattern}
For a frequency $k$ let $\bar u(k)$ and $\bar\omega(k)$ be the averages of $u$ and $\omega$ over the $n$ coordinates
$x_i(k)$. If $\pat(k)=\pi$ then, by Lemma~\ref{lem:pattern} (three, respectively four, terms each within $\eps$),
\begin{equation}\label{eq:paterr}
 \abs{U_\pi-\bar u(k)}\le3\eps,\qquad\abs{W_\pi-\bar\omega(k)}\le4\eps .
\end{equation}

\begin{lemma}[domination]\label{lem:B2}
Let $\pi$ be a sampled pattern and $k$ a frequency with $\pat(k)=\pi$. Put $\mathsf R_1(x)=p_1(\Delta-u(x))$ and
$\mathsf R_2(x)=p_2(a_\Gamma\omega(x)-u(x))$, and let $e=1$ if $W_\pi\le\Gamma$ and $e=2$ if $W_\pi>\Gamma$. Then
\begin{equation}\label{eq:dom}
 -\alpha_\pi U_\pi\ \le\ -\Delta+\frac1n\sum_{i=1}^n\mathsf R_e(x_i(k))+c_e\,\eps,\qquad c_1=3p_1,\quad c_2=p_2(4a_\Gamma+3).
\end{equation}
By \textup{(N11)}, $\max(c_1,c_2)\le6\max(p_1,p_2)$.
\end{lemma}
\begin{proof}
\emph{Case $W_\pi\le\Gamma$.} Then $\alpha_\pi=1$ and, the pattern being sampled, $\Delta-U_\pi>0$. Since $p_1\ge1$
\textup{(N12)}, $-U_\pi=-\Delta+(\Delta-U_\pi)\le-\Delta+p_1(\Delta-U_\pi)\le-\Delta+p_1(\Delta-\bar u(k))+3p_1\eps$
by~\eqref{eq:paterr}.

\emph{Case $W_\pi>\Gamma$.} Then $\alpha_\pi=\Gamma/W_\pi\in(0,1)$ and $\alpha_\pi a_\Gamma W_\pi=\Delta$. The pattern
being sampled, $\alpha_\pi U_\pi<\Delta=\alpha_\pi a_\Gamma W_\pi$, so $a_\Gamma W_\pi-U_\pi>0$. Since
$\alpha_\pi\le1\le p_2$,
\[
 -\alpha_\pi U_\pi=-\Delta+\alpha_\pi(a_\Gamma W_\pi-U_\pi)\le-\Delta+p_2(a_\Gamma W_\pi-U_\pi)
 \le-\Delta+p_2(a_\Gamma\bar\omega(k)-\bar u(k))+p_2(4a_\Gamma+3)\eps .
\]
In both cases the middle term is $\frac1n\sum_i\mathsf R_e(x_i(k))$. If $a_\Gamma\le\tfrac34$ then
$4a_\Gamma+3\le6$.
\end{proof}
The positivity used in the two cases comes from the sampling condition, not from~\eqref{eq:order}; the lemma is
false for direct patterns and for frequencies of other patterns, and it is applied to neither.

\begin{lemma}[one sampled pattern]\label{lem:B3}
Let the phases be independent and uniform and independent of the tables and of the sampling draws. For a sampled
pattern $\pi$ and a residue $\rho$ let $\mathsf E_{\pi,\rho}$ be the difference between the ideal sampled estimate
and $v_{\pi,\rho}$. Then, with $\beta_3=\eta-\rho_f-(18+6\max(p_1,p_2))\eps>0$ \textup{(N13)},
\[
 \E\bigl[\mathbf 1_{\cG}\,\abs{\mathsf E_{\pi,\rho}}^2\bigr]\le2^{-\beta_3n+2}.
\]
\end{lemma}
\begin{proof}
By Lemma~\ref{lem:cond}(b), on $\cG$ the conditional expectation of $\abs{\mathsf E_{\pi,\rho}}^2$ given the phases
and tables is at most~\eqref{eq:mse}. The phases must not now be averaged ``conditionally on $\cG$'': under that
conditioning they are no longer independent. Instead we bound the integrand pointwise. For every $k<P$,
by~\eqref{eq:Lgood}, Lemma~\ref{lem:fourier} and Lemma~\ref{lem:B2},
\begin{equation}\label{eq:pointwise}
 \mathbf 1_{\cG}\,\mathbf 1_{\{\pat(k)=\pi\}}\,L_\pi\,\abs{J_\rho(k)}^2
 \ \le\ 2^{(\eta-\Delta+(18+c_e)\eps)n}\prod_{i=1}^n\Lambda(x_i(k))\,2^{\mathsf R_e(x_i(k))},
\end{equation}
where $e\in\{1,2\}$ is determined by $\pi$. The right side is nonnegative and involves neither $\cG$ nor the pattern
of $k$, so the inequality may be integrated against the original law. For fixed $k$ the $x_i(k)$ are independent and
uniform (Lemma~\ref{lem:fourier}), so by Lemma~\ref{lem:B1} the expectation of the product is
$\bigl(P^{-1}\sum_x\Lambda(x)2^{\mathsf R_e(x)}\bigr)^n<2^{\rho_fn}$. Summing over the $P$ frequencies (linearity; no
independence between frequencies is used),
\[
 \E\bigl[\mathbf 1_{\cG}\abs{\mathsf E_{\pi,\rho}}^2\bigr]\le\frac{1}{NP^2}\cdot P\cdot2^{(\eta-\Delta+(18+c_e)\eps+\rho_f)n}
 \le2^{(\rho_f+(18+c_e)\eps-\Delta-\tau)n+2},
\]
using $N\ge2^{\tau n}$ and $P\ge2^{\eta n-2}$. Finally $\Delta+\tau=\eta$ and $c_e\le6\max(p_1,p_2)$.
\end{proof}

\subsection{Proof of Theorem~B}
\paragraph{The procedure.} For each $\pi\in\cP$ in turn: compute $\alpha_\pi,H_\pi,B_\pi$ and the $R_\beta$; draw and
store the block tables; build the two entry lists, returning \textsc{fail} before either list would exceed
\begin{equation}\label{eq:Cn}
 C_n=2^{\ceil{(\tau+20\eps)n}}
\end{equation}
entries; sort them and form the prefix sums. For a direct pattern enumerate the matches, returning \textsc{fail}
before enumerating more than $C_n$ of them. For a sampled pattern draw $N$ matches (or return $0$ if $L_\pi=0$).
For each enumerated or sampled match compute $k$, test $k<P$, compute $\pat(k)$, $\mu_\pi(k)$ and, for each
$\rho\in\cT$, the value $J_\rho(k)$. All draws use the samplers of Lemma~\ref{lem:samplers} (\textsc{fail} on a
sampler failure), and the complex arithmetic is that of Lemmas~\ref{lem:phase}--\ref{lem:sums}.

\paragraph{Ideal accuracy.} Let the phases be independent and uniform. For a direct pattern
$\mathsf E_{\pi,\rho}=0$ on $\cG$ (Lemma~\ref{lem:cond}(a)), as it is for a sampled pattern with $L_\pi=0$. By the
Cauchy--Schwarz inequality over the $Q_n$ patterns and Lemma~\ref{lem:B3},
\[
 \E\Bigl[\mathbf 1_{\cG}\Bigl\lvert\sum_{\pi\in\cP}\mathsf E_{\pi,\rho}\Bigr\rvert^2\Bigr]
 \le Q_n\sum_{\pi}\E\bigl[\mathbf 1_{\cG}\abs{\mathsf E_{\pi,\rho}}^2\bigr]\le Q_n^2\,2^{-\beta_3n+2}.
\]
By Markov's inequality and a union bound over $\cT$, the probability that $\cG$ holds and some
$\abs{V^\circ_\rho-V_\rho}$ exceeds $\tfrac1{16}$ is at most $256\,n^d\,Q_n^2\,2^{-\beta_3n+2}$. No independence between
pattern errors or between targets is used. With~\eqref{eq:Qn} this tends to $0$, and so does $\epsilon_{\cG}(n)$.

\paragraph{Caps and samplers.} On $\cG$ no cap is reached: the entry counts are at most
$2^{(\tau+14\eps)n}\le C_n$ and a direct pattern has $L_\pi\le2^{(\tau+18\eps)n}\le C_n$. The procedure makes at most
$2^{2n}$ sampler calls on any execution (see the count below), so by~\eqref{eq:couple} the real run coincides with
the ideal one except with probability $o(1)$.

\paragraph{Arithmetic.} On every execution that reaches the summation, the number of terms is at most
$Q_n\max(N,C_n)\le2^{n}$; each term is a product of at most $n+1$ factors of modulus at most $2$, evaluated to within
$2^{-2n^2+2n}$ (Lemma~\ref{lem:sums}(a)), times an exact rational weight $1/(P\mu_\pi(k))$ or $L_\pi/(NP\mu_\pi(k))$ of
absolute value at most $L_\pi\le C_n^2\le2^{n}$ whose numerator and denominator have $O(n)$ bits (on every execution
that passes the entry caps, $L_\pi\le C_n^2$; and $\mu_\pi(k)\ge1$). By Lemma~\ref{lem:sums}(b) with $K=1$ the total
arithmetic error is at most $2^{-n^2+5n}\le\tfrac1{16}$ for $n\ge8$. Hence
$\abs{\hat V_\rho-V_\rho}\le\tfrac1{16}+\tfrac1{16}$ outside the exceptional events, which proves part~(1) with
\[
 \epsilon_B(n)=\epsilon_{\cG}(n)+2^{10+d\log n+2c_{\cP}\sqrt n-\beta_3n}+o(1).
\]

\paragraph{Resources (every execution).} For one pattern: computing its rational parameters costs $\poly(n)$;
drawing and storing the block tables costs $2^{c_R\sqrt n}\poly(n)$ by~\eqref{eq:sumR}; each of the at most
$2^{\tau n+3}$ digits per side costs $2^{c_R\sqrt n}\poly(n)$ before its entries are emitted; at most $C_n$ entries
are emitted per side, each at polynomial cost; sorting, prefix sums and each random index cost polynomially per
entry or sample, the integers involved ($L_\pi\le C_n^2$ and its prefix sums) having $O(n)$ bits; each of the at
most $\max(N,C_n)$ sampled or enumerated matches costs $2^{c_R\sqrt n}\poly(n)$ for its multiplicity, $\poly(n)$ for
its pattern, and $\poly(n)$ per target for $J_\rho(k)$. Multiplying by $Q_n\le2^{c_{\cP}\sqrt n}$ patterns gives
part~(2) with $c_B=c_{\cP}+c_R+1$. The number of sampler calls is at most the number of elementary operations,
which is below $2^{2n}$ for large $n$. \qed

\begin{remark}[what is used about the filter]\label{rem:onefilter}
The proof uses the filter only through Lemma~\ref{lem:order} (order and bounds), Lemma~\ref{lem:eval} (evaluation)
and the two finite moments~\eqref{eq:finmom}. In the language of~\cite{Base}, whose Theorem~B uses a family of $43$
filters and ten regions of $W_\pi$ with multiplier vectors $z$, the present case is the family with one filter and
the vectors $z=(p_1,0)$ for $W_\pi\le\Gamma$ and $z=(p_2,0,\dots)$ for $W_\pi>\Gamma$; then all ten certificates
reduce to \textup{(M1)} and \textup{(M2)}. We have written the argument directly for one filter, so the regions and
the comparison rows of~\cite{Base} do not appear, and the general ``$6\eps$ per row'' accounting, which
needs $a_\Gamma\le\tfrac34$ (otherwise a row combining a $3\eps$ and a $4\eps$ error costs $7\eps$), is replaced by
the exact constants $c_1,c_2$ of Lemma~\ref{lem:B2}. With $7\max(p_1,p_2)$ in place of $6\max(p_1,p_2)$ the
premise \textup{(N13)} would still hold at our point (slack $7.2\cdot10^{-7}$).
\end{remark}
\section{The final algorithm and its correctness}\label{sec:final}

\subsection{The algorithm}
Let $n_\ast$ be a fixed integer, specified in Section~\ref{sec:nstar}.
\begin{procedure}[main algorithm]\label{proc:main}
\begin{enumerate}[label=\arabic*.]
\item If $n<n_\ast$, or if the guard $G_0\ceil{\log(b+1)}\le n$ fails, run exact meet in the middle on the original
instance (Section~\ref{sec:fallback}) and stop. (Here $b$ is the bit length of the \emph{original} instance.)
\item Draw the tie weights and form $c_i$ and $t_s$ by~\eqref{eq:transform}.
\item Run the preliminary search (Procedure~\ref{proc:prelim}). If it returns \textsc{yes} or \textsc{no}, return
that answer.
\item Draw a prime $P$ of $[2^{\floor{\eta n}},2^{\floor{\eta n}+1})$ by Lemma~\ref{lem:samplers}(c), and independent
uniform phases $\varphi_1,\dots,\varphi_n\in\Z/P\Z$ by Lemma~\ref{lem:samplers}(a); store them. Put $\bc_i=c_i\bmod P$.
If a sampler fails, return \textsc{no}.
\item Run the capped alias procedure (Procedure~\ref{proc:Acap}) on $(c,t,P,\varphi)$. If it fails, return \textsc{no}.
\item Run the checksum procedure of Theorem~B (with $d=3$) on $(P,\bc,\varphi)$ and the target list
$\cT=\{t_s\bmod P:0\le s\le nN_0\}$. If it fails, return \textsc{no}.
\item For each $s$ put $\hat B_s=\hat V_{t_s\bmod P}-\hat A_s$. Return \textsc{yes} if $\abs{\hat B_s}^2\ge\tfrac14$ for
some $s$, and \textsc{no} otherwise. (The real and imaginary parts are dyadic rationals with $n^2+O(n)$ fractional
bits, so the comparison is exact.)
\end{enumerate}
\end{procedure}
The procedures of Steps~3, 5 and~6 read only their declared inputs and their own fresh random words; in
particular Step~5 does not modify the stored phases, and Steps~4--6 do not read the random words used in Step~3.

\subsection{Exact cancellation}
\begin{lemma}\label{lem:cancel}
For every prime $P$, all phases and every $s$,
\begin{equation}\label{eq:cancel}
 V_{t_s\bmod P}-\sum_{I\in\cA_s(P)}\chi(I)=B_s:=\sum_{I\subseteq[n]:\ c(I)=t_s}\chi(I).
\end{equation}
If no $I$ has $c(I)=t_s$ then $B_s=0$; if exactly one has, then $\abs{B_s}=1$.
\end{lemma}
\begin{proof}
Since $\bc(I)\equiv c(I)\pmod P$, the sum $V_{t_s\bmod P}$ of~\eqref{eq:V} runs over the $I$ with
$c(I)\equiv t_s\pmod P$; these are the aliases of $t_s$ together with the $I$ with $c(I)=t_s$. Each $\chi(I)$ has
modulus $1$.
\end{proof}

\subsection{Correctness}
For fixed tie weights $u$ the statement \textup{(LD)} is either true or false; the algorithm cannot test it.
Write $\mathrm{ISO}(u)$ for the event of Lemma~\ref{lem:iso}(b) (some $t_s$ has exactly one representing subset).
Put
\[
 \epsilon_{\mathrm{pre}}(n)=\frac8{\zeta^2n}+n2^{-n^3},\qquad
 \epsilon_{\mathrm{main}}(n)=\epsilon_A(n)+\epsilon_B(n)+e^{-c_\ast n^2}+2n^32^{-n^3},
\]
with $\epsilon_A,\epsilon_B$ from Theorems~A and~B; all depend on $n$ only and tend to $0$.

\begin{proposition}[correctness for large guarded inputs]\label{prop:correct}
Let $n\ge\max(n_A,n_B,1/\zeta)$ and let the input pass the guard and Step~1.
\begin{enumerate}[label=\textup{(\alph*)}]
\item On a \textsc{no} instance the algorithm returns \textsc{yes} with probability at most
$\max(\epsilon_{\mathrm{pre}}(n),\epsilon_{\mathrm{main}}(n))$.
\item On a \textsc{yes} instance it returns \textsc{no} with probability at most
$\tfrac1{12}+\max(\epsilon_{\mathrm{pre}}(n),\epsilon_{\mathrm{main}}(n))$.
\end{enumerate}
\end{proposition}
\begin{proof}
\emph{An extended experiment.} The random words read in Step~3 and those read in Steps~4--6 are disjoint and
therefore independent. Consider the experiment in which, after Step~3, Steps~4--7 are carried out \emph{in every
case} (also when Step~3 has returned an answer), reading their own words. The output of the algorithm is a function
of this experiment: it is the answer of Step~3 if that is \textsc{yes} or \textsc{no}, and otherwise the answer of
Step~7 (or \textsc{no} if Steps~4--6 fail). In the extended experiment the joint law of everything in Steps~4--7 is
the law of Steps~4--7 run on $(c,t)$ with fresh randomness; it is \emph{not} conditioned on the outcome of Step~3.
No statement below is made conditionally on \textsc{continue}.

Fix the tie weights $u$. In the extended experiment let $\mathcal S$ be the event that: no sampler of Step~4 fails;
the alias procedure does not fail and $\abs{\hat A_s-\sum_{I\in\cA_s(P)}\chi(I)}\le\tfrac1{16}$ for all $s$; and the
checksum procedure does not fail and $\abs{\hat V_\rho-V_\rho}\le\tfrac18$ for all $\rho\in\cT$. On $\mathcal S$,
by~\eqref{eq:cancel},
\begin{equation}\label{eq:resid}
 \abs{\hat B_s-B_s}\le\tfrac3{16}<\tfrac14\qquad\text{for every }s .
\end{equation}

\emph{Claim: if \textup{(LD)} holds for $u$, then $\Pr[\mathcal S^c\mid u]\le\epsilon_{\mathrm{main}}(n)$.} Couple
Step~4 with ideal draws (Section~\ref{sec:samplers}): there are a prime $P^\circ$ uniform among the primes of the
interval and, given $P^\circ$, independent uniform phases $\varphi^\circ$, such that $(P,\varphi)=(P^\circ,\varphi^\circ)$
unless a sampler of Step~4 fails, which has probability at most $e^{-c_\ast n^2}+2n^32^{-n^3}$. We bound the
probabilities of the remaining parts of $\mathcal S^c$ for the run on $(P^\circ,\varphi^\circ)$; this involves no
conditioning on the samplers' success. \emph{Alias part.} The random words of the alias procedure are independent
of $(P^\circ,\varphi^\circ)$, and $P^\circ$ is uniform among the primes of its interval. Let $\mathcal H$ be the
event of Theorem~A(1) for the run with the prime $P^\circ$; it is determined by $P^\circ$ and the alias
procedure's words, and $\Pr[\mathcal H^c]\le\epsilon_A(n)$. On $\mathcal H$ the alias part of $\mathcal S$ holds
for every phase vector, in particular for $\varphi^\circ$; no conditioning on $\varphi^\circ$ is needed, and it
does not matter that $\varphi^\circ$ depends on $P^\circ$. So the alias part of $\mathcal S$ fails
with probability at most $\epsilon_A(n)$. \emph{Checksum part.} Condition on $P^\circ$. The phases are independent
and uniform on $\Z/P^\circ\Z$ and independent of the checksum procedure's own words; $\abs{\log P^\circ-\eta n}\le1$;
the target list has at most $24n^2+1\le n^3$ residues and is a function of $(c,t,P^\circ)$. By
Theorem~B(1) the checksum part of $\mathcal S$ fails with probability at most $\epsilon_B(n)$. That the
alias procedure reads the same phases is irrelevant here: the checksum event is a function of
$(P^\circ,\bc,\varphi^\circ)$ and the checksum's own words only. The union bound proves the claim.

\emph{Case 1: \textup{(LD)} holds for $u$.} If Step~3 returns an answer, it is correct
(Proposition~\ref{prop:P}(a)). Otherwise the output is that of Steps~4--7. On $\mathcal S$: for a \textsc{no}
instance every $B_s=0$ (Lemma~\ref{lem:iso}(a) and Lemma~\ref{lem:cancel}), so $\abs{\hat B_s}<\tfrac14$ for all $s$
by~\eqref{eq:resid} and the output is \textsc{no}; for a \textsc{yes} instance with $\mathrm{ISO}(u)$, some
$\abs{B_s}=1$, so $\abs{\hat B_s}>\tfrac34$ and the output is \textsc{yes}. Hence, given $u$, the algorithm errs with
probability at most $\epsilon_{\mathrm{main}}(n)$ on a \textsc{no} instance, and at most
$\epsilon_{\mathrm{main}}(n)+\mathbf 1_{\neg\mathrm{ISO}(u)}$ on a \textsc{yes} instance.

\emph{Case 2: \textup{(LD)} fails for $u$.} By Proposition~\ref{prop:P}(c), Step~3 returns a correct answer except
with probability at most $\epsilon_{\mathrm{pre}}(n)$. So, given $u$, the algorithm errs with probability at most
$\epsilon_{\mathrm{pre}}(n)$, on either kind of instance.

In both cases the conditional error probability given $u$ is at most
\[
 \max(\epsilon_{\mathrm{pre}}(n),\epsilon_{\mathrm{main}}(n))+\mathbf 1_{\{\textsc{yes}\text{ instance}\}}\,\mathbf 1_{\neg\mathrm{ISO}(u)} .
\]
This bound is pointwise in $u$, so the dependence between $\mathrm{ISO}(u)$ and the truth of \textup{(LD)} does not
matter. Average over $u$ and use $\Pr[\neg\mathrm{ISO}]\le\tfrac1{12}$.
\end{proof}

\begin{remark}[how a wrong answer can arise]\label{rem:errors}
A \textsc{no} instance is answered \textsc{yes} only if Step~7 finds $\abs{\hat B_s}\ge\tfrac12$ although $B_s=0$,
which requires an inaccurate alias sum or checksum; \textsc{fail} outcomes are converted to \textsc{no} and the
preliminary search never errs. A \textsc{yes} instance is answered \textsc{no} if isolation fails, if a main
procedure fails or is inaccurate, or if \textup{(LD)} is false and the preliminary search nevertheless returns
\textsc{continue} and the main part then errs. The error is two-sided. Failure probabilities are carried to the
final decision only through the union bounds above; no expected-time bound is used in the correctness proof except
inside Theorem~A(1), where it bounds the probability of reaching the cap.
\end{remark}

\subsection{The cutoff $n_\ast$ and part (a) of Theorem~\ref{thm:main}}\label{sec:nstar}
Let $n_\ast$ be an integer such that for all $n\ge n_\ast$: (i) $n\ge\max(n_A,n_B,1/\zeta)$ and the thresholds of
\textup{(S1)--(S6)}, Lemma~\ref{lem:cov} and Proposition~\ref{prop:time} hold, and $2^{\eta n-1}\ge\max(P_0,P_{\eps/8})$;
(ii) $\max(\epsilon_{\mathrm{pre}}(n),\epsilon_{\mathrm{main}}(n))\le\tfrac1{12}$. Such an integer exists and depends
only on the fixed constants, because every one of these quantities has been bounded by an explicit function of $n$
tending to its limit; Section~\ref{sec:limits} estimates its size. Then by Proposition~\ref{prop:correct} the
algorithm answers correctly with probability at least $1-\tfrac1{12}-\tfrac1{12}=\tfrac56\ge\tfrac23$ on every input
with $n\ge n_\ast$ that passes the guard. On all other inputs it runs the exact fallback. Together with the halting
statement of Proposition~\ref{prop:time} this proves Theorem~\ref{thm:main}(a).

\subsection{The fallback}\label{sec:fallback}
Split $[n]$ into two halves, list all subset sums $a(I)$ of each half (including the empty set), sort one list and,
for each element $x$ of the other, search for $t-x$ by binary search. This decides the original instance exactly,
uses $O(n2^{n/2})$ operations on integers of at most $b+\ceil{\log(n+1)}\le w$ bits (one word each), and at most
$4\cdot2^{\ceil{n/2}}+O(1)$ words of storage with an in-place merge sort buffer of one list length.
\section{Running time and resource accounting}\label{sec:time}

\begin{proposition}[time, halting, address space]\label{prop:time}
There are constants $K_0,K$ such that the following hold.
\begin{enumerate}[label=\textup{(\alph*)}]
\item Every execution of Procedure~\ref{proc:main} halts.
\item Every execution on an input with $n\ge n_\ast$ that passes the guard uses at most
\begin{equation}\label{eq:alltime}
 K_0\,2^{(\tau_0+20\eps)n+2n^{0.7}}\,(n+b+2)^{K}
\end{equation}
word operations. Every execution on any other input uses at most $K_0\,2^{n/2}(n+b+2)^K$ word operations.
\item The number of memory words ever allocated is at most the number of operations, memory is allocated
sequentially, and for every input ($n\ge2$, $b\ge1$) every address used is below $2^{w}$.
\end{enumerate}
\end{proposition}
\begin{proof}
(b) We go through the steps for $n\ge n_\ast$ and a guarded input; every bound is for every execution.
\begin{itemize}
\item \emph{Step 1}: reading the input, computing $b$ by shifts, and the two comparisons cost $\poly(n+b)$. The
comparison of $n$ with the constant $n_\ast$ is a comparison with a fixed multiword literal.
\item \emph{Step 2}: $n\log N_0$ random bits and $O(n^2)$ arithmetic operations on integers of $O(b+\log n)$ bits.
\item \emph{Step 3}: at most $2^{\tau_0n}\poly(n+b+2)$ (Proposition~\ref{prop:P}(b)).
\item \emph{Step 4}: at most $n^3$ candidates, each tested by at most $2^{(\floor{\eta n}+1)/2}$ trial divisions of
$(\eta n+1)$-bit integers (Lemma~\ref{lem:samplers}(c)): at most $2^{\eta n/2}\poly(n)$, and $\eta/2<\tau_0$ by
\textup{(N9)}. The phases and the residues $\bc_i$ cost $\poly(n+b)$.
\item \emph{Step 5}: at most $2^{\tau_0n+2n^{0.7}}\poly(n+b+2)$ (Theorem~A(2)).
\item \emph{Step 6}: at most $2^{(\tau_0+20\eps)n+c_B\sqrt n}\poly(n)$ (Theorem~B(2), $\abs\cT\le24n^2+1$).
\item \emph{Step 7}: $O(n^2)$ subtractions and comparisons of dyadic numbers with $n^2+O(n)$ bits.
\end{itemize}
Since $c_B\sqrt n\le n^{0.7}$ for large $n$ (guaranteed by the choice of $n_\ast$), the sum is at
most~\eqref{eq:alltime}. For inputs with $n<n_\ast$ or failing the guard, Step~1 costs $\poly(n+b)$ and the fallback
costs $O(n2^{n/2})$ operations on one-word integers.

(a) Each step has just been bounded on every execution: all loops are bounded by the list caps, the step cap of
Procedure~\ref{proc:Acap}, the entry and match caps~\eqref{eq:Cn}, the fixed numbers of samples and trials, and the
bounded samplers. (An unbounded rejection sampler would give only almost-sure halting; this is why the samplers of
Lemma~\ref{lem:samplers} are capped.)

(c) Allocate memory by a single increasing pointer, a new word being allocated only when it is first written; then
every address is below $n+2$ plus the number of words written, which is at most the operation count. For
$n\ge n_\ast$ and a guarded input, by~\eqref{eq:alltime} the logarithm of the operation count is at most
$\log K_0+0.49n+2n^{0.7}+K\log(n+b+2)$. For $x\ge x_0(K)$ one has $K\log x\le x$; taking $n_\ast\ge x_0(K)$ and
$n_\ast$ so large that $\log K_0+2n^{0.7}+4\le n/2$, this is at most $n+(n+b+2)-4\le w-2$ for every $b\ge1$, since
$w\ge4n+4b$. (The term $4b$ of the word length pays for the factor $(n+b+2)^K$ whatever $b$ is; the guard is not
used here.) For the fallback the storage is $4\cdot2^{\ceil{n/2}}+O(1)$ words, and
$\log(4\cdot2^{\ceil{n/2}}+n+64)\le n/2+8\le4n+4\le w$ for all $n\ge2$, $b\ge1$.
\end{proof}

\paragraph{Proof of Theorem~\ref{thm:main}(b).} Fix $c\ge1$ and $T'>T=\tau_0+20\eps$. Let $N\ge n_\ast$ be such that
for all $n\ge N$: $G_0\ceil{\log(n^c+1)}\le n$ (so every input with $b\le n^c$ passes the guard) and
$2n^{0.7}+K\log(n+n^c+2)\le(T'-T)n$. Then by~\eqref{eq:alltime} every execution on an input with $n\ge N$ and
$b\le n^c$ uses at most $K_0\,2^{T'n}$ word operations. With $T'=0.4875755>0.4875754036$ this is the bound
$O(2^{0.4875755n})$. The program does not depend on $c$ or $T'$; only $N$ and the constant do. \qed

\begin{remark}[expected versus worst-case time]\label{rem:worst}
The bound~\eqref{eq:alltime} is a bound on every execution, including executions that return a wrong answer. It is
obtained from caps, not from expectations. The only expectation bound in the paper, Lemma~\ref{lem:A7}, is used to
show that the cap of the alias procedure is rarely reached; if it is reached the procedure fails and the algorithm
answers \textsc{no}, which costs correctness probability (accounted for in $\epsilon_A$), not time.
\end{remark}

\begin{remark}[polynomial factors]\label{rem:polyfactors}
On guarded inputs $b$ may be as large as $2^{n/G_0}$, so the factor $(n+b+2)^K$ in~\eqref{eq:alltime} can itself be
exponential in $n$. This is harmless for Theorem~\ref{thm:main}(b), which concerns $b\le n^c$, and it is why the
theorem is stated for polynomial bit length. Inside the proof of Theorem~A the bit length enters the
\emph{probability} of reaching the cap only to the first power, through the factor $b+2$ of Lemma~\ref{lem:A7},
and there the guard and \textup{(N1)} absorb it.
\end{remark}
\section{Numerical certification}\label{sec:cert}
This section says what is certified, by what means, and how each numerical premise of Section~\ref{sec:params}
corresponds to a line of the certificate. Certification of a numerical inequality and proof of a lemma are
different things; this section concerns only the former.

\subsection{Method}
A real number is represented by a ball (midpoint and rigorous radius) in the Arb library, through
\texttt{python-flint}, at 200 bits; every operation enlarges the radius so that the true value stays inside. An
inequality is certified when the whole ball lies on the correct side. The scalar premises \textup{(E0)},
\textup{(C1)--(C8)}, \textup{(N1)--(N14)} are evaluated directly from the exact decimals of Section~\ref{sec:params},
with $h(x)=-(x\ln x+(1-x)\ln(1-x))/\ln2$. The constant $\eps_0$ enters as the exact decimal $0.0000013156$; no
computed quantity is substituted for it, so every enclosure containing $\eps_0$ has the direction required by the
inequality it certifies.

For the moments \textup{(M1)}, \textup{(M2)} the certificate evaluates the formulas~\eqref{eq:uomega} on the
continuous circle. With $F=2^r$, $G=2^s$, $Z_f=\int F$, $Z_g=\int G$, $Z_{f,\kappa}=\int F^\kappa$,
$Z_{g,\kappa}=\int G^\kappa$, $\mathcal P(x)=\int F(x+w)G(w)\,dw$ and
$\mathcal Y(x)=\int F(x+w)G(w)\,(r(x+w)+s(w))\,dw$, it forms
\[
 \mathsf L=\log\mathcal P-\log(Z_fZ_g),\qquad\mathsf B=\mathcal Y/\mathcal P-\log(Z_fZ_g),\qquad
 \mathsf H=\log Z_{f,\kappa}+\log Z_{g,\kappa}-\kappa\log(Z_fZ_g),
\]
and $u=\mathsf L+(\kappa-1)\mathsf B-\mathsf H$, $\omega=(2-\kappa)\mathsf B-2\mathsf L+\mathsf H$. In the notation
of~\eqref{eq:means}, $\mathcal P=q$, $\mathcal Y=Y_r+Y_s$, $Z_{f,\kappa}=N_r$, $Z_{g,\kappa}=N_s$; the terms
$\log(Z_fZ_g)$ cancel, and these expressions are identical to~\eqref{eq:uomega}.

\paragraph{Orientation.} The certificate uses the convention ``$F$ at $x+w$, $G$ at $w$'' of
Section~\ref{sec:filter}: $\mathcal P$ is a \emph{correlation} of the two densities. Two other conventions must be
distinguished from it and from each other. \emph{Exchanging the densities}, that is, using
$\int G(x+w)F(w)\,dw=\mathcal P(-x)$, only reflects $x$: it replaces $u(x)$, $\omega(x)$ by $u(-x)$, $\omega(-x)$,
and because $\Lambda$ is even the two moments are unchanged. \emph{Replacing the correlation by the convolution}
$\int F(x-w)G(w)\,dw$, which amounts to reflecting one of the two densities, is a different filter. With the
asymmetric densities of Appendix~\ref{app:filter} it gives $\log M_1=3.27728\ldots$ and $\log M_2=3.27609\ldots$
(certified enclosures from the same code with the sine coefficients of $s$ negated), far above $\rho_c$; so
\textup{(M1)} and \textup{(M2)} hold for the stated convention only. Reflecting both densities again changes
nothing.

\paragraph{Quadrature.} Each integral is a trapezoid sum with a rigorous error term
from the bound of Trefethen and Weideman~\cite[Theorem~3.2]{TW}: if $g$ is $2\pi$-periodic, analytic and
$\abs g\le M$ in the strip
$\abs{\operatorname{Im}z}<\mathsf a$, the $N$-point trapezoid mean differs from the mean of $g$ by at most
$2M/(e^{\mathsf aN}-1)$. (The theorem is stated in~\cite{TW} for the integral over a period, with the constant
$4\pi M$; dividing by $2\pi$ gives the bound for means.) What has to be supplied is, for every integrand, a strip
on which it is analytic and a bound for it there. For a real trigonometric polynomial $r$ as in
Section~\ref{sec:filter} and $y\ge0$ write
\[
 S_r(y)=\sum_{j=1}^{J}\bigl(\abs{r_j}+\abs{r'_j}\bigr)\cosh(jy),
\]
so that $\abs{r(z)}\le S_r(\abs{\operatorname{Im}z})$ for complex $z$, because
$\abs{\cos jz},\abs{\sin jz}\le\cosh(j\operatorname{Im}z)$. Consequently $\abs{2^{r(z)}}=2^{\operatorname{Re}r(z)}\le2^{S_r(\abs{\operatorname{Im}z})}$.

\emph{Inner integrals} (strip half-width $\sigma$, variable $w$, real $x$). The integrands of $Z_f$, $Z_{f,\kappa}$,
$\mathcal P(x)$ and $\mathcal Y(x)$ are entire and $2\pi$-periodic in $w$, and for $\abs{\operatorname{Im}w}<\sigma$
they are bounded by $2^{S_r(\sigma)}$, $2^{\kappa S_r(\sigma)}$, $2^{S_r(\sigma)+S_s(\sigma)}$ and
$(S_r(\sigma)+S_s(\sigma))\,2^{S_r(\sigma)+S_s(\sigma)}$ respectively (likewise for $g$). No further argument is
needed for them.

\emph{Outer integrals} (strip half-width $2\beta$, variable $x$). Here the integrands
$\Lambda\,2^{p_1(\Delta-u)}$ and $\Lambda\,2^{p_2(a_\Gamma\omega-u)}$ contain $\log\mathcal P$ and
$\mathcal Y/\mathcal P$, so one must show that $\mathcal P$ has no zero in the strip, fix the branch of the
logarithm, and bound $\operatorname{Re}u$ and $\operatorname{Re}\omega$ there. This is the content of the next
lemma. Write $f=F/Z_f$ and $g=G/Z_g$; these are entire, $2\pi$-periodic, positive on the real axis and of mean
$1$.

\begin{lemma}[strip bounds for the moment integrands]\label{lem:strip}
Let $\beta>0$ and put
\[
 n_f=\int_{\Tc}\abs{f(w+\mathrm i\beta)}^2dw,\qquad n_g=\int_{\Tc}\abs{g(w+\mathrm i\beta)}^2dw,\qquad
 q_0=\sqrt{(n_f-1)(n_g-1)} .
\]
Assume $q_0<1$. Then for every complex $z$ with $\abs{\operatorname{Im}z}<2\beta$:
\begin{enumerate}[label=\textup{(\alph*)}]
\item $\bigl\lvert\mathcal P(z)/(Z_fZ_g)-1\bigr\rvert\le q_0$. Hence $\mathcal P$ has no zero in the strip;
$\mathsf L(z)=\operatorname{Log}\bigl(\mathcal P(z)/(Z_fZ_g)\bigr)$, with the principal branch (base $2$), is
analytic there and coincides on the real axis with the real function $\mathsf L$ of the certificate; and
$\abs{\operatorname{Re}\mathsf L(z)}\le\lambda_0:=-\log(1-q_0)$.
\item $\bigl\lvert\mathcal Y(z)/\mathcal P(z)\bigr\rvert\le\sqrt{n_fn_g}\,\bigl(S_r(\beta)+S_s(\beta)\bigr)/(1-q_0)$, so that
$\abs{\mathsf B(z)}\le B^\ast$ with
\[
 B^\ast=\frac{\sqrt{n_fn_g}\,\bigl(S_r(\beta)+S_s(\beta)\bigr)+(1+q_0)\abs{\log(Z_fZ_g)}}{1-q_0}.
\]
\item $u$ and $\omega$ extend to analytic functions on the strip with
\[
 \abs{\operatorname{Re}u}\le U^\ast:=\lambda_0+\abs{\kappa-1}B^\ast+\abs{\mathsf H},\qquad
 \abs{\operatorname{Re}\omega}\le W^\ast:=2\lambda_0+\abs{2-\kappa}B^\ast+\abs{\mathsf H},
\]
and the two moment integrands are analytic on the strip and bounded there by
$(2+2\cosh2\beta)\,2^{p_1(\Delta+U^\ast)}$ and $(2+2\cosh2\beta)\,2^{p_2(a_\Gamma W^\ast+U^\ast)}$ respectively.
\end{enumerate}
\end{lemma}
\begin{proof}
\emph{Norms on shifted lines.} Since $f$ is entire and $2\pi$-periodic it has an absolutely convergent Fourier
series $f(w)=\sum_{j\in\Z}\hat f_je^{\mathrm ijw}$, valid for complex $w$, with $\hat f_0=\int f=1$. For real $t$,
$f(w+\mathrm it)=\sum_j\hat f_je^{-jt}e^{\mathrm ijw}$; so the mean of $w\mapsto f(w+\mathrm it)$ over $\Tc$ is $1$
for every $t$, and by Parseval
\[
 \int_{\Tc}\abs{f(w+\mathrm it)}^2dw=\sum_j\abs{\hat f_j}^2e^{-2jt}=1+\sum_{j\ge1}\abs{\hat f_j}^2\,2\cosh(2jt),
\]
where the last step uses $\hat f_{-j}=\overline{\hat f_j}$, which holds because $f$ is real on the real axis. This
is an even function of $t$, nondecreasing in $\abs t$. Hence for $\abs t\le\beta$,
\begin{equation}\label{eq:shiftnorm}
 \int_{\Tc}\abs{f(w+\mathrm it)}^2dw\le n_f,\qquad\int_{\Tc}\abs{f(w+\mathrm it)-1}^2dw=\int_{\Tc}\abs{f(w+\mathrm it)}^2dw-1\le n_f-1,
\end{equation}
and the same for $g$ with $n_g$.

\emph{Contour shift.} Let $z=x+\mathrm iy$ with $\abs y<2\beta$. The function
$w\mapsto f(z+w)g(w)$ is entire and $2\pi$-periodic, so its integral over a period is the same along every
horizontal line (Cauchy's theorem on a rectangle; the vertical sides cancel by periodicity). Taking the line
$\operatorname{Im}w=-y/2$,
\[
 \frac{\mathcal P(z)}{Z_fZ_g}=\int_{\Tc}f(z+w)g(w)\,dw=\int_{\Tc}f\bigl(x+w+\tfrac{\mathrm iy}2\bigr)\,g\bigl(w-\tfrac{\mathrm iy}2\bigr)\,dw .
\]
Both factors now have imaginary displacement $\abs y/2<\beta$, and both have mean $1$. Therefore
\[
 \frac{\mathcal P(z)}{Z_fZ_g}-1=\int_{\Tc}\Bigl(f\bigl(x+w+\tfrac{\mathrm iy}2\bigr)-1\Bigr)\Bigl(g\bigl(w-\tfrac{\mathrm iy}2\bigr)-1\Bigr)dw,
\]
and the Cauchy--Schwarz inequality with~\eqref{eq:shiftnorm} bounds the modulus by
$\sqrt{(n_f-1)(n_g-1)}=q_0$, which is the inequality of (a). So $\mathcal P(z)/(Z_fZ_g)$ lies in the closed disc of radius $q_0<1$ about $1$, which
does not contain $0$ and on which the principal logarithm is analytic; on the real axis the quotient is a positive
real number and the principal logarithm is the real one. As the modulus of the quotient lies in $[1-q_0,1+q_0]$,
$\abs{\operatorname{Re}\mathsf L}\le\max\bigl(-\log(1-q_0),\log(1+q_0)\bigr)=-\log(1-q_0)$.

(b) The same shift applied to $\mathcal Y$ gives
\[
 \frac{\mathcal Y(z)}{Z_fZ_g}=\int_{\Tc}f\bigl(x+w+\tfrac{\mathrm iy}2\bigr)g\bigl(w-\tfrac{\mathrm iy}2\bigr)
 \Bigl(r\bigl(x+w+\tfrac{\mathrm iy}2\bigr)+s\bigl(w-\tfrac{\mathrm iy}2\bigr)\Bigr)dw .
\]
The last factor is at most $S_r(\beta)+S_s(\beta)$ in modulus, and by the Cauchy--Schwarz inequality
and~\eqref{eq:shiftnorm} the integral of $\abs{f(\cdot)g(\cdot)}$ is at most $\sqrt{n_fn_g}$. Divide by
$\abs{\mathcal P(z)/(Z_fZ_g)}\ge1-q_0$. Finally $\mathsf B=\mathcal Y/\mathcal P-\log(Z_fZ_g)$ and
$(1+q_0)/(1-q_0)\ge1$, which gives the (slightly generous) bound $B^\ast$ used by the certificate.

(c) $u=\mathsf L+(\kappa-1)\mathsf B-\mathsf H$ and $\omega=(2-\kappa)\mathsf B-2\mathsf L+\mathsf H$, where
$\mathsf H$ is a real constant; by (a) and (b) both are analytic on the strip, with the stated bounds for their
real parts. For real $\theta$ and complex $\xi$, $\abs{2^{\theta\xi}}=2^{\theta\operatorname{Re}\xi}$; so
$\abs{2^{p_1(\Delta-u)}}\le2^{p_1(\Delta+U^\ast)}$ and
$\abs{2^{p_2(a_\Gamma\omega-u)}}\le2^{p_2(a_\Gamma W^\ast+U^\ast)}$, using $p_1,p_2,a_\Gamma,\Delta>0$. Finally
$\abs{2+2\cos z}\le2+2\cosh(\operatorname{Im}z)\le2+2\cosh2\beta$.
\end{proof}

\emph{How the certificate uses the lemma.} The number $n_f$ is itself an integral of an entire periodic function:
$\abs{F(w+\mathrm i\beta)}^2=2^{2r_\beta(w)}$ with the real trigonometric polynomial
$r_\beta(w)=\operatorname{Re}r(w+\mathrm i\beta)=\sum_j\cosh(j\beta)\,(r_j\cos jw+r'_j\sin jw)$, so
$n_f=Z_f^{-2}\int2^{2r_\beta}$ is enclosed by a trapezoid sum with the error term for the bound
$2^{2S_{r_\beta}(\sigma)}$; likewise $n_g$. The script then forms an upper bound for $q_0$, checks $q_0<1$, and
forms upper bounds for $B^\ast$, $U^\ast$, $W^\ast$ and for the two outer error terms
$2M/(e^{2\beta N}-1)$ with $M$ from~(c). The values of $u$ and $\omega$ at the outer nodes are balls that already
contain the inner quadrature errors. At the stored parameters ($\beta=0.1$, $\sigma=0.25$, $512$ inner nodes for
$\mathcal P$ and $\mathcal Y$, and $N=16384$ nodes for the outer integrals, the normalizers and $n_f,n_g$) the run
gives $q_0\le0.37084$, $U^\ast\le1.3258$, $W^\ast\le19.178$, inner error terms
below $5\cdot10^{-51}$ and outer error terms below $2^{-4700}$.

\subsection{Results}
The certified lower bounds of the slacks (right side minus left side, rounded down) are as follows. For
\textup{(M1)} and \textup{(M2)}, marked $\ast$, the number is the \emph{logarithmic} margin $\rho_c-\log M_i$, where
$M_i$ is the integral on the left of the premise; $M_i<2^{\rho_c}$ is equivalent to $\log M_i<\rho_c$.
\begin{center}
\begin{tabular}{llll}
\toprule
Premise & Slack $\ge$ & Premise & Slack $\ge$\\
\midrule
(C1) & $1.3160\cdot10^{-6}$ & (N3) & $0.0574$\\
(C2) & $1.3157\cdot10^{-6}$ & (N4) & $0.967$\\
(C3) & $1.3157\cdot10^{-6}$ & (N5) & $0.854$\\
(C4) & $1.3156\cdot10^{-6}$ & (N6) & $0.0325$\\
(C5) & $1.3157\cdot10^{-6}$ & (N7) & $0.930$\\
(C6) & $0.0411,\ 0.446,\ 0.00394$ & (N8) & $0.460,\ 0.434$\\
(C7) & $0.0162$ & (N9) & $0.229$\\
(C8) & $0.139$ & (N10) & $0.0286,\ 0.458$\\
(E0) & $7.02\cdot10^{-11}$ & (N11) & $0.687$\\
(N1) & $1.2956\cdot10^{-6}$ (exact) & (N12) & $7.45,\ 7.27$\\
(N2) & $0.0123$ & (N13) & $8.124\cdot10^{-7}$\\
(M1)$^\ast$ & $1.3154\cdot10^{-6}$ & (N14) & $10^{-6}$ (exact)\\
(M2)$^\ast$ & $1.4940\cdot10^{-6}$ & &\\
\bottomrule
\end{tabular}
\end{center}
The line \textup{(E0)} is the least, over the ten lines of \textup{(C1)--(C8)}, of the slack minus $\eps_0$; it is
attained by \textup{(C4)}, whose slack is $1.31567023384\ldots\cdot10^{-6}$.
The moments are $\log M_1=0.51626832585128515\ldots$ and $\log M_2=0.51626814726466968\ldots$ with ball radii below
$10^{-20}$, against $\rho_c=\eta-2.5\cdot10^{-6}=0.5162696413$.

\paragraph{Correspondence with the certificate file.} The script \texttt{certify\_selfcontained.py} supplied with
this manuscript prints its margins under the labels of Section~\ref{sec:params}: ten lines \textup{(E0)}, one for
each line of \textup{(C1)--(C8)} (\textup{(C6)} has three parts); the ten lines \textup{(C1)--(C8)} themselves;
seventeen lines \textup{(N1)--(N14)} (\textup{(N8)}, \textup{(N10)} and \textup{(N12)} have two parts each); and
the two moments. Eight further lines X1--X8 certify the auxiliary constants of the filter that are quoted in the
text: $\|r\|_1\le3.64$, $\|s\|_1\le3.60$, $R\le7.24$, $C_0\le15$ (Lemma~\ref{lem:order}); $C_1<2^{32}$,
$\beta_0<270$ and $\log P_0<339$ (Lemma~\ref{lem:B1} and the remark after it); and the variant of \textup{(N13)}
mentioned in Remark~\ref{rem:onefilter}, which is not a premise. The script computes
$\|r\|_1=3.63121121$, $\|s\|_1=3.59861571$, $R^1=24.49012608$, $C_0=14.9608\ldots$,
$C_1=3.62406\ldots\cdot10^{9}$ and $\beta_0=266.006\ldots$ For $P_0=\ceil A$, with $A$ the quantity under the
ceiling in Lemma~\ref{lem:B1}, the script certifies $\log(A+1)<339$, which suffices because $\ceil A\le A+1$; it
finds $\log A=326.727\ldots$ The moment computation is done
by \texttt{certify\_c5.py}, unchanged from~\cite{V}.

Conversely, every numerical inequality used in
Sections~\ref{sec:model}--\ref{sec:time} is one of \textup{(E0)}, \textup{(C1)--(C8)}, \textup{(N1)--(N14)},
\textup{(M1)}, \textup{(M2)}, one of the auxiliary bounds X1--X7, or one of the following facts of exact decimal
arithmetic, which the script checks with exact rationals: $\tau_0<\tfrac12$; $1/G_0=\zeta$;
$\tau_0+20\eps=0.4875754036$; $0.4875754036<0.4875755$; $\tau_0+20\eps<0.49$ (address space,
Proposition~\ref{prop:time}); $a\le\tfrac16$ (in \textup{(S5)}); $\eta<1$ (Lemma~\ref{lem:A5}, Step~4); and the
values of $r$, $\Delta$, $\Gamma$ and $\rho_c$ displayed in Section~\ref{sec:params}.

\paragraph{Relation to the certificate of~\cite{V}.} The certificate of~\cite{V} verifies the same inequalities at
the same parameter point, with the same moment code. It differs in the treatment of $\eps_0$: there $\eps_0$ was
computed as the least of the \textup{(C)}-slacks, whereas here $\eps_0$ is the fixed decimal of
Section~\ref{sec:params} and \textup{(E0)} is a premise, so that $\eps_0$ enters every other line as an exact number
and \textup{(N1)} is a fact of decimal arithmetic. The margins of~\cite{V} that do not involve $\eps_0$ are
reproduced to within one unit in the last place of a double (the present script rounds downward).

\paragraph{Reproduction.} From the directory \texttt{certificates/}, with Python~3.12 and \texttt{python-flint}
0.9.0, the stored run is
\begin{center}
\texttt{python certify\_selfcontained.py candidate\_base.json 512 32 0.1 0.25}
\end{center}
The arguments are the number of inner nodes ($512$), the ratio of outer to inner nodes ($32$, so $N=16384$),
$\beta$ (the outer strip has half-width $2\beta=0.2$) and the inner strip half-width $\sigma=0.25$; the working
precision is $200$ bits, and the run takes about ten seconds. Its output is \texttt{cert\_selfcontained\_512.json}.
The margins in that file are directed lower bounds (rounded toward $-\infty$). The script reads from
\texttt{candidate\_base.json} only the decimal strings of the parameters and of the coefficient arrays
\texttt{rc}, \texttt{sc}, whose order is that of Appendix~\ref{app:filter}; the fields \texttt{floating\_margins}
and \texttt{floating\_T} of that file are metadata of the floating-point optimization and are not read.

\paragraph{Independent computations.} According to~\cite{V}, the moments were computed by the script above and,
independently, by two separately written interval implementations (agreement to 20 digits; one with an analytic
quadrature bound, one with ball radius below $4\cdot10^{-55}$), and the scalar conditions were re-derived from the
sources by two reviewers. A further directed-interval computation by Astra (in its referee
report on~\cite{V}) gives logarithmic moment margins exceeding $1.3154487148\cdot10^{-6}$ and $1.4940353303\cdot10^{-6}$. For
this paper the script \texttt{certify\_selfcontained.py} was run as described above, and the reviewer re-ran it with
identical output. The reviewer also compared both logarithmic moments with a separately written floating-point
computation (an FFT correlation at $512$ to $4096$ points), which agreed to about fourteen decimal places; that
is a check of the formulas and of the orientation at floating-point accuracy, not a certificate.
All runs of the certificate share the moment code \texttt{certify\_c5.py}; they are same-code reproductions, and
the independence of the moment values rests on the separately written implementations reported in~\cite{V}. The
output file contains directed bounds, and the table above is rounded down.

\paragraph{Scope.} The certificate proves the inequalities listed above and nothing else. It does not check the
lemmas, the finite-circle transfer (Lemma~\ref{lem:B1}, proved in the text), or that the list of premises is
complete; completeness is the claim of the previous paragraph and is a statement about the written proof, to be
checked by the reader.
\section{Precisely scoped limits of the analysis}\label{sec:limits}
Each limitation below is stated together with the assumption that produces it and with what is known about
removing that assumption. None of them is a lower bound for Subset Sum.

\subsection{The result is asymptotic, and the cutoff is astronomically large}
Theorem~\ref{thm:main} is an asymptotic statement. The algorithm runs plain meet in the middle below $n_\ast$, and the
written proof needs, among others, the following thresholds (order-of-magnitude estimates, with
$C_0\approx14.96$, $\eps=\zeta=10^{-8}$, $c_{\cP}\approx201$):
\begin{center}
\begin{tabular}{p{0.52\linewidth}p{0.2\linewidth}l}
\toprule
Requirement & Where & $n\gtrsim$\\
\midrule
block concentration $n_\beta\ge6C_0^2/\eps^2$ with $n_\beta\approx\sqrt n$ & Lemma~\ref{lem:cov}(i) & $2\cdot10^{38}$\\
$3Q_n2^{-6\eps n+N_b+3}$ small, i.e.\ $6\eps n\gg c_{\cP}\sqrt n$ & Lemma~\ref{lem:cov}(iii) & $10^{19}$\\
$8/(\zeta^2n)\le\tfrac1{12}$ & Proposition~\ref{prop:P}(c), Section~\ref{sec:nstar} & $10^{18}$\\
$Q_n^22^{-\beta_3n}$ small with $\beta_3\approx8.1\cdot10^{-7}$ & Theorem~B & $3\cdot10^{17}$\\
$(\eps_0-\zeta-1/G_0)\,n\gg c_7n^{0.6}$ & Theorem~A & $10^{17}$\\
$(2n^2+3)e^{-n^{0.2}/2}$ small & Lemma~\ref{lem:A3} & $10^{10}$\\
the guard $G_0\ceil{\log(b+1)}\le n$ & Section~\ref{sec:params} & $10^{8}$\\
\bottomrule
\end{tabular}
\end{center}
Hence $n_\ast>10^{38}$ in this proof: the first row alone needs
$n\ge(6C_0^2/\eps^2)^2\approx1.8\cdot10^{38}$. Below $n_\ast$ the program runs meet in the middle, so on every input
that could be written down it takes the classical time $2^{n/2}\poly(n+b)$, and the proof gives no estimate
of a crossover point for inputs of feasible size. No statement about the size of the constant $C$ in
Theorem~\ref{thm:main}(b) follows from this: the theorem allows a threshold $N$, so the finite range of inputs
handled by the fallback need not be absorbed into $C$, and an upper bound on a running time does not yield a lower
bound on a constant. This proof supplies a bound of the form $C\,2^{T'n}$ only from $n\ge N\ge n_\ast$ on; a bound
valid from a stated small $n$ would have to be derived separately. \emph{Assumption responsible:} the small accuracy constant $\eps=10^{-8}$ (shared with~\cite{OAI}), used with
Chebyshev's inequality on blocks of size $\sqrt n$, and the thin numerical slacks ($10^{-6}$) at the frozen point.
\emph{Removable?} Partly: a point with larger slacks, an exponential tail bound in place of Chebyshev, and a larger
$\eps$ would lower $n_\ast$ at the price of a slightly larger exponent; we have not optimized this, and no choice
is known that makes the cutoff practical.

\subsection{Three closed-form bounds}
\begin{proposition}\label{prop:three}
Write $\eta=\tau_0+\Delta$ and $l+r=1-\mu$. The premises \textup{(C2)--(C5)} imply
\begin{enumerate}[label=\textup{(\roman*)}]
\item $\tau_0>\tfrac12-\mu\bigl(\gamma-\tfrac12h(a)\bigr)$ \quad(add \textup{(C2)} and \textup{(C3)});
\item $2\tau_0>1-\Delta+\mu\bigl[(h(a)-\gamma)+(h(\tfrac12-a)-1)\bigr]$ \quad(\textup{(C4)});
\item $\tau_0>\tfrac12\bigl(\tfrac12+\mu\,h(\tfrac12-a)\bigr)$ \quad(\textup{(C5)}).
\end{enumerate}
\end{proposition}
\begin{proof}
(i) $l+r+(h(a)-\gamma)\mu+(1-\gamma)\mu<2\tau_0$ and $l+r=1-\mu$ give $1+(h(a)-2\gamma)\mu<2\tau_0$. (ii) Substitute
$\eta=\tau_0+\Delta$ in \textup{(C4)}. (iii) is \textup{(C5)}.
\end{proof}
These are exact algebra on the premises of the present proof. They describe this proof's budget, not the problem:
at the frozen point (ii) and (iii) are the tight ones. \emph{Assumption responsible:} the specific accounting of
Lemma~\ref{lem:A7} (every trial lists all masks; every visited tuple is expanded). \emph{Removable?} Unknown;
\cite{V} reports numerical experiments on several variants, none of which is a theorem.

\subsection{A counting floor for alias enumeration}
\begin{proposition}[counting floor]\label{prop:floor}
Consider any procedure that, for a prime $P$ of about $\eta n$ bits, inserts a name for every alias class of a
target into a dictionary, spending at least one counted step per class, within a budget of $2^{(\tau_0+o(1))n}$
steps, and is combined with a checksum at the same modulus that requires $\eta-\tau_0\le\Delta^\ast$. If for some
input and target there are at least $2^{(1-\eta-o(1))n}$ alias classes, then
$\tau_0\ge1-\eta-o(1)\ge1-\tau_0-\Delta^\ast-o(1)$, that is $\tau_0\ge(1-\Delta^\ast)/2-o(1)$.
\end{proposition}
\begin{proof}
The number of steps is at least the number of classes inserted, so $\tau_0\ge1-\eta-o(1)$; and
$\eta\le\tau_0+\Delta^\ast$.
\end{proof}
At $\Delta^\ast=\Delta=0.0286969377$ the floor is $0.48565\ldots$. \emph{Assumptions responsible, and their status:}
(A1) every alias class is inserted individually --- removable in
principle (an exact signed aggregation over classes is known in a restricted model, but no algorithm that uses it
within the budget is known); (A2) the checksum uses the same modulus --- other combinations untested; (A3) the
checksum capacity is $\Delta^\ast$ --- see below; (A4) the input has $2^{(1-\eta-o(1))n}$ alias classes --- this is a
property of inputs whose subset sums are nearly equidistributed modulo $P$, not of all inputs (the weights
$P,2P,4P,\dots$ have distinct subset sums and no aliases of a target $\equiv1$). The proposition says nothing about
procedures that do not enumerate classes.

\subsection{A lower bound on the first moment}
\begin{lemma}[first-moment bound; Lemma~1 of~\cite{V}]\label{lem:first}
For all strictly positive continuous densities (in particular for every pair $r,s$ and every $\kappa>0$ in
Section~\ref{sec:filter}), every $\Delta\ge0$ and $p>0$,
\[
 \log\int_\Tc\Lambda(x)\,2^{p(\Delta-u(x))}\,dx\ \ge\ 1+p\Delta+(1+p)\log\int_\Tc(1+\cos x)^{1/(1+p)}\,dx .
\]
\end{lemma}
\begin{proof}
By Jensen's inequality for the probability density $K_x$,
$u(x)=\av{K_x\log\bigl(f_\ast(x+\cdot)g_\ast/K_x\bigr)}\le\log Q_\ast(x)$ with $Q_\ast(x)=\av{f_\ast(x+\cdot)g_\ast}$,
and $Q_\ast$ is a probability density on $\Tc$. Let $\varpi(x)=1+\cos x$, also a probability density. Then the
left integral equals $2^{1+p\Delta}\int\varpi\,2^{-pu}\ge2^{1+p\Delta}\int\varpi\,Q_\ast^{-p}$. By H\"older's inequality
with exponents $1+p$ and $(1+p)/p$,
$\int\varpi^{1/(1+p)}=\int(\varpi Q_\ast^{-p})^{1/(1+p)}Q_\ast^{p/(1+p)}\le\bigl(\int\varpi Q_\ast^{-p}\bigr)^{1/(1+p)}$.
\end{proof}
So \textup{(M1)} can hold only if $p_1\Delta<\rho_c-1-(1+p_1)\log\int(1+\cos x)^{1/(1+p_1)}dx$: the checksum capacity
$\Delta$ of any filter in this moment format is bounded. \emph{Assumption responsible:} the error is controlled
through a first-moment condition with one multiplier and a coordinatewise product filter. \emph{Removable?} In
principle, by an estimator outside the moment format; none is known. The consequences drawn from this lemma
in~\cite{V} for the original construction of~\cite{OAI} are not used or re-proved here.

\subsection{What is not claimed}
The paper does not claim: the conditional exponent $0.4873407528$ of~\cite{V} (which needs the open statement
\textup{[S]}); optimality of the parameter point; any lower bound; a derandomization; or a practical algorithm. The
correctness guarantee is two-sided error at most $1/3$; it can be amplified by repetition in the usual way.
\section{Verification status and supplementary files}\label{sec:verification}
The proof is divided into the statement groups S1--S29 of the table below. The codes mean:
\begin{description}[leftmargin=2.2em,style=nextline]
\item[W] A complete proof is written in this manuscript. It was produced by Claude, which had read the
corresponding passages of~\cite{OAI,Base}; it is a \emph{source-aware reconstruction}, not an independent proof.
\item[B] Two \emph{blind} derivations exist: two fresh instances of the same AI system were each given a
statement-only packet (no proof text, instructed not to open the sources) and wrote proofs from scratch. Their
conclusions agree with each other and with the proof written here; the comparison was made by the builder of this
manuscript, not by a referee. Blindness was by instruction and self-declaration, and both provers are the same
model as the builder: this is not cross-model independence. The packets and the derivations are working
documents and are not part of the supplement.
\item[R] Re-derived before this manuscript by two different AI systems (Claude and Astra),
each with partial exposure to the source proof.
\item[X] Re-derived for this manuscript by Astra from a statement-only packet, without
opening the manuscript during the derivation. That system had read an earlier draft, so the derivation is
independent in authorship but not blind. The comparison with the proof written here was made by the builder.
\item[C] Published classical theorem, cited.
\item[N] Numerical inequality, certified in ball arithmetic (Section~\ref{sec:cert}).
\end{description}

\paragraph{Cross-model review.} The manuscript was reviewed by Astra with the sources, the
earlier preprint~\cite{V} and the certificate files open; this is a \emph{source-aware} review, not a blind
re-derivation and not a formal verification. It traced the dependency chain of the main theorem and checked, among
others: the two cases of Corollary~\ref{cor:D} and its use in Lemma~\ref{lem:A5} under a signing fixed before the
half-block is drawn; the charging convention of Section~\ref{sec:procA} and item~7 of Lemma~\ref{lem:A7}; the event
$\mathcal H$ of Theorem~A; the change of measure and the joint match count of Lemma~\ref{lem:cov};
Lemmas~\ref{lem:B2} and~\ref{lem:B3}; the inverse-multiplicity factor in both branches of the estimator;
Proposition~\ref{prop:correct}; the proof of Lemma~\ref{lem:strip} against the moment code; the use of
\textup{(E0)}; and the final exponent. It compared the 64 coefficients of Appendix~\ref{app:filter} with the
parameter file and re-ran the certificate, reproducing every margin. The corrections it requested are
incorporated in this text. It did not write a second interval implementation of the moments.

\paragraph{Cross-model re-derivation and audit.} The same system then did three things, in separate sessions.
(i) From statement-only packets, and without opening the manuscript during the derivation, it proved
Lemmas~\ref{lem:cover}--\ref{lem:good}, Corollary~\ref{cor:D} and Lemma~\ref{lem:A5} (rows S8 and S13, code X),
using no hypothesis beyond those stated; for the collision step of Lemma~\ref{lem:A5} it used Jensen's inequality
where the proof here uses Markov's. It had read an earlier draft, so these derivations are independent in
authorship but not blind. (ii) With the manuscript open it audited Lemmas~\ref{lem:cov}--\ref{lem:B3} and
Theorem~B, looking for a missing hypothesis, a wrong constant or an unjustified independence, and judged each
sound. (iii) In the same way it audited the rest of the chain: Sections~\ref{sec:final} and~\ref{sec:time} and the
proof of Theorem~\ref{thm:main}; the rest of Section~\ref{sec:alias} with Lemma~\ref{lem:samplers}
and~\eqref{eq:couple}; Section~\ref{sec:filter} with Lemmas~\ref{lem:fourier}--\ref{lem:balance};
Lemmas~\ref{lem:iso}--\ref{lem:fmono}, \ref{lem:primes}, \ref{lem:phase}, \ref{lem:sums} and
Proposition~\ref{prop:P}; and the premises of Section~\ref{sec:params} against the proofs. It reported no gap in a
proof and no use of a statement of~\cite{OAI} or~\cite{Base} that is not proved here; the clarifications it
recommended are incorporated in this text. The comparison of its derivations with the proofs here was made by the
builder. No row has been assessed by a human referee.

\begin{center}
\small
\begin{longtable}{lp{0.27\linewidth}p{0.40\linewidth}p{0.18\linewidth}}
\toprule
ID & Where & Statement & Status\\
\midrule
\endhead
S1 & Lemma~\ref{lem:iso} & isolation & W, B\\
S2 & Lemma~\ref{lem:def} & deficiency; sign invariance & W, B\\
S3 & Lemmas~\ref{lem:azuma}--\ref{lem:fmono} & random subsets: variance, tail, monotone mean & W, B\\
S4 & Lemma~\ref{lem:primes} & random-prime collisions (prime count: C) & W, B, C\\
S5 & Proposition~\ref{prop:P} & capped preliminary search, \textup{(LD)} dichotomy & W, B\\
S6 & Lemma~\ref{lem:samplers}, \eqref{eq:couple} & bounded samplers, coupling & W\\
S7 & Lemmas~\ref{lem:phase},~\ref{lem:sums} & finite-precision sums & W, B (sums)\\
S8 & Lemmas~\ref{lem:cover}--\ref{lem:good}, Cor.~\ref{cor:D} & small layers, Corollary D & W, R, X\\
S9 & \eqref{eq:signcF}, \eqref{eq:classsum}, Lemma~\ref{lem:A1} & sign identity, classes, soundness & W, B\\
S10 & Lemma~\ref{lem:A2} & key and class counts & W, B\\
S11 & Lemma~\ref{lem:A3} & trial count & W, B\\
S12 & Lemma~\ref{lem:A4} & bad half-blocks & W, B\\
S13 & Lemma~\ref{lem:A5} & one trial & W, R, X\\
S14 & Corollary~\ref{cor:A6} & coverage & W, B\\
S15 & Lemma~\ref{lem:A7} & expected work & W, B\\
S16 & Theorem~A & Theorem A & W, B (given S13)\\
S17 & Lemma~\ref{lem:order} & filter functions, entropy order, bounds & W, B\\
S18 & Lemmas~\ref{lem:riemann}--\ref{lem:eval} & finite circle, with sine terms; evaluation & W, B\\
S19 & Lemma~\ref{lem:fourier} & Fourier form & W, B\\
S20 & Section~\ref{sec:tables}, Lemmas~\ref{lem:pattern},~\ref{lem:balance} & patterns, digits, tables & W, B\\
S21 & Lemma~\ref{lem:cov} & coverage and sizes; joint moment & W, B\\
S22 & Lemmas~\ref{lem:mult},~\ref{lem:cond} & multiplicity, direct sum, sampling error & W, B\\
S23 & \eqref{eq:paterr} & pattern errors & W, B\\
S24 & Lemmas~\ref{lem:B2},~\ref{lem:B3} & domination; one sampled pattern & W, B, R\\
S25 & Section~\ref{sec:checksum}, last proof & resources and precision & W, B\\
S26 & Theorem~B & Theorem B, one filter & W, B\\
S27 & Lemma~\ref{lem:cancel}, Prop.~\ref{prop:correct} & alias subtraction, decision rule & W, B\\
S28 & Section~\ref{sec:fallback}, Prop.~\ref{prop:time} & fallback, address space, time on every execution & W, B\\
S29 & Theorem~\ref{thm:main} & main theorem & W, B\\
N & Section~\ref{sec:cert} & premises (E0), (C1)--(C8), (N1)--(N14), (M1), (M2) & N\\
N$'$ & Lemma~\ref{lem:strip} & strip bounds behind the moment certificate & W\\
\bottomrule
\end{longtable}
\end{center}

\paragraph{Supplementary files.} The files are in the folder \texttt{selfcontained-2026} of the repository
\url{https://github.com/jv1000/subset-sum}; the supplement of~\cite{V} is the folder \texttt{exponent-2026} of the
same repository. Only the files under \texttt{certificates/} enter the proof.
(1) \texttt{certificates/candidate\_base.json}: the parameter point and filter
coefficients (also printed in Appendix~\ref{app:filter}).
(2) \texttt{certificates/certify\_selfcontained.py}, \texttt{certificates/certify\_c5.py} and
\texttt{certificates/cert\_selfcontained\_512.json}: the certificate of this paper and its output.

\paragraph{Use of AI tools.} The proofs in this manuscript were drafted by Claude under the author's
direction, from the sources and from the author's earlier preprint; the blind derivations were produced by further
instances of the same system; the review, the re-derivations and the audits described above were made by Astra, and the
corrections that answer them were written by Claude. The author is responsible for the
content. The role of AI systems is described here
in full so that a reader can weigh the verification claims accordingly.

\appendix
\section{The filter}\label{app:filter}
The log-density polynomials of Section~\ref{sec:filter} have degree $J=16$:
\[
 r(x)=\sum_{j=1}^{16}\bigl(r_j\cos jx+r'_j\sin jx\bigr),\qquad s(x)=\sum_{j=1}^{16}\bigl(s_j\cos jx+s'_j\sin jx\bigr),
\]
with the coefficients below, which are exact decimals. In \texttt{candidate\_base.json} the array \texttt{rc} holds
$r_1,\dots,r_{16},r'_1,\dots,r'_{16}$ in this order, and \texttt{sc} holds $s_1,\dots,s_{16},s'_1,\dots,s'_{16}$.
The filter is used with $\kappa=1.03465831$ and with the orientation of~\eqref{eq:fg}: $r$ is evaluated at $x+w$
and $s$ at $w$.
\begin{center}
\begin{tabular}{rrrrr}
\toprule
$j$ & $r_j$ (cos) & $r'_j$ (sin) & $s_j$ (cos) & $s'_j$ (sin)\\
\midrule
1 & $-$0.21919156 & 0.91131919 & $-$0.20782817 & 0.96462148\\
2 & $-$0.57160580 & 0.00973572 & 0.61194965 & 0.01416858\\
3 & 0.02024951 & 0.47988091 & 0.02603977 & 0.38660217\\
4 & 0.03552214 & $-$0.34819083 & 0.00846337 & 0.31951898\\
5 & $-$0.29988253 & $-$0.07175384 & $-$0.22018634 & $-$0.07671638\\
6 & $-$0.20387757 & 0.04897708 & 0.21025833 & $-$0.00950575\\
7 & 0.12066750 & 0.01245896 & 0.20455466 & $-$0.02875410\\
8 & $-$0.04068323 & $-$0.11710680 & 0.02047196 & 0.10160127\\
9 & $-$0.03536274 & $-$0.01505095 & $-$0.04648472 & $-$0.01822490\\
10 & $-$0.00662503 & $-$0.02221718 & $-$0.01884839 & $-$0.03524125\\
11 & $-$0.00401488 & 0.00282094 & 0.02926117 & $-$0.01087451\\
12 & 0.00734024 & 0.00166169 & 0.00689286 & 0.00271495\\
13 & 0.00751173 & $-$0.00656995 & $-$0.00353414 & 0.00602714\\
14 & $-$0.00392649 & $-$0.00332303 & $-$0.00297744 & $-$0.00342132\\
15 & 0.00011920 & 0.00231101 & $-$0.00005725 & $-$0.00147364\\
16 & 0.00117642 & $-$0.00007656 & 0.00128320 & $-$0.00005787\\
\bottomrule
\end{tabular}
\end{center}
The sums of the absolute values of the two columns of $r$ and of $s$ are $\|r\|_1=3.63121121$ and
$\|s\|_1=3.59861571$, and $R^1=\sum_jj(\abs{r_j}+\abs{r'_j}+\abs{s_j}+\abs{s'_j})=24.49012608$.

\begin{thebibliography}{99}
\raggedright
\makeatletter\g@addto@macro\UrlBreaks{\do\-}\makeatother
\bibitem{OAI} OpenAI, \emph{Subset Sum in Time $O(2^{0.49n})$}, manuscript with \LaTeX\ source, 4 October 2026.\newline
\url{https://github.com/openai/math/tree/main/preprints/Subset-Sum-in-Time-2-power-0-49n-October-4-2026}.
\bibitem{Base} A.~Kapila, \emph{Subset Sum in time $O(2^{0.4877n})$: a complete proof}, manuscript
(\texttt{subset-sum/PROOF\_0.4877.tex}), 14 October 2026.\newline \url{https://github.com/a-kapila/maths-proofs}, commit
\texttt{44c34961}.
\bibitem{V} J.~Voitle, \emph{A certified exponent of $0.48758$ for worst-case Subset Sum, and limits of the
alias--checksum analysis}, preprint, hexagon:2611.00178, November 2026.
\bibitem{RS} J.~B.~Rosser and L.~Schoenfeld, Approximate formulas for some functions of prime numbers,
\emph{Illinois J.\ Math.}\ 6 (1962), 64--94.
\bibitem{TW} L.~N.~Trefethen and J.~A.~C.~Weideman, The exponentially convergent trapezoidal rule,
\emph{SIAM Review} 56 (2014), 385--458.
\bibitem{MVV} K.~Mulmuley, U.~V.~Vazirani and V.~V.~Vazirani, Matching is as easy as matrix inversion,
\emph{Combinatorica} 7 (1987), 105--113.
\bibitem{HS} E.~Horowitz and S.~Sahni, Computing partitions with applications to the knapsack problem,
\emph{J.\ ACM} 21 (1974), 277--292.
\bibitem{KLM} R.~M.~Karp, M.~Luby and N.~Madras, Monte-Carlo approximation algorithms for enumeration problems,
\emph{J.\ Algorithms} 10 (1989), 429--448.
\bibitem{CJRS} X.~Chen, Y.~Jin, T.~Randolph and R.~A.~Servedio, Subset Sum in time $2^{n/2}/\poly(n)$, in
\emph{APPROX/RANDOM 2023}; arXiv:2301.07134.
\bibitem{LY} B.~Li and L.~Ye, \emph{Subset Sum via Partial Match}, preprint, arXiv:2610.10446v2, 7 October 2026,
revised 2 November 2026.
\bibitem{W} J.~Wharton, \emph{subset-sum-bound}, manuscript and repository, 20 October 2026.\newline
\url{https://github.com/JoshuaQFEX/subset-sum-bound}, commit \texttt{bed8e18}.
\bibitem{B} Y.~Bills, \emph{subset-sum-conditional-refinement}, manuscript and repository, v1.0.0, 21 October 2026,
and v1.1.0 (verification supplement), 22 October 2026.\newline
\url{https://github.com/yannouli/subset-sum-conditional-refinement}, commit \texttt{590bb68}.
\end{thebibliography}
\end{document}
